%% ============================================================================
%%  CAPITOLO 4 — Limiti di funzioni. Funzioni continue
%%  Fonte principale per argomenti e dimostrazioni: Marcellini–Sbordone, cap. 4
%%  (§29 premessa, §30 definizioni, §31 legame con le successioni, §32 esempi e
%%  proprietà, §33 funzioni continue, §34 discontinuità, §35 teoremi sulle
%%  funzioni continue; appendice §36 bisezione, §37 teorema di Weierstrass,
%%  §38 funzioni monotòne e inverse).
%% ============================================================================
\aperturacapitolo{Limiti di funzioni. Funzioni \texorpdfstring{\mbox{continue}}{continue}}{cap:funzioni-continue}

\epigraph{La funzione $f(x)$ resterà continua rispetto a $x$ tra i limiti dati se, tra
questi limiti, un incremento infinitamente piccolo della variabile produce sempre un
incremento infinitamente piccolo della funzione stessa.}{Augustin-Louis Cauchy,
\emph{Cours d'analyse} (1821)}

\lettrine[lines=3, findent=2pt, nindent=0pt]{N}{el capitolo precedente} abbiamo studiato
il comportamento di una successione quando l'indice $n$ cresce senza fine. Ora facciamo
lo stesso passo per le funzioni di una variabile reale: vogliamo capire come si comporta
$f(x)$ quando $x$ si avvicina a un punto $x_{0}$, eventualmente un punto in cui $f$ non è
nemmeno definita, oppure quando $x$ cresce o decresce illimitatamente. È la domanda che
si pone chiunque voglia sapere quale valore «dovrebbe» avere una grandezza fisica in una
condizione che non si può realizzare esattamente: la velocità di un corpo in un istante,
il coefficiente di portanza a incidenza nulla calcolato da misure a incidenze diverse da
zero, il valore di una formula nel punto in cui il denominatore si annulla.

La strada che percorreremo è quella del testo di riferimento \cite{marcellini-sbordone}:
\emph{definiremo il limite di una funzione per mezzo dei limiti di successioni}, in modo
da ereditare senza fatica tutti i risultati del \cref{cap:successioni}; poi
dimostreremo che questa definizione equivale a quella con $\epsilon$ e $\delta$, dovuta a
Cauchy e Weierstrass. Sui limiti costruiremo il concetto di \emph{funzione continua},
che formalizza l'idea intuitiva di un grafico «tracciato senza staccare la matita dal
foglio», e arriveremo ai grandi teoremi di esistenza dello schema della
\cref{fig:schema-teoremi}: il teorema degli zeri, quello dei valori intermedi e il
teorema di Weierstrass. Sono teoremi che garantiscono l'esistenza di qualcosa (una
radice, un valore, un massimo) senza dire dove si trovi: ma la dimostrazione del primo,
con il metodo di bisezione, è anche un algoritmo per calcolarla.

\medskip
\noindent\textbf{In questo capitolo.}
\begin{itemize}
  \item La funzione $\frac{\sen x}{x}$ e l'idea di limite di una funzione
  (\S\ref{sec:premessa-funzioni}).
  \item Intervalli, intorni, punti di accumulazione; la definizione di limite, finito e
  infinito, destro e sinistro (\S\ref{sec:def-limite-funzioni}).
  \item L'equivalenza tra la definizione «con le successioni» e quella con $\epsilon$ e
  $\delta$ (\S\ref{sec:legame-successioni}).
  \item Esempi notevoli; operazioni, confronti e limiti di funzioni composte
  (\S\ref{sec:proprieta-limiti-funzioni}).
  \item Funzioni continue e loro proprietà; le discontinuità e il prolungamento per
  continuità (\S\ref{sec:continue}, \S\ref{sec:discontinuita}).
  \item Permanenza del segno, teorema degli zeri, teoremi dei valori intermedi e di
  Weierstrass, criterio di invertibilità (\S\ref{sec:teoremi-continue}).
  \item Appendice: il metodo di bisezione, la dimostrazione del teorema di Weierstrass,
  le funzioni monotòne e le funzioni inverse; la continuità di esponenziali, logaritmi e
  potenze, e i limiti notevoli che ne seguono (\S\ref{sec:bisezione},
  \S\ref{sec:dim-weierstrass}, \S\ref{sec:monotone-inverse}).
  \item Approfondimento: su un intervallo le funzioni continue e iniettive sono
  strettamente monotòne (\S\ref{subsec:continue-iniettive}).
  \item Punti fondamentali del capitolo (\S\ref{sec:sintesi-cap04}).
\end{itemize}

%% ----------------------------------------------------------------------------
%%  §4.1  Premessa: la funzione sen x / x
%% ----------------------------------------------------------------------------
\section{Premessa: la funzione \texorpdfstring{$\frac{\sen x}{x}$}{sen x / x}}
\label{sec:premessa-funzioni}

Consideriamo la funzione
\begin{equation}\label{eq:senx-x}
  f(x)=\frac{\sen x}{x},
\end{equation}
che è definita per ogni $x\in\R\setminus\{0\}$. Allo scopo di disegnarne il grafico,
osserviamo preliminarmente che $-1\le\sen x\le1$ per ogni $x\in\R$
\eqref{eq:sen-cos-limitate}; dividendo tutti i membri per $x>0$ (che conserva il verso
delle disuguaglianze, \cref{prop:ordine-prodotto}) otteniamo
\begin{equation}\label{eq:senx-stima}
  -\frac1x\le f(x)\le\frac1x \qquad\perogni x>0 .
\end{equation}
Se $x<0$ si ottiene la situazione analoga $\frac1x\le f(x)\le-\frac1x$; comunque, essendo
$f$ una \emph{funzione pari}, cioè
\[
  f(-x)=\frac{\sen(-x)}{-x}=\frac{-\sen x}{-x}=f(x) \qquad\perogni x\in\R\setminus\{0\},
\]
è sufficiente studiarne le proprietà per $x>0$, riportando poi il disegno anche per
$x<0$ per simmetria rispetto all'asse $y$.

Le funzioni $y=-1/x$ e $y=1/x$, che appaiono nella stima \eqref{eq:senx-stima}, hanno per
grafico dei rami di iperbole. Tenendo conto del segno di $f(x)$, che per $x>0$ è lo
stesso segno di $\sen x$, e del fatto che $f$ si annulla esattamente nei punti $x=k\pi$
con $k\in\Z\setminus\{0\}$, si ottiene per $f$ un grafico come quello della
\cref{fig:senx-iperboli}: un'oscillazione che si smorza, costretta tra le due iperboli.
Il disegno è significativo per $x$ sufficientemente grande; è invece del tutto
indeterminato per $x$ «vicino» a zero, dove le due iperboli si allontanano e non
forniscono più alcuna informazione utile.

\begin{figure}[htb]
\centering
\begin{tikzpicture}
\begin{axis}[width=\textwidth, height=5.2cm, xmin=0, xmax=13.6, ymin=-1.15, ymax=1.6,
  axis lines=middle, xtick={3.14159,6.28318,9.42478,12.56637},
  xticklabels={$\pi$,$2\pi$,$3\pi$,$4\pi$}, ytick={-1,1},
  xlabel={$x$}, ylabel={$y$}, every axis x label/.style={at={(axis description cs:1,0.42)},anchor=west},
  tick label style={font=\scriptsize}]
  \addplot[tratteggio, domain=0.62:13.4, samples=120] {1/x};
  \addplot[tratteggio, domain=0.62:13.4, samples=120] {-1/x};
  \node[anchor=south west, font=\scriptsize] at (axis cs:1.75,0.6) {$y=\frac1x$};
  \node[anchor=north west, font=\scriptsize] at (axis cs:1.75,-0.6) {$y=-\frac1x$};
  \addplot[curva, domain=0.62:13.4, samples=240] {sin(deg(x))/x};
\end{axis}
\end{tikzpicture}
\caption{Il grafico di $f(x)=\frac{\sen x}{x}$ per $x>0$, compreso tra i rami di iperbole
$y=\pm\frac1x$ (tratteggiati). Il grafico tocca le iperboli dove $\sen x=\pm1$ e
attraversa l'asse $x$ nei punti $k\pi$. Vicino a $x=0$ le iperboli non dicono nulla.}
\label{fig:senx-iperboli}
\end{figure}

\subsection{Che cosa succede vicino a zero?}

Ricordiamo che la funzione \eqref{eq:senx-x} non è definita per $x_{0}=0$, cioè non è
calcolabile $f(0)$. È invece possibile calcolare valori di $f(x)$ per $x$ «vicino» a
zero; ad esempio, usando la \cref{tab:valori-notevoli}, si trova
\begin{gather*}
  f\Bigl(\frac\pi2\Bigr)=\frac2\pi\approx0{,}637,\qquad
  f\Bigl(\frac\pi3\Bigr)=\frac{3\sqrt3}{2\pi}\approx0{,}827,\\
  f\Bigl(\frac\pi4\Bigr)=\frac{2\sqrt2}{\pi}\approx0{,}900,\qquad
  f\Bigl(\frac\pi6\Bigr)=\frac3\pi\approx0{,}955,
\end{gather*}
e con un calcolatore si ottiene l'ulteriore tavola di valori prodotta dal programma
seguente.

\begin{python}[Valori di $\frac{\sen x}{x}$ vicino a zero]
import math

f = lambda x: math.sin(x) / x    # definita per x != 0, ed è pari

valori = [math.pi/2, math.pi/3, math.pi/4, math.pi/6,
          0.1, 0.01, 0.001, 1e-4, 1e-8]
for x in valori:
    print(f"x = {x:.8f}    f(x) = {f(x):.12f}")
print("f pari sui valori provati:",
      all(f(-x) == f(x) for x in valori))
\end{python}
\begin{risultato}
x = 1.57079633    f(x) = 0.636619772368
x = 1.04719755    f(x) = 0.826993343133
x = 0.78539816    f(x) = 0.900316316157
x = 0.52359878    f(x) = 0.954929658551
x = 0.10000000    f(x) = 0.998334166468
x = 0.01000000    f(x) = 0.999983333417
x = 0.00100000    f(x) = 0.999999833333
x = 0.00010000    f(x) = 0.999999998333
x = 0.00000001    f(x) = 1.000000000000
f pari sui valori provati: True
\end{risultato}

Nella tabella sono stati scelti valori di $x$ \emph{sufficientemente vicini} a $x_{0}=0$?
La domanda, così posta, non ha senso: uno, due, nove o mille valori di $x$ non sono, in
assoluto, né vicini né lontani da zero. Certamente, sulla base della tabella si può
\emph{intuire} che $f(x)$, per $x$ vicino a zero, assuma valori vicini al numero $1$, del
tipo $1-\epsilon$ con $\epsilon>0$ «piccolo». Non è però possibile escludere che, per
$x$ ancora più vicino a zero rispetto ai valori già considerati, $f(x)$ cambi
comportamento.

\needspace{8\baselineskip}
\begin{attenzione}[Una tabella non è una dimostrazione]
La funzione
\[
  g(x)=\frac{\sen x}{x}+\sen\frac\pi x \qquad(x\neq0)
\]
assume \emph{esattamente} gli stessi valori di $f$ in tutti i punti $x=10^{-k}$,
$k\in\N$, perché in quei punti $\sen\frac\pi x=\sen(10^{k}\pi)=0$. Una tabella costruita
con i valori $0{,}1$; $0{,}01$; $0{,}001$; $\dots$ non distinguerebbe $g$ da $f$. Eppure,
nei punti $x=\frac2{4m+1}$, che si avvicinano a zero al crescere di $m\in\N$, si ha
$\sen\frac\pi x=\sen\bigl(2m\pi+\frac\pi2\bigr)=1$, e quindi lì $g(x)=f(x)+1$ è vicino
a $2$, non a $1$. Come vedremo nel \cref{sec:legame-successioni}, $g$ \emph{non} ha
limite per $x\to0$. La morale è che il comportamento di una funzione vicino a un punto
non si accerta con un numero finito di prove, ma con un ragionamento valido per
\emph{tutti} i modi di avvicinarsi al punto.
\end{attenzione}

\subsection{La tabella ideale}

Una formulazione rigorosa del comportamento di $f(x)$ per $x$ «vicino» a $x_{0}$ si
ottiene nel modo seguente. Si considera una tabella ideale, illimitata a destra, del tipo
\begin{center}
\renewcommand{\arraystretch}{1.4}
\begin{tabular}{c|ccccc}
  $x$ & $x_{1}$ & $x_{2}$ & $x_{3}$ & $\cdots$ & $x_{n}\ \cdots$\\ \hline
  $f(x)$ & $y_{1}=f(x_{1})$ & $y_{2}=f(x_{2})$ & $y_{3}=f(x_{3})$ & $\cdots$ & $y_{n}=f(x_{n})\ \cdots$
\end{tabular}
\end{center}
cioè si considera una generica successione $x_{n}$ che converge a $x_{0}$ (il termine
$x_{n}$ è «vicino» a $x_{0}$ se $n$ è «grande»), con $x_{n}\neq x_{0}$, e la
corrispondente successione $y_{n}=f(x_{n})$ dei valori assunti dalla funzione. Se $y_{n}$
converge a un numero $\ell$, \emph{e se il numero $\ell$ non cambia qualunque sia la
successione $x_{n}$ che converge a $x_{0}$}, allora si dice che la funzione $f$
\emph{ammette limite}, uguale a $\ell$, per $x\to x_{0}$. La clausola in corsivo è
essenziale: è proprio quella che la funzione $g$ dell'esempio precedente non soddisfa.

Torniamo alla funzione \eqref{eq:senx-x}. Abbiamo già dimostrato nel
\cref{teo:sen-x-su-x} che
\begin{equation}\label{eq:senx-x-succ}
  \lim_{n\to+\infty}f(x_{n})=\lim_{n\to+\infty}\frac{\sen x_{n}}{x_{n}}=1
\end{equation}
\emph{qualunque sia} la successione $x_{n}$ che converge a $x_{0}=0$, con $x_{n}\neq0$ per
ogni $n\in\N$. In accordo con quanto detto sopra, in termini di limiti di funzioni la
\eqref{eq:senx-x-succ} equivale a dire che $f(x)$ ha limite $\ell=1$ per $x\to0$; in
simboli
\begin{equation}\label{eq:lim-senx-x}
  \lim_{x\to0}f(x)=\lim_{x\to0}\frac{\sen x}{x}=1 .
\end{equation}

Nella \cref{fig:senx-x} è riportato il grafico della funzione $\frac{\sen x}{x}$; la
funzione non è definita per $x_{0}=0$ (nel grafico ciò è evidenziato con un «tondino»
vuoto), però ha un comportamento regolare anche nelle vicinanze di $x_{0}=0$, e i valori
assunti da $f(x)$ sono vicini a $\ell=1$ quando $x$ è vicino a $x_{0}=0$.

\begin{figure}[htb]
\centering
\begin{tikzpicture}
\begin{axis}[width=\textwidth, height=5.2cm, xmin=-10.2, xmax=10.2, ymin=-0.45, ymax=1.3,
  axis lines=middle, xtick={-9.42478,-6.28318,-3.14159,3.14159,6.28318,9.42478},
  xticklabels={$-3\pi$,$-2\pi$,$-\pi$,$\pi$,$2\pi$,$3\pi$}, ytick={1},
  xlabel={$x$}, ylabel={$y$}, tick label style={font=\scriptsize},
  every axis x label/.style={at={(axis description cs:1,0.25)},anchor=west}]
  \addplot[curva, domain=-10:-0.01, samples=200] {sin(deg(x))/x};
  \addplot[curva, domain=0.01:10, samples=200] {sin(deg(x))/x};
  \addplot[tratteggio, domain=0:1] {1};
  \draw[blu, thick, fill=white] (axis cs:0,1) circle (2.2pt);
  \node[font=\scriptsize, anchor=south west] at (axis cs:1.5,0.75) {$f(x)=\dfrac{\sen x}{x}$};
\end{axis}
\end{tikzpicture}
\caption{Il grafico di $\frac{\sen x}{x}$. La funzione non è definita in $x=0$ (tondino
vuoto), ma i suoi valori si avvicinano a $1$ quando $x$ si avvicina a zero, da destra e
da sinistra: è la \eqref{eq:lim-senx-x}.}
\label{fig:senx-x}
\end{figure}

Viceversa, anche se la struttura analitica può sembrare a prima vista simile, il limite
\begin{equation}\label{eq:sen-1-x-anticipo}
  \lim_{x\to0}\sen\frac1x \qquad\text{non esiste}
\end{equation}
(lo dimostreremo nel \cref{sec:legame-successioni}). Nella \cref{fig:sen-1-x} è
riportato il grafico, eseguito al computer, della funzione $\sen\frac1x$: si noti in
particolare il comportamento caotico della funzione nelle vicinanze di $x_{0}=0$, dove
il grafico oscilla infinite volte tra $-1$ e $1$. Accanto è disegnato il grafico di
$x\sen\frac1x$, che oscilla anch'esso infinite volte, ma è «schiacciato» verso zero tra
i grafici di $y=\pm\abs x$: questa funzione, come vedremo, ha limite $0$.

\begin{figure}[htb]
\centering
\includegraphics[width=\textwidth]{sen_1_x}
\caption{A sinistra, il grafico di $\sen\frac1x$: vicino a zero la funzione assume
infinite volte tutti i valori tra $-1$ e $1$, e non ha limite. A destra, il grafico di
$x\sen\frac1x$, compreso tra i grafici di $y=\pm\abs x$: le oscillazioni sono smorzate e la
funzione tende a $0$ (\texttt{python/fig\_sen\_1\_x.py}).}
\label{fig:sen-1-x}
\end{figure}

\begin{idea}[Un valore «proibito»: la velocità istantanea]
La situazione di un limite in un punto in cui la funzione non è definita è tutt'altro che
artificiosa. Un carico sganciato da un aereo in volo orizzontale cade, trascurando la
resistenza dell'aria, percorrendo in verticale lo spazio $s(t)=\frac12gt^{2}$. La sua
velocità media verticale nell'intervallo di tempo tra $t_{0}$ e $t_{0}+h$ è
\[
  v_{m}(h)=\frac{s(t_{0}+h)-s(t_{0})}{h}=\frac{\frac12g(t_{0}^{2}+2t_{0}h+h^{2})-\frac12gt_{0}^{2}}{h}
  =gt_{0}+\frac g2\,h ,
\]
una formula che ha senso solo per $h\neq0$: per $h=0$ si otterrebbe $0/0$. La
\emph{velocità istantanea} all'istante $t_{0}$ è, per definizione, il limite di
$v_{m}(h)$ per $h\to0$, che qui vale $gt_{0}$: per ogni successione $h_{n}\to0$ con
$h_{n}\neq0$ si ha $v_{m}(h_{n})=gt_{0}+\frac g2h_{n}\to gt_{0}$. Questo è il germe del
concetto di \emph{derivata}, che studieremo nel \cref{cap:derivate}: la derivata è il
limite di un rapporto in un punto in cui il rapporto non è definito.
\end{idea}

Nel paragrafo che segue formalizziamo la definizione di limite di funzione secondo le
idee sopra esposte.

%% ----------------------------------------------------------------------------
%%  §4.2  Definizioni: intervalli, intorni, punti di accumulazione, limiti
%% ----------------------------------------------------------------------------
\section{Definizioni}
\label{sec:def-limite-funzioni}

\subsection{Intervalli e intorni}

Gli insiemi su cui definiremo le nostre funzioni sono quasi sempre intervalli o unioni
di intervalli. Fissiamo quindi la terminologia, che finora abbiamo usato in modo
informale.

\begin{definizione}[Intervalli]\label{def:intervalli}\index{intervallo}
Se $a,b$ sono due numeri reali con $a<b$, per indicare un \emph{intervallo} di estremi
$a$, $b$ si usano le notazioni
\begin{align}
  [a,b]&=\{x\in\R : a\le x\le b\}, & (a,b)&=\{x\in\R : a<x<b\}, \label{eq:intervalli-1}\\
  [a,b)&=\{x\in\R : a\le x<b\}, & (a,b]&=\{x\in\R : a<x\le b\}. \label{eq:intervalli-2}
\end{align}
L'intervallo $[a,b]$ si dice \emph{chiuso}, mentre $(a,b)$ è detto \emph{aperto};
$[a,b)$ è detto \emph{chiuso a sinistra} e \emph{aperto a destra} (analogamente per
$(a,b]$). Gli intervalli sopra scritti si dicono \emph{limitati}. Si considerano anche
gli \emph{intervalli illimitati}
\begin{equation}\label{eq:intervalli-illimitati}
\begin{aligned}
  [a,+\infty)&=\{x\in\R : x\ge a\}, & (a,+\infty)&=\{x\in\R : x>a\},\\
  (-\infty,b]&=\{x\in\R : x\le b\}, & (-\infty,b)&=\{x\in\R : x<b\},\\
  (-\infty,+\infty)&=\R .
\end{aligned}
\end{equation}
I numeri $a$, $b$ sono detti \emph{estremi} dell'intervallo, anche nel caso in cui non
ne facciano parte.
\end{definizione}

Si noti che i simboli $+\infty$ e $-\infty$ \emph{non} sono numeri reali: la scrittura
$(a,+\infty)$ è soltanto un modo compatto per indicare la semiretta dei numeri maggiori
di $a$. Talvolta conviene però aggiungere formalmente a $\R$ i due simboli e considerare
la \emph{retta reale ampliata} $\overline{\R}=\R\cup\{-\infty,+\infty\}$, con la
convenzione $-\infty<x<+\infty$ per ogni $x\in\R$; ciò permette di enunciare in un colpo
solo definizioni che altrimenti andrebbero ripetute caso per caso.

\begin{definizione}[Intorni]\label{def:intorni}\index{intorno}
Sia $x_{0}\in\R$. Per ogni $\delta>0$ l'intervallo aperto $(x_{0}-\delta,x_{0}+\delta)$
si chiama \emph{intorno} di $x_{0}$ di raggio $\delta$ \eqref{eq:intorno}; più in
generale, si chiama intorno di $x_{0}$ ogni insieme che contiene un intervallo aperto
contenente $x_{0}$. Gli intervalli $[x_{0},x_{0}+\delta)$ e $(x_{0}-\delta,x_{0}]$ si
chiamano rispettivamente \emph{intorno destro} e \emph{intorno sinistro} di $x_{0}$.
Si chiama \emph{intorno di $+\infty$} ogni semiretta $(k,+\infty)$, e \emph{intorno di
$-\infty$} ogni semiretta $(-\infty,k)$, con $k\in\R$.
\end{definizione}

Due proprietà degli intorni, evidenti ma usate di continuo, sono le seguenti:
l'intersezione di due intorni di $x_{0}$ è ancora un intorno di $x_{0}$ (basta prendere
il più piccolo dei due raggi); due punti distinti $x_{0}\neq x_{1}$ hanno intorni
disgiunti (basta prendere come raggio comune $\abs{x_{1}-x_{0}}/2$).

\subsection{Punti di accumulazione}

Per parlare del limite di $f(x)$ per $x\to x_{0}$ dobbiamo poterci avvicinare a $x_{0}$
restando nel dominio di $f$, e senza mai toccare $x_{0}$. La nozione che formalizza
questa esigenza è la seguente.

\begin{definizione}[Punto di accumulazione]\label{def:punto-accumulazione}\index{punto!di accumulazione}
Sia $A\subseteq\R$. Un numero reale $x_{0}$ si dice \emph{punto di accumulazione} per
$A$ se in ogni intorno di $x_{0}$, cioè in ogni insieme $(x_{0}-\delta,x_{0}+\delta)$ con
$\delta>0$, cade almeno un punto di $A$ \emph{distinto da $x_{0}$}. Si dice che $+\infty$
(rispettivamente $-\infty$) è punto di accumulazione per $A$ se $A$ non è limitato
superiormente (rispettivamente inferiormente), cioè se ogni intorno di $+\infty$
(rispettivamente di $-\infty$) contiene punti di $A$.
\end{definizione}

\begin{marginfigure}
\centering
\begin{tikzpicture}[x=4cm]
  \draw[grigio, ->] (-0.08,0) -- (1.12,0);
  \foreach \n in {1,...,40} { \fill[blu] ({1/\n},0) circle (0.9pt); }
  \draw[rame, thick] (0,-0.12) -- (0,0.12);
  \node[below, font=\scriptsize] at (0,-0.1) {$0$};
  \node[below, font=\scriptsize] at (1,-0.06) {$1$};
  \node[below, font=\scriptsize] at (0.5,-0.06) {$\frac12$};
  \node[below, font=\scriptsize] at (0.3333,-0.06) {$\frac13$};
\end{tikzpicture}
\caption{L'insieme $\{1/n : n\in\N\}$: i suoi punti si addensano verso $0$, che è
l'unico punto di accumulazione e non appartiene all'insieme; ogni punto $1/n$ è invece
isolato.}
\label{fig:accumulazione}
\end{marginfigure}

\begin{esempio}\label{es:accumulazione}
\begin{enumerate}[label=\textup{(\roman*)}]
  \item Sia $A=(0,1]$. Ogni $x_{0}\in[0,1]$ è punto di accumulazione per $A$: se
  $x_{0}=0$, l'intorno $(-\delta,\delta)$ contiene il punto $\min\{\delta/2,1\}\in A$,
  che è diverso da $0$; se $0<x_{0}\le1$, l'intorno $(x_{0}-\delta,x_{0}+\delta)$
  contiene il punto $x_{0}-\frac12\min\{\delta,x_{0}\}$, che sta in $A$ ed è diverso da
  $x_{0}$. Nessun altro numero è punto di accumulazione: se $x_{0}>1$ l'intorno di
  raggio $x_{0}-1$ non incontra $A$, e se $x_{0}<0$ non lo incontra l'intorno di raggio
  $\abs{x_{0}}$. Si noti che $0$ è di accumulazione pur non appartenendo ad $A$.
  \item L'insieme $\N$ non ha punti di accumulazione reali: l'intorno di raggio $\frac12$
  di un qualsiasi numero reale $x_{0}$ contiene al più un numero naturale $m$; se
  $m\neq x_{0}$, l'intorno di raggio $\min\{\frac12,\abs{m-x_{0}}\}$ non contiene alcun
  punto di $\N$ diverso da $x_{0}$. Invece $+\infty$ è
  punto di accumulazione per $\N$, che non è limitato superiormente (\cref{teo:archimede}).
  \item L'insieme $A=\{1/n : n\in\N\}$ ha come unico punto di accumulazione $0$
  (\cref{fig:accumulazione}). Ogni punto $1/n$ di $A$ è invece \emph{isolato}: l'intorno
  di $1/n$ di raggio $\frac1n-\frac1{n+1}$ non contiene altri punti di $A$.
  \item Se $A=\R\setminus\{0\}$, ogni numero reale, compreso $0$, è punto di accumulazione
  per $A$; inoltre lo sono anche $+\infty$ e $-\infty$.
\end{enumerate}
\end{esempio}

Il legame con le successioni è dato dalla proposizione seguente, che sarà la chiave di
tutto il capitolo.

\begin{proposizione}\label{prop:caratt-accumulazione}
Siano $A\subseteq\R$ e $x_{0}\in\overline{\R}$. Allora $x_{0}$ è punto di accumulazione per
$A$ se e solo se esiste una successione $x_{n}$ di punti di $A$, con $x_{n}\neq x_{0}$
per ogni $n\in\N$, che tende a $x_{0}$.
\end{proposizione}

\begin{schema}
\Ipotesi{$A\subseteq\R$, \ $x_{0}\in\R$ (oppure $x_{0}=\pm\infty$).}
\Tesi{$x_{0}$ di accumulazione per $A$ \ $\sse$ \ $\esiste x_{n}\in A\setminus\{x_{0}\}$ con
$x_{n}\to x_{0}$.}
\end{schema}

\begin{strategia}
Per costruire la successione basta usare la definizione con una famiglia di intorni che
si restringono: gli intorni di raggio $\frac1n$. In ciascuno di essi c'è un punto di $A$
diverso da $x_{0}$, e lo prendiamo come $x_{n}$; i carabinieri faranno il resto.
Viceversa, se una successione di punti di $A$ diversi da $x_{0}$ converge a $x_{0}$,
ogni intorno di $x_{0}$ ne contiene definitivamente i termini.
\end{strategia}

\begin{dimostrazione}
Sia $x_{0}\in\R$ di accumulazione per $A$. Per ogni $n\in\N$ l'intorno di raggio
$\frac1n$ contiene un punto di $A$ diverso da $x_{0}$; scegliamone uno e chiamiamolo
$x_{n}$. Allora $x_{n}\in A$, $x_{n}\neq x_{0}$ e
\[
  x_{0}-\frac1n<x_{n}<x_{0}+\frac1n \qquad\perogni n\in\N ;
\]
poiché $\frac1n\to0$, per il teorema dei carabinieri (\cref{teo:carabinieri}) si ha
$x_{n}\to x_{0}$. Se $x_{0}=+\infty$, l'insieme $A$ non è limitato superiormente, e per
ogni $n$ esiste $x_{n}\in A$ con $x_{n}>n$; per il confronto (\cref{teo:confronto-infinito})
$x_{n}\to+\infty$. Il caso $x_{0}=-\infty$ è analogo.

Viceversa, sia $x_{n}\in A\setminus\{x_{0}\}$ con $x_{n}\to x_{0}\in\R$, e fissiamo
$\delta>0$. Per la definizione di limite esiste $\nu$ tale che $\abs{x_{n}-x_{0}}<\delta$
per $n>\nu$: il punto $x_{\nu+1}$, ad esempio, sta in $A$, è diverso da $x_{0}$ e
appartiene all'intorno $(x_{0}-\delta,x_{0}+\delta)$. Se $x_{0}=+\infty$, per ogni $k$ si
ha definitivamente $x_{n}>k$, e quindi la semiretta $(k,+\infty)$ contiene punti di $A$.
\end{dimostrazione}

Nel seguito del capitolo considereremo, come il testo di riferimento, soprattutto
funzioni il cui dominio $A$ è un intervallo (o un'unione finita di intervalli), e punti
$x_{0}$ che appartengono ad $A$ oppure sono estremi di uno degli intervalli che lo
compongono. Ad esempio, se $f$ è definita nell'insieme $A=\R\setminus\{0\}=(-\infty,0)\cup
(0,+\infty)$, il punto $x_{0}$ prescelto per il calcolo del limite può appartenere a uno
dei due intervalli oppure essere l'estremo $0$: in questo esempio $x_{0}$ può essere un
qualunque numero reale. In tutti questi casi $x_{0}$ è punto di accumulazione per $A$,
per l'argomento dell'\cref{es:accumulazione}(i).

\subsection{Limite finito per \texorpdfstring{$x\to x_{0}$}{x che tende a x0}}

\begin{definizione}[Limite di una funzione]\label{def:limite-funzione}\index{limite!di una funzione}
Siano $f:A\to\R$ una funzione, $x_{0}\in\R$ un punto di accumulazione per $A$ e
$\ell\in\R$. Si dice che $f(x)$ \emph{ha limite uguale a $\ell$} (o \emph{tende}, o
\emph{converge}, ad $\ell$) per $x$ che tende a $x_{0}$ se, qualunque sia la successione
$x_{n}\to x_{0}$, con $x_{n}\in A$ e $x_{n}\neq x_{0}$ per ogni $n$, risulta
$f(x_{n})\to\ell$. In simboli:
\begin{equation}\label{eq:def-limite-funzione}
  \lim_{x\to x_{0}}f(x)=\ell \quad\sse\quad
  \perogni x_{n}\in A\setminus\{x_{0}\}\ \text{con}\ x_{n}\to x_{0}:\quad f(x_{n})\to\ell .
\end{equation}
Si scrive anche $f(x)\to\ell$ per $x\to x_{0}$.
\end{definizione}

Due osservazioni chiariscono il significato delle ipotesi.

\begin{osservazione}\label{oss:valore-irrilevante}
Il valore $f(x_{0})$, se anche $x_{0}$ appartiene al dominio, \emph{non ha alcun ruolo}
nella definizione: le successioni considerate hanno $x_{n}\neq x_{0}$. Ad esempio, la
funzione che vale $1$ per $x\neq0$ e $5$ per $x=0$ ha limite $1$ per $x\to0$. Il limite
descrive il comportamento di $f$ \emph{vicino} a $x_{0}$, non \emph{in} $x_{0}$; il
confronto tra le due cose sarà il tema del \cref{sec:continue}.
\end{osservazione}

\begin{osservazione}\label{oss:perche-accumulazione}
L'ipotesi che $x_{0}$ sia di accumulazione per $A$ non è un dettaglio tecnico. Se non lo
fosse, per la \cref{prop:caratt-accumulazione} non esisterebbe \emph{nessuna}
successione $x_{n}\in A\setminus\{x_{0}\}$ convergente a $x_{0}$, e la condizione
\eqref{eq:def-limite-funzione} sarebbe soddisfatta, per mancanza di controesempi, da
\emph{ogni} numero $\ell$: il «limite» perderebbe qualunque significato. Grazie a questa
ipotesi, invece, il limite, se esiste, è unico.
\end{osservazione}

\begin{proposizione}[Unicità del limite]\label{prop:unicita-limite-funzioni}\index{limite!unicità}
Se $\displaystyle\lim_{x\to x_{0}}f(x)=\ell$ e $\displaystyle\lim_{x\to x_{0}}f(x)=\ell'$,
allora $\ell=\ell'$.
\end{proposizione}

\begin{schema}
\Ipotesi{$f:A\to\R$; \ $x_{0}$ di accumulazione per $A$; \ $f(x)\to\ell$ e $f(x)\to\ell'$
per $x\to x_{0}$.}
\Tesi{$\ell=\ell'$.}
\end{schema}

\begin{strategia}
Si riduce tutto all'unicità del limite di una successione. Basta disporre di
\emph{una} successione $x_{n}\in A\setminus\{x_{0}\}$ convergente a $x_{0}$: lungo di
essa $f(x_{n})$ deve tendere sia a $\ell$ sia a $\ell'$. Che una tale successione esista
è garantito proprio dall'ipotesi che $x_{0}$ sia di accumulazione.
\end{strategia}

\begin{dimostrazione}
Poiché $x_{0}$ è di accumulazione per $A$, per la \cref{prop:caratt-accumulazione}
esiste almeno una successione $x_{n}\in A\setminus\{x_{0}\}$ con $x_{n}\to x_{0}$. Per
la definizione di limite la successione $f(x_{n})$ converge sia a $\ell$ sia a $\ell'$;
per l'unicità del limite di una successione (\cref{teo:unicita-limite}) è $\ell=\ell'$.
\end{dimostrazione}

\begin{esempio}\label{es:limiti-elementari}
\begin{enumerate}[label=\textup{(\roman*)}]
  \item Per la \eqref{eq:lim-senx-x}, $\displaystyle\lim_{x\to0}\frac{\sen x}{x}=1$.
  \item Per il \cref{teo:lim-sen-cos}, se $x_{n}\to0$ allora $\sen x_{n}\to0$ e
  $\cos x_{n}\to1$; dunque
  \begin{equation}\label{eq:lim-sen-cos-0}
    \lim_{x\to0}\sen x=0, \qquad \lim_{x\to0}\cos x=1 .
  \end{equation}
  \item Se $f(x)=c$ è costante, $\lim_{x\to x_{0}}f(x)=c$; se $f(x)=x$,
  $\lim_{x\to x_{0}}f(x)=x_{0}$: infatti $f(x_{n})=c\to c$, e rispettivamente
  $f(x_{n})=x_{n}\to x_{0}$.
  \item La funzione $f(x)=\frac{x^{2}-4}{x-2}$ non è definita in $x_{0}=2$, che però è di
  accumulazione per il dominio $\R\setminus\{2\}$. Per $x\neq2$ si ha
  $f(x)=\frac{(x-2)(x+2)}{x-2}=x+2$; quindi, se $x_{n}\to2$ con $x_{n}\neq2$,
  $f(x_{n})=x_{n}+2\to4$, cioè $\lim_{x\to2}f(x)=4$. La semplificazione è lecita
  proprio perché nella definizione di limite si considerano solo punti $x\neq2$.
\end{enumerate}
\end{esempio}

\subsection{La definizione con \texorpdfstring{$\epsilon$ e $\delta$}{epsilon e delta}}

La \cref{def:limite-funzione} è comoda per ereditare i risultati sulle successioni, ma
in essa compare un quantificatore «per ogni successione» che non è facile da verificare
direttamente. Possiamo formulare la definizione di limite per mezzo di disuguaglianze,
come già fatto per le successioni con i simboli $\epsilon$, $\nu$. I simboli usati
classicamente per i limiti di funzioni sono $\epsilon$ e $\delta$ (delta).

\begin{teorema}[Caratterizzazione con $\epsilon$ e $\delta$]\label{teo:eps-delta}
Siano $f:A\to\R$, $x_{0}\in\R$ di accumulazione per $A$ e $\ell\in\R$. Si ha
$\lim_{x\to x_{0}}f(x)=\ell$ se e soltanto se, qualunque sia $\epsilon>0$, esiste un
numero $\delta>0$ tale che $\ell-\epsilon<f(x)<\ell+\epsilon$ per ogni $x\in A$, con
$x\neq x_{0}$ e $x_{0}-\delta<x<x_{0}+\delta$. In simboli:
\begin{equation}\label{eq:eps-delta}
\begin{split}
  \lim_{x\to x_{0}}f(x)=\ell \quad\sse\quad
  &\perogni\epsilon>0\ \esiste\delta>0 :\\
  &\abs{f(x)-\ell}<\epsilon \quad\perogni x\in A,\ 0<\abs{x-x_{0}}<\delta .
\end{split}
\end{equation}
\end{teorema}

Il teorema è dimostrato nel \cref{sec:legame-successioni}. La \cref{fig:eps-delta} ne
illustra il significato: comunque si fissi una striscia orizzontale di semiampiezza
$\epsilon$ attorno alla retta $y=\ell$, si può trovare una striscia verticale di
semiampiezza $\delta$ attorno alla retta $x=x_{0}$ tale che la parte di grafico contenuta
nella striscia verticale (escluso al più il punto di ascissa $x_{0}$) sta tutta dentro
la striscia orizzontale.

\begin{figure}[htb]
\centering
\begin{tikzpicture}
\begin{axis}[width=0.9\textwidth, height=6cm, xmin=0, xmax=4.2, ymin=0, ymax=3.2,
  axis lines=left, xtick={2}, xticklabels={$x_{0}$}, ytick={1.8}, yticklabels={$\ell$},
  xlabel={$x$}, ylabel={$y$}, clip=false]
  \fill[verde!15] (axis cs:0,1.45) rectangle (axis cs:4.2,2.15);
  \fill[blu!12] (axis cs:1.62,0) rectangle (axis cs:2.38,3.2);
  \addplot[tratteggio, domain=0:4.2] {1.8};
  \addplot[curva, domain=0.2:1.97, samples=80] {1.8+0.9*(x-2)+0.12*(x-2)^3*sin(deg(5*x))};
  \addplot[curva, domain=2.03:4, samples=80] {1.8+0.9*(x-2)+0.12*(x-2)^3*sin(deg(5*x))};
  \draw[blu, thick, fill=white] (axis cs:2,1.8) circle (2pt);
  \draw[<->, verde!50!black] (axis cs:3.95,1.8) -- node[right, font=\scriptsize] {$\epsilon$} (axis cs:3.95,2.15);
  \draw[<->, verde!50!black] (axis cs:3.95,1.45) -- node[right, font=\scriptsize] {$\epsilon$} (axis cs:3.95,1.8);
  \draw[<->, blu] (axis cs:1.62,0.25) -- node[below, font=\scriptsize] {$\delta$} (axis cs:2,0.25);
  \draw[<->, blu] (axis cs:2,0.25) -- node[below, font=\scriptsize] {$\delta$} (axis cs:2.38,0.25);
\end{axis}
\end{tikzpicture}
\caption{Il significato della \eqref{eq:eps-delta}: fissata la striscia orizzontale
$\ell-\epsilon<y<\ell+\epsilon$ (verde), esiste una striscia verticale
$x_{0}-\delta<x<x_{0}+\delta$ (azzurra) in cui il grafico, escluso al più il punto di
ascissa $x_{0}$, resta dentro la striscia orizzontale. Se si restringe $\epsilon$, in
generale bisogna restringere anche $\delta$.}
\label{fig:eps-delta}
\end{figure}

\begin{attenzione}[$\delta$ dipende da $\epsilon$]
Come per le successioni, l'ordine dei quantificatori è essenziale: prima si fissa
$\epsilon$, poi si trova $\delta$, che in generale dipende da $\epsilon$ (e da $x_{0}$).
Scambiando l'ordine, «esiste $\delta$ tale che per ogni $\epsilon$\dots», si otterrebbe
che $\abs{f(x)-\ell}$ è minore di ogni numero positivo, cioè $f(x)=\ell$ per tutti gli
$x\neq x_{0}$ di un intorno: una richiesta molto più forte, che non ha nulla a che fare
con il concetto di limite.
\end{attenzione}

\begin{esempio}\label{es:x-quadro-eps-delta}
Verifichiamo con la \eqref{eq:eps-delta} che $\lim_{x\to3}x^{2}=9$. Dobbiamo rendere
piccola la quantità $\abs{x^{2}-9}=\abs{x-3}\cdot\abs{x+3}$. Il primo fattore è
piccolo quando $x$ è vicino a $3$; il secondo non è piccolo, ma è \emph{limitato} se
$x$ resta vicino a $3$: se $\abs{x-3}<1$, cioè $2<x<4$, si ha $\abs{x+3}<7$. Fissato
$\epsilon>0$ poniamo quindi
\[
  \delta=\min\Bigl\{1,\frac\epsilon7\Bigr\}.
\]
Se $0<\abs{x-3}<\delta$, allora $\abs{x-3}<1$, quindi $\abs{x+3}<7$, e
$\abs{x^{2}-9}<7\abs{x-3}<7\cdot\frac\epsilon7=\epsilon$. La scelta del minimo tra due
valori è un artificio tipico: il primo vincolo serve a tenere sotto controllo i fattori
«innocui», il secondo a rendere piccolo il fattore che tende a zero.
\end{esempio}

\begin{idea}[Tolleranze in ingresso e in uscita]
La definizione con $\epsilon$ e $\delta$ ha un'interpretazione naturale per chi progetta.
Si pensi a un sistema che a una grandezza di ingresso $x$ (la posizione della manetta,
l'angolo di un'aletta di compensazione) associa una grandezza di uscita $f(x)$ (la
spinta, il momento di cerniera). Il committente fissa la \emph{tolleranza in uscita}
$\epsilon$: l'uscita deve differire dal valore nominale $\ell$ meno di $\epsilon$.
Il progettista deve allora garantire una \emph{tolleranza in ingresso} $\delta$
abbastanza stretta: se l'ingresso differisce dal valore nominale $x_{0}$ meno di
$\delta$, l'uscita rispetta la specifica. Che $f(x)\to\ell$ per $x\to x_{0}$ significa
esattamente che \emph{qualunque} specifica $\epsilon$ può essere soddisfatta con un'opportuna
precisione $\delta$ sull'ingresso; e, come nell'\cref{es:x-quadro-eps-delta},
specifiche più severe richiedono in generale lavorazioni più precise.
\end{idea}

\subsection{Limiti infiniti e limiti all'infinito}

Valgono analoghe definizioni quando il limite $\ell$, oppure il punto $x_{0}$, o entrambi,
sono infiniti: basta sostituire nella \cref{def:limite-funzione} le parole «converge a»
con «tende a $\pm\infty$», secondo le definizioni del \cref{cap:successioni}. Così, ad
esempio, se $x_{0}\in\R$ è di accumulazione per $A$, e se $A$ non è limitato
superiormente,
\begin{align}
  \lim_{x\to x_{0}}f(x)=+\infty \quad&\sse\quad \perogni x_{n}\in A\setminus\{x_{0}\},\ x_{n}\to x_{0}:\ f(x_{n})\to+\infty;
  \label{eq:def-lim-piu-inf}\\
  \lim_{x\to+\infty}f(x)=\ell \quad&\sse\quad \perogni x_{n}\in A,\ x_{n}\to+\infty:\ f(x_{n})\to\ell;
  \label{eq:def-lim-all-inf}\\
  \lim_{x\to+\infty}f(x)=+\infty \quad&\sse\quad \perogni x_{n}\in A,\ x_{n}\to+\infty:\ f(x_{n})\to+\infty .
  \label{eq:def-lim-inf-inf}
\end{align}
Il lettore formuli i casi corrispondenti con $-\infty$ al posto di $+\infty$. Anche
queste definizioni ammettono una formulazione con le disuguaglianze, analoga alla
\eqref{eq:eps-delta}: la riassumiamo nella \cref{tab:intorni}. In ogni caso la struttura
è la stessa:
\begin{equation}\label{eq:def-intorni}
\begin{gathered}
  \text{per ogni intorno } V \text{ di } \ell \ \ \text{esiste un intorno } U \text{ di } x_{0} :\\
  f(x)\in V \quad\perogni x\in(U\cap A)\setminus\{x_{0}\},
\end{gathered}
\end{equation}
dove gli intorni sono quelli della \cref{def:intorni}, e per i punti infiniti la
condizione $x\neq x_{0}$ è automatica.

\begin{table}[htb]
\centering
\small
\renewcommand{\arraystretch}{1.35}
\begin{tabular}{l l}
\toprule
punto $x_{0}$ & «$x$ vicino a $x_{0}$»\\
\midrule
$x_{0}\in\R$ & $\esiste\delta>0$: \ $0<\abs{x-x_{0}}<\delta$\\
$x_{0}=+\infty$ & $\esiste k$: \ $x>k$\\
$x_{0}=-\infty$ & $\esiste k$: \ $x<k$\\
\bottomrule
\end{tabular}
\qquad
\begin{tabular}{l l}
\toprule
limite $\ell$ & «$f(x)$ vicino a $\ell$»\\
\midrule
$\ell\in\R$ & $\perogni\epsilon>0$: \ $\abs{f(x)-\ell}<\epsilon$\\
$\ell=+\infty$ & $\perogni M>0$: \ $f(x)>M$\\
$\ell=-\infty$ & $\perogni M>0$: \ $f(x)<-M$\\
\bottomrule
\end{tabular}
\caption{Le nove definizioni di limite con le disuguaglianze si ottengono combinando una
riga della tabella di sinistra con una riga della tabella di destra secondo lo schema
\eqref{eq:def-intorni}: «per ogni [condizione su $f(x)$] esiste [condizione su $x$], con
$x\in A$». Ad esempio, $\lim_{x\to+\infty}f(x)=\ell$ significa: per ogni $\epsilon>0$
esiste $k$ tale che $\abs{f(x)-\ell}<\epsilon$ per ogni $x\in A$ con $x>k$.}
\label{tab:intorni}
\end{table}

\begin{esempio}\label{es:limiti-infiniti}
\begin{enumerate}[label=\textup{(\roman*)}]
  \item $\displaystyle\lim_{x\to0}\frac1{x^{2}}=+\infty$. Con le successioni: se
  $x_{n}\to0$ e $x_{n}\neq0$, allora $x_{n}^{2}\to0$ con $x_{n}^{2}>0$, e quindi
  $1/x_{n}^{2}\to+\infty$ (per la \eqref{eq:inf7} e l'avviso «$1/0$ non è $+\infty$»
  del \cref{sec:forme-indeterminate}). Con le
  disuguaglianze: fissato $M>0$, si ha $\frac1{x^{2}}>M$ se $0<\abs x<\delta=\frac1{\sqrt M}$.
  \item $\displaystyle\lim_{x\to+\infty}\frac1x=0$: fissato $\epsilon>0$, per $x>k=\frac1\epsilon$
  si ha $0<\frac1x<\epsilon$.
  \item $\displaystyle\lim_{x\to+\infty}x^{2}=+\infty$: fissato $M>0$, per $x>k=\sqrt M$
  si ha $x^{2}>M$.
  \item Il limite di $\frac1x$ per $x\to0$ \emph{non esiste}: lungo $x_{n}=\frac1n$ si ha
  $\frac1{x_{n}}=n\to+\infty$, lungo $x_{n}'=-\frac1n$ si ha $\frac1{x_{n}'}=-n\to-\infty$,
  mentre la definizione richiede lo stesso limite lungo tutte le successioni.
\end{enumerate}
\end{esempio}

\subsection{Limite destro e limite sinistro}

L'ultimo esempio suggerisce di distinguere da quale parte ci si avvicina a $x_{0}$. È utile
considerare i cosiddetti \emph{limite destro} ($x\to x_{0}^{+}$) e \emph{limite
sinistro} ($x\to x_{0}^{-}$), quando ci si avvicina al punto $x_{0}$ per valori di
$x\in A$ rispettivamente solo maggiori di $x_{0}$, o solo minori.

\begin{definizione}[Limite destro e sinistro]\label{def:limite-dx-sx}\index{limite!destro e sinistro}
Sia $x_{0}\in\R$ di accumulazione per $A\cap(x_{0},+\infty)$, e sia $\ell\in\R$. Si dice
che $\ell$ è il \emph{limite destro} di $f$ in $x_{0}$ se
\begin{equation}\label{eq:def-limite-dx}
\begin{aligned}
  \lim_{x\to x_{0}^{+}}f(x)=\ell \quad&\sse\quad \perogni x_{n}\to x_{0},\ x_{n}\in A,\ x_{n}>x_{0}:\ f(x_{n})\to\ell\\
  &\sse\quad \perogni\epsilon>0\ \esiste\delta>0 :\ \abs{f(x)-\ell}<\epsilon\\
  &\phantom{\sse}\quad \perogni x\in A,\ x_{0}<x<x_{0}+\delta .
\end{aligned}
\end{equation}
Analogamente, se $x_{0}$ è di accumulazione per $A\cap(-\infty,x_{0})$,
\begin{equation}\label{eq:def-limite-sx}
\begin{aligned}
  \lim_{x\to x_{0}^{-}}f(x)=\ell \quad&\sse\quad \perogni x_{n}\to x_{0},\ x_{n}\in A,\ x_{n}<x_{0}:\ f(x_{n})\to\ell\\
  &\sse\quad \perogni\epsilon>0\ \esiste\delta>0 :\ \abs{f(x)-\ell}<\epsilon\\
  &\phantom{\sse}\quad \perogni x\in A,\ x_{0}-\delta<x<x_{0} .
\end{aligned}
\end{equation}
Il lettore formuli i casi con limite infinito.
\end{definizione}

In altre parole, il limite destro di $f$ in $x_{0}$ non è altro che il limite in $x_{0}$
della restrizione di $f$ all'insieme $A\cap(x_{0},+\infty)$: l'equivalenza tra le due
righe della \eqref{eq:def-limite-dx} è quindi un caso particolare del \cref{teo:eps-delta}.

\begin{marginfigure}
\centering
\begin{tikzpicture}[scale=0.9]
  \draw[->, grigio] (-2,0) -- (2.2,0) node[below, font=\scriptsize, black] {$x$};
  \draw[->, grigio] (0,-1.5) -- (0,1.6) node[left, font=\scriptsize, black] {$y$};
  \draw[blu, very thick] (0,1) -- (1.9,1);
  \draw[blu, very thick] (-1.9,-1) -- (0,-1);
  \draw[blu, thick, fill=white] (0,1) circle (2pt);
  \draw[blu, thick, fill=white] (0,-1) circle (2pt);
  \node[left, font=\scriptsize] at (0,1) {$1$};
  \node[right, font=\scriptsize] at (0,-1) {$-1$};
\end{tikzpicture}
\caption{La funzione $\frac{\abs x}{x}$: limite destro $1$ e limite sinistro $-1$ in
$x_{0}=0$; il limite (bilatero) non esiste.}
\label{fig:modulo-su-x}
\end{marginfigure}

\begin{esempio}\label{es:modulo-su-x}
Per la funzione $f(x)=\frac{\abs x}{x}$, definita per $x\neq0$, si ha $f(x)=1$ per
$x>0$ e $f(x)=-1$ per $x<0$ (\cref{fig:modulo-su-x}). Quindi
\begin{equation}\label{eq:modulo-su-x}
  \lim_{x\to0^{+}}\frac{\abs x}{x}=1, \qquad \lim_{x\to0^{-}}\frac{\abs x}{x}=-1,
  \qquad \lim_{x\to0}\frac{\abs x}{x}\ \text{non esiste}:
\end{equation}
la successione $x_{n}=\frac{(-1)^{n}}n$ tende a zero, ma $f(x_{n})=(-1)^{n}$ non è
regolare. Analogamente $\lim_{x\to0^{+}}\frac1x=+\infty$ e $\lim_{x\to0^{-}}\frac1x=-\infty$.
\end{esempio}

\begin{notastorica}[Due definizioni, una sola idea]
La definizione di limite con le disuguaglianze si è affermata attraverso le lezioni
berlinesi di Karl Weierstrass, negli anni Sessanta dell'Ottocento, ma la sostanza era già
nel \emph{Résumé des leçons} di Cauchy (1823), dove le lettere $\epsilon$ e $\delta$
compaiono per misurare l'incremento della variabile e lo scarto del rapporto
incrementale dalla derivata. La definizione
«per successioni» porta invece il nome di Eduard Heine, autore degli \emph{Elemente der
Functionenlehre} (1872), un'esposizione dei fondamenti dell'analisi scritta sotto
l'influenza delle lezioni di Weierstrass e in dialogo con Georg Cantor, nella quale
anche i numeri reali sono costruiti a partire dalle successioni. Si parla così di
\emph{definizione di Heine} (sequenziale) e di \emph{definizione di Cauchy}
($\epsilon$-$\delta$).
Il fatto che le due definizioni siano equivalenti, che dimostreremo tra poco, non è
banale: la parte delicata richiede di \emph{scegliere} infiniti punti, uno per ogni
$n$, e sottintende un principio logico (una forma debole dell'assioma della scelta) che
fu messo in luce solo all'inizio del Novecento.
\end{notastorica}

%% ----------------------------------------------------------------------------
%%  §4.3  Legame tra limiti di funzioni e limiti di successioni
%% ----------------------------------------------------------------------------
\section{Legame tra limiti di funzioni e limiti di successioni}
\label{sec:legame-successioni}

In questo paragrafo dimostriamo il \cref{teo:eps-delta}, cioè l'equivalenza tra la
definizione di limite «per successioni» (\cref{def:limite-funzione}) e quella con le
disuguaglianze. Fissati $f:A\to\R$, un punto $x_{0}\in\R$ di accumulazione per $A$ e
$\ell\in\R$, si tratta di provare che le due relazioni
\begin{align}
  &\perogni x_{n}\to x_{0},\quad x_{n}\in A\setminus\{x_{0}\}\ \perogni n\in\N
  \quad\implica\quad f(x_{n})\to\ell; \label{eq:ponte-succ}\\
  &\perogni\epsilon>0\ \esiste\delta>0 :\quad x\in A,\ \ 0<\abs{x-x_{0}}<\delta
  \quad\implica\quad \abs{f(x)-\ell}<\epsilon \label{eq:ponte-eps}
\end{align}
sono fra loro equivalenti.

\begin{schema}
\Ipotesi{$f:A\to\R$; \ $x_{0}\in\R$ di accumulazione per $A$; \ $\ell\in\R$.}
\Tesi{\eqref{eq:ponte-succ} \ $\sse$ \ \eqref{eq:ponte-eps}.}
\end{schema}

\begin{strategia}
Le due implicazioni hanno natura molto diversa. Che la \eqref{eq:ponte-eps} implichi la
\eqref{eq:ponte-succ} è quasi immediato: se una successione $x_{n}$ tende a $x_{0}$, i
suoi termini finiscono per stare nell'intorno di raggio $\delta$, e lì la
\eqref{eq:ponte-eps} controlla i valori di $f$. Il viceversa è più sottile, perché
dalla conoscenza di $f$ lungo le successioni dobbiamo ricavare un'informazione su
\emph{tutti} i punti di un intorno. Lo dimostriamo per assurdo: se la
\eqref{eq:ponte-eps} fosse falsa, esisterebbe una tolleranza $\epsilon_{0}$ violata in
ogni intorno di $x_{0}$, per quanto piccolo; scegliendo un punto «cattivo» in ciascuno
degli intorni di raggio $\frac1n$ costruiremo una successione che tende a $x_{0}$ lungo
la quale $f$ \emph{non} tende a $\ell$, contro la \eqref{eq:ponte-succ}.
\end{strategia}

\begin{dimostrazione}[Dimostrazione del \cref{teo:eps-delta}]
\emph{La \eqref{eq:ponte-eps} implica la \eqref{eq:ponte-succ}.} Consideriamo una
generica successione $x_{n}$ di punti di $A$, convergente a $x_{0}$, con $x_{n}\neq x_{0}$
per ogni $n\in\N$, e fissiamo $\epsilon>0$; sia $\delta>0$ il numero per cui vale la
\eqref{eq:ponte-eps}. Per la definizione di limite di successione (applicata con $\delta$
al posto di $\epsilon$) esiste un indice $\nu$ per cui $\abs{x_{n}-x_{0}}<\delta$ per ogni
$n>\nu$; inoltre, essendo $x_{n}\neq x_{0}$, in definitiva si ha
\begin{equation}\label{eq:ponte-passo1}
  x_{n}\in A, \qquad 0<\abs{x_{n}-x_{0}}<\delta \qquad\perogni n>\nu .
\end{equation}
Per l'ipotesi \eqref{eq:ponte-eps}, applicata con $x=x_{n}$, segue allora
\begin{equation}\label{eq:ponte-passo2}
  \abs{f(x_{n})-\ell}<\epsilon \qquad\perogni n>\nu,
\end{equation}
che, in base alla definizione di limite di successione, significa $f(x_{n})\to\ell$ per
$n\to+\infty$.

\emph{La \eqref{eq:ponte-succ} implica la \eqref{eq:ponte-eps}.} Procediamo per assurdo.
La \eqref{eq:ponte-eps} ha la forma «per ogni $\epsilon$ esiste $\delta$ tale che per
ogni $x$ vale una certa proprietà»; la sua negazione si ottiene scambiando i
quantificatori «per ogni» ed «esiste» e negando la proprietà finale. Contraddire la
\eqref{eq:ponte-eps} equivale dunque ad affermare che
\begin{equation}\label{eq:ponte-negazione}
\begin{split}
  \esiste\epsilon_{0}>0 :\ \perogni\delta>0\ \esiste x\in A :\quad &0<\abs{x-x_{0}}<\delta,\\
  &\abs{f(x)-\ell}\ge\epsilon_{0} .
\end{split}
\end{equation}
Poniamo $\delta=\frac1n$, con $n\in\N$, e indichiamo con $x=x_{n}$ un punto che compare
nella \eqref{eq:ponte-negazione} in corrispondenza di $\delta=\frac1n$. Otteniamo
\begin{equation}\label{eq:ponte-succ-cattiva}
  \esiste\epsilon_{0}>0 :\ \perogni n\in\N\ \esiste x_{n}\in A :\quad
  0<\abs{x_{n}-x_{0}}<\frac1n, \quad \abs{f(x_{n})-\ell}\ge\epsilon_{0} .
\end{equation}
Risulta in particolare
\begin{equation}\label{eq:ponte-carabinieri}
  x_{n}\neq x_{0}, \qquad x_{0}-\frac1n<x_{n}<x_{0}+\frac1n \qquad\perogni n\in\N;
\end{equation}
perciò $x_{n}\in A\setminus\{x_{0}\}$ per ogni $n$ e $x_{n}\to x_{0}$, per il teorema dei
carabinieri. Però $f(x_{n})$ non converge a $\ell$, perché la disuguaglianza
$\abs{f(x_{n})-\ell}\ge\epsilon_{0}$, valida per \emph{ogni} $n\in\N$, contrasta con la
definizione di limite di successione scritta con $\epsilon=\epsilon_{0}$ (non può
esistere alcun $\nu$ oltre il quale $\abs{f(x_{n})-\ell}<\epsilon_{0}$). Ciò contraddice
la \eqref{eq:ponte-succ}.
\end{dimostrazione}

\begin{osservazione}[Il caso dei punti e dei limiti infiniti]\label{oss:ponte-infinito}
La stessa dimostrazione vale, quasi parola per parola, quando $x_{0}$ oppure $\ell$, o
entrambi, sono infiniti, cioè prova l'equivalenza tra le definizioni
\eqref{eq:def-lim-piu-inf}--\eqref{eq:def-lim-inf-inf} e le corrispondenti formulazioni
della \cref{tab:intorni}. Il modo più rapido di vederlo è usare il linguaggio degli
intorni \eqref{eq:def-intorni}, osservando che le definizioni di limite di successione,
finito o infinito, del \cref{cap:successioni} si riassumono in un'unica frase:
\[
  x_{n}\to x_{0}\in\overline{\R} \quad\sse\quad
  \begin{gathered}[t]
    \text{per ogni intorno } U \text{ di } x_{0},\\
    x_{n}\in U \text{ definitivamente.}
  \end{gathered}
\]
Con questa frase, la prima parte della dimostrazione si ripete identica. Nella seconda,
la negazione della \eqref{eq:def-intorni} fornisce un intorno $V_{0}$ di $\ell$ tale che
ogni intorno $U$ di $x_{0}$ contiene un punto $x\in A\setminus\{x_{0}\}$ con
$f(x)\notin V_{0}$; scegliendo come $U$ l'intorno di raggio $\frac1n$ se $x_{0}\in\R$,
la semiretta $(n,+\infty)$ se $x_{0}=+\infty$ e la semiretta $(-\infty,-n)$ se
$x_{0}=-\infty$, si costruisce una successione $x_{n}\to x_{0}$ con $f(x_{n})\notin V_{0}$
per ogni $n$, che quindi non tende a $\ell$. Il lettore scriva per esteso il caso
$x_{0}=+\infty$, $\ell\in\R$ (\cref{es:ponte-infinito}).
\end{osservazione}

\subsection{Il limite è una proprietà locale}

Una conseguenza immediata della caratterizzazione con gli intorni è che il limite di $f$
per $x\to x_{0}$ dipende soltanto dai valori di $f$ \emph{vicino} a $x_{0}$: modificare
$f$ lontano da $x_{0}$ non cambia nulla.

\begin{lemma}[Località del limite]\label{lem:localita}
Siano $f:A\to\R$, $B\subseteq A$, e $x_{0}\in\overline{\R}$ di accumulazione per $B$; sia
$\ell\in\overline{\R}$.
\begin{enumerate}[label=\textup{(\roman*)}]
  \item Se $\lim_{x\to x_{0}}f(x)=\ell$, anche la restrizione di $f$ a $B$ tende a $\ell$
  per $x\to x_{0}$.
  \item Se esiste un intorno $U_{0}$ di $x_{0}$ tale che $U_{0}\cap A\subseteq B$, vale anche il
  viceversa: se la restrizione di $f$ a $B$ tende a $\ell$, allora
  $\lim_{x\to x_{0}}f(x)=\ell$.
\end{enumerate}
\end{lemma}

\begin{schema}
\Ipotesi{$B\subseteq A$; \ $x_{0}$ di accumulazione per $B$; \ (i) $f(x)\to\ell$; \
(ii) la restrizione di $f$ a $B$ tende a $\ell$, e $U_{0}\cap A\subseteq B$ per un
intorno $U_{0}$ di $x_{0}$.}
\Tesi{(i) la restrizione di $f$ a $B$ tende a $\ell$; \ (ii) $f(x)\to\ell$ per
$x\to x_{0}$.}
\end{schema}

\begin{strategia}
Si lavora con la definizione «per successioni». La (i) è immediata, perché le
successioni da controllare per la restrizione sono una parte di quelle da controllare
per $f$. Per la (ii) si osserva che una successione di punti di $A$ che tende a $x_{0}$
entra definitivamente nell'intorno $U_{0}$, e quindi in $B$: tolti i primi termini, che
non influiscono sul limite, è una successione di punti di $B$.
\end{strategia}

\begin{dimostrazione}
(i) Ogni successione di punti di $B\setminus\{x_{0}\}$ che tende a $x_{0}$ è anche una
successione di punti di $A\setminus\{x_{0}\}$: lungo di essa $f$ tende a $\ell$ per
ipotesi.

(ii) Sia $x_{n}\in A\setminus\{x_{0}\}$ con $x_{n}\to x_{0}$. Per la frase
dell'\cref{oss:ponte-infinito}, $x_{n}\in U_{0}$ definitivamente, diciamo per $n>\nu$;
quindi $x_{n}\in B$ per $n>\nu$. La successione traslata $x_{\nu+k}$, $k\in\N$, è
formata da punti di $B\setminus\{x_{0}\}$ e tende a $x_{0}$: per ipotesi
$f(x_{\nu+k})\to\ell$ per $k\to+\infty$. Ne segue $f(x_{n})\to\ell$, perché le due
successioni $f(x_{n})$ e $f(x_{\nu+k})$ differiscono soltanto per i primi $\nu$ termini:
se, ad esempio con $\ell\in\R$, si ha $\abs{f(x_{\nu+k})-\ell}<\epsilon$ per ogni $k>k_{0}$,
allora $\abs{f(x_{n})-\ell}<\epsilon$ per ogni $n>\nu+k_{0}$.
\end{dimostrazione}

La parte (ii) non vale senza l'ipotesi sull'intorno: la funzione $\frac{\abs x}{x}$,
ristretta a $B=(0,+\infty)$, tende a $1$ per $x\to0$, ma su tutto il suo dominio non ha
limite. Il limite destro e il limite sinistro sono proprio i limiti delle restrizioni a
$A\cap(x_{0},+\infty)$ e ad $A\cap(-\infty,x_{0})$; la proposizione seguente dice come
ricostruire da essi il limite.

\begin{proposizione}\label{prop:dx-sx}\index{limite!destro e sinistro}
Sia $x_{0}\in\R$ di accumulazione sia per $A\cap(x_{0},+\infty)$ sia per
$A\cap(-\infty,x_{0})$, e sia $\ell\in\overline{\R}$. Allora
\[
  \lim_{x\to x_{0}}f(x)=\ell \quad\sse\quad
  \lim_{x\to x_{0}^{+}}f(x)=\lim_{x\to x_{0}^{-}}f(x)=\ell .
\]
\end{proposizione}

\begin{schema}
\Ipotesi{$x_{0}\in\R$ di accumulazione per $A\cap(x_{0},+\infty)$ e per
$A\cap(-\infty,x_{0})$; \ $\ell\in\overline{\R}$.}
\Tesi{$f(x)\to\ell$ per $x\to x_{0}$ \ $\sse$ \ limite destro e limite sinistro esistono
e valgono entrambi $\ell$.}
\end{schema}

\begin{strategia}
Limite destro e sinistro sono limiti di restrizioni di $f$, e quindi la parte
«$\implica$» è la (i) del lemma di località. Per il viceversa si usano le disuguaglianze:
ciascuno dei due limiti unilaterali fornisce un $\delta$ dal proprio lato di $x_{0}$, e il
più piccolo dei due va bene da entrambi i lati.
\end{strategia}

\begin{dimostrazione}
L'implicazione «$\implica$» segue dal \cref{lem:localita}(i). Viceversa, supponiamo
$\ell\in\R$ (il caso infinito è analogo, con $M$ al posto di $\epsilon$) e fissiamo
$\epsilon>0$. Per la \eqref{eq:def-limite-dx} esiste $\delta_{1}>0$ tale che
$\abs{f(x)-\ell}<\epsilon$ per $x\in A$, $x_{0}<x<x_{0}+\delta_{1}$; per la
\eqref{eq:def-limite-sx} esiste $\delta_{2}>0$ tale che $\abs{f(x)-\ell}<\epsilon$ per
$x\in A$, $x_{0}-\delta_{2}<x<x_{0}$. Posto $\delta=\min\{\delta_{1},\delta_{2}\}$, per ogni
$x\in A$ con $0<\abs{x-x_{0}}<\delta$ si ha $x$ nel primo o nel secondo caso, e quindi
$\abs{f(x)-\ell}<\epsilon$. Per il \cref{teo:eps-delta}, $f(x)\to\ell$.
\end{dimostrazione}

\subsection{Come si dimostra che un limite non esiste}

La definizione «per successioni» è particolarmente efficace per dimostrare che un limite
\emph{non} esiste: basta esibire due successioni lungo le quali la funzione si comporta
in modo diverso.

\begin{corollario}[Criterio di non esistenza]\label{cor:non-esistenza}
Se esistono due successioni $x_{n}$, $x_{n}'$ di punti di $A\setminus\{x_{0}\}$, entrambe
tendenti a $x_{0}$, tali che $f(x_{n})\to\lambda$ e $f(x_{n}')\to\lambda'$ con
$\lambda\neq\lambda'$, allora il limite di $f(x)$ per $x\to x_{0}$ non esiste. Lo stesso
accade se esiste una sola successione $x_{n}\in A\setminus\{x_{0}\}$, $x_{n}\to x_{0}$,
lungo la quale $f(x_{n})$ non è regolare.
\end{corollario}

\begin{schema}
\Ipotesi{$x_{n},x_{n}'\in A\setminus\{x_{0}\}$, \ $x_{n}\to x_{0}$, $x_{n}'\to x_{0}$, \
$f(x_{n})\to\lambda$, $f(x_{n}')\to\lambda'$, $\lambda\neq\lambda'$ \ (oppure: $f(x_{n})$
non regolare).}
\Tesi{$\lim_{x\to x_{0}}f(x)$ non esiste.}
\end{schema}

\begin{strategia}
È la negazione della \cref{def:limite-funzione}: se il limite esistesse, la funzione
dovrebbe tendere a esso lungo \emph{tutte} le successioni ammissibili. Si ragiona per
assurdo e si conclude con l'unicità del limite di una successione.
\end{strategia}

\begin{dimostrazione}
Se esistesse $\lim_{x\to x_{0}}f(x)=\ell$, per la \cref{def:limite-funzione} sarebbe
$f(x_{n})\to\ell$ e $f(x_{n}')\to\ell$, quindi $\lambda=\ell=\lambda'$ per l'unicità del
limite di successione; nel secondo caso $f(x_{n})$ avrebbe limite $\ell$.
\end{dimostrazione}

\begin{esempio}\label{es:sen-cos-infinito}
I limiti
\begin{equation}\label{eq:sen-cos-infinito}
  \lim_{x\to+\infty}\sen x, \qquad \lim_{x\to+\infty}\cos x \qquad\text{non esistono.}
\end{equation}
Limitiamoci al primo dei due. Se esistesse il limite $\ell$, dovremmo avere che
$\sen x_{n}$ tende sempre allo stesso valore $\ell$ qualunque sia la successione
$x_{n}\to+\infty$. Ma ponendo $x_{n}=2\pi n$ e $x_{n}'=2\pi n+\frac\pi2$, entrambe
divergenti a $+\infty$, risulta $\sen x_{n}=0\to0$, mentre $\sen x_{n}'=1\to1$. Per il
coseno si usano le stesse successioni, scambiate di ruolo.
\end{esempio}

\begin{esempio}\label{es:sen-1-x}
Analogamente si dimostra la \eqref{eq:sen-1-x-anticipo}:
\begin{equation}\label{eq:sen-1-x}
  \lim_{x\to0}\sen\frac1x \qquad\text{non esiste.}
\end{equation}
La successione $x_{n}=\frac1{n\pi}$ converge a zero per $n\to+\infty$; in
corrispondenza la funzione $f(x)=\sen\frac1x$ assume i valori
\begin{equation}\label{eq:sen-1-x-a}
  f(x_{n})=\sen\frac1{x_{n}}=\sen(n\pi)=0 .
\end{equation}
Perciò $f(x_{n})$ è la successione costante uguale a zero, e converge a zero.
Analogamente, posto $x_{n}'=\frac1{\pi/2+2n\pi}$, risulta
\begin{equation}\label{eq:sen-1-x-b}
  f(x_{n}')=\sen\frac1{x_{n}'}=\sen\Bigl(\frac\pi2+2n\pi\Bigr)=1;
\end{equation}
quindi $x_{n}'\to0$, ma $f(x_{n}')\to1$, in contrasto con il fatto che $f(x_{n})\to0$.
Ne segue che $f(x)$ non ammette limite per $x\to0$ (si veda il grafico nella
\cref{fig:sen-1-x}). Con lo stesso argomento si vede che la funzione $g$ dell'avviso
«Una tabella non è una dimostrazione» del \cref{sec:premessa-funzioni} non ha limite per
$x\to0$: lungo $x_{m}=\frac1m$ si ha $\sen\frac\pi{x_{m}}=\sen(m\pi)=0$ e
$g(x_{m})=\frac{\sen x_{m}}{x_{m}}\to1$, mentre lungo $x_{m}'=\frac2{4m+1}$ si ha
$g(x_{m}')=\frac{\sen x_{m}'}{x_{m}'}+1\to2$.
\end{esempio}

Il programma seguente calcola i valori di $\sen\frac1x$ lungo le due successioni
dell'\cref{es:sen-1-x}.

\begin{python}[Due successioni, due limiti diversi]
import math

f = lambda x: math.sin(1 / x)

# x_n = 1/(n pi) e x'_n = 1/(pi/2 + 2 n pi) tendono entrambe a 0
print("      n      x_n        f(x_n)       x'_n        f(x'_n)")
for n in (1, 10, 100, 1000, 10**6):
    x = 1 / (n * math.pi)
    y = 1 / (math.pi / 2 + 2 * n * math.pi)
    print(f"{n:7d}   {x:.4e}  {f(x):+.2e}   {y:.4e}   {f(y):.12f}")
\end{python}
\begin{risultato}
      n      x_n        f(x_n)       x'_n        f(x'_n)
      1   3.1831e-01  +1.22e-16   1.2732e-01   1.000000000000
     10   3.1831e-02  -1.22e-15   1.5527e-02   1.000000000000
    100   3.1831e-03  +1.96e-15   1.5876e-03   1.000000000000
   1000   3.1831e-04  -3.21e-13   1.5912e-04   1.000000000000
1000000   3.1831e-07  -2.23e-10   1.5915e-07   1.000000000000
\end{risultato}

Si noti che il calcolatore non restituisce esattamente $\sen(n\pi)=0$: il numero $\pi$
è rappresentato con un errore relativo dell'ordine di $10^{-16}$, e l'argomento
$1/x_{n}\approx n\pi$ porta con sé un errore assoluto che cresce con $n$ (per
$n=10^{6}$ si arriva a $2\cdot10^{-10}$). È un'ulteriore conferma che la questione
dell'esistenza di un limite non si decide con il calcolo numerico: il programma
\emph{illustra} il ragionamento dell'\cref{es:sen-1-x}, ma non lo sostituisce.

%% ----------------------------------------------------------------------------
%%  §4.4  Esempi e proprietà dei limiti di funzioni
%% ----------------------------------------------------------------------------
\section{Esempi e proprietà dei limiti di funzioni}
\label{sec:proprieta-limiti-funzioni}

Dato che i limiti di funzioni sono definiti a partire dai limiti di successioni, per
essi valgono le proprietà già dimostrate nel \cref{cap:successioni}. In questo paragrafo
ricaviamo alcuni esempi fondamentali e le regole di calcolo: operazioni, confronti,
limiti di funzioni composte.

\subsection{Esponenziali e logaritmi}

Esempi di limiti molto importanti, e che quindi occorre sapere bene, sono i limiti della
funzione esponenziale, che derivano dal limite della successione geometrica
(\cref{teo:lim-geometrica}).

\begin{proposizione}[Limiti dell'esponenziale]\label{prop:lim-esponenziale}\index{funzione!esponenziale!limiti}
Risulta
\begin{equation}\label{eq:lim-esponenziale}
  \lim_{x\to+\infty}a^{x}=
  \begin{cases}
    +\infty & \text{se } a>1,\\
    0 & \text{se } 0<a<1,
  \end{cases}
  \quad
  \lim_{x\to-\infty}a^{x}=
  \begin{cases}
    0 & \text{se } a>1,\\
    +\infty & \text{se } 0<a<1.
  \end{cases}
\end{equation}
In particolare, per la base $e$,
\begin{equation}\label{eq:lim-e-x}
  \lim_{x\to+\infty}e^{x}=+\infty, \qquad \lim_{x\to-\infty}e^{x}=0 .
\end{equation}
\end{proposizione}

\begin{schema}
\Ipotesi{$a>0$, \ $a\neq1$; potenze a esponente reale (\cref{def:potenza-reale}).}
\Tesi{i limiti \eqref{eq:lim-esponenziale}.}
\end{schema}

\begin{strategia}
Conosciamo già il comportamento di $a^{n}$ per $n$ intero; per passare a un esponente
reale qualsiasi basta la monotonia della funzione esponenziale: se $x$ supera l'intero
$\bar n$, allora $a^{x}$ supera $a^{\bar n}$. Useremo la formulazione con le
disuguaglianze della \cref{tab:intorni}. Il caso $0<a<1$ si riconduce al caso $a>1$
attraverso la definizione $a^{x}=(1/a)^{-x}$.
\end{strategia}

\begin{dimostrazione}
Sia $a>1$. Fissato $M>0$, poiché $a^{n}\to+\infty$ (\cref{teo:lim-geometrica}) esiste
$\bar n\in\N$ con $a^{\bar n}>M$. Per ogni $x>\bar n$, la stretta monotonia
dell'esponenziale (\cref{teo:esponenziale}(ii)) dà $a^{x}>a^{\bar n}>M$: è la
definizione di $\lim_{x\to+\infty}a^{x}=+\infty$ con $k=\bar n$. Fissato poi $\epsilon>0$,
poiché $a^{-n}=\frac1{a^{n}}\to0$ (\eqref{eq:inf5}), esiste $\bar n$ con
$a^{-\bar n}<\epsilon$; per ogni $x<-\bar n$ si ha allora $0<a^{x}<a^{-\bar n}<\epsilon$,
cioè $a^{x}\to0$ per $x\to-\infty$.

Sia ora $0<a<1$, e sia $x_{n}\to+\infty$. Allora $-x_{n}\to-\infty$ e, per quanto appena
provato con la base $1/a>1$, $a^{x_{n}}=(1/a)^{-x_{n}}\to0$; per la definizione «per
successioni», $a^{x}\to0$ per $x\to+\infty$. Allo stesso modo, se $x_{n}\to-\infty$,
allora $-x_{n}\to+\infty$ e $a^{x_{n}}=(1/a)^{-x_{n}}\to+\infty$.
\end{dimostrazione}

Il lettore controlli i limiti \eqref{eq:lim-esponenziale} sui grafici della
\cref{fig:esponenziali}; controlli anche i seguenti importanti limiti della funzione
$\log x$ (in base $e$), visibili nella \cref{fig:logaritmi}.

\begin{proposizione}[Limiti del logaritmo]\label{prop:lim-logaritmo}\index{funzione!logaritmo!limiti}
\begin{equation}\label{eq:lim-log}
  \lim_{x\to+\infty}\log x=+\infty, \qquad \lim_{x\to0^{+}}\log x=-\infty .
\end{equation}
\end{proposizione}

\begin{schema}
\Ipotesi{$\log$ definito su $(0,+\infty)$, strettamente crescente, con $\log e^{y}=y$.}
\Tesi{i limiti \eqref{eq:lim-log}.}
\end{schema}

\begin{strategia}
Si usa la formulazione con le disuguaglianze. Per superare una soglia $M$ basta che
$x$ superi il numero di cui $M$ è il logaritmo, cioè $e^{M}$: la monotonia fa il resto.
Allo stesso modo, $\log x<-M$ appena $0<x<e^{-M}$.
\end{strategia}

\begin{dimostrazione}
Il logaritmo è definito su tutto $(0,+\infty)$ (\cref{teo:esponenziale}(iii)), è
strettamente crescente (\cref{prop:log-monotono}) e $\log e^{y}=y$ per ogni $y\in\R$
\eqref{eq:log-inversa}. Fissato $M>0$: se $x>e^{M}$ si ha $\log x>\log e^{M}=M$; se
$0<x<e^{-M}$ si ha $\log x<\log e^{-M}=-M$. Sono le due definizioni richieste, con
$k=e^{M}$ e con $\delta=e^{-M}$ (si osservi che il dominio del logaritmo sta tutto a
destra di $0$, per cui il limite per $x\to0$ e il limite destro coincidono).
\end{dimostrazione}

Altri limiti notevoli, conseguenza della definizione di limite di funzione e dei limiti
di successioni del \cref{teo:e-generalizzato}, sono i seguenti.

\begin{proposizione}\label{prop:e-funzioni}
\begin{equation}\label{eq:e-funzioni}
  \lim_{x\to+\infty}\Bigl(1+\frac1x\Bigr)^{x}=e, \qquad
  \lim_{x\to-\infty}\Bigl(1+\frac1x\Bigr)^{x}=e .
\end{equation}
\end{proposizione}

\begin{schema}
\Ipotesi{$\bigl(1+\frac1{a_{n}}\bigr)^{a_{n}}\to e$ per ogni successione $a_{n}\to\pm\infty$
(\cref{teo:e-generalizzato}).}
\Tesi{i limiti \eqref{eq:e-funzioni}.}
\end{schema}

\begin{strategia}
Il lavoro è già stato fatto nel \cref{cap:successioni}: la definizione «per
successioni» chiede esattamente il comportamento lungo ogni successione divergente di
punti del dominio, ed è ciò che afferma il \cref{teo:e-generalizzato}. Occorre solo
controllare che il dominio non sia limitato, perché $\pm\infty$ sia di accumulazione.
\end{strategia}

\begin{dimostrazione}
La funzione è definita dove $1+\frac1x>0$, cioè per $x<-1$ oppure $x>0$: il suo dominio
non è limitato né superiormente né inferiormente. Se $x_{n}$ è una qualsiasi successione
di punti del dominio che diverge a $+\infty$, la \eqref{eq:e-piu-inf} del
\cref{teo:e-generalizzato}, applicata con $a_{n}=x_{n}$, dà
$\bigl(1+\frac1{x_{n}}\bigr)^{x_{n}}\to e$: è proprio la definizione
\eqref{eq:def-lim-all-inf}. Per $x\to-\infty$ si usa allo stesso modo la
\eqref{eq:e-meno-inf}.
\end{dimostrazione}

Si noti come il \cref{teo:e-generalizzato}, che nel \cref{cap:successioni} sembrava un
complemento tecnico, sia esattamente ciò che serve: la definizione di limite di funzione
richiede il comportamento lungo \emph{tutte} le successioni divergenti, non soltanto
lungo $x_{n}=n$.

\subsection{Operazioni con i limiti}

Il limite della somma, differenza, prodotto, quoziente di due funzioni è
rispettivamente uguale alla somma, differenza, prodotto, quoziente (se il denominatore è
diverso da zero) dei due limiti, purché non si presenti una delle forme indeterminate.
Precisamente:

\begin{teorema}[Operazioni con i limiti di funzioni]\label{teo:operazioni-funzioni}\index{limite!operazioni}
Siano $f,g:A\to\R$, $x_{0}\in\overline{\R}$ di accumulazione per $A$, e supponiamo che
$f(x)\to\ell_{1}$, $g(x)\to\ell_{2}$ per $x\to x_{0}$, con $\ell_{1},\ell_{2}\in\R$. Allora
\begin{equation}\label{eq:operazioni-funzioni}
\begin{gathered}
  f(x)\pm g(x)\to\ell_{1}\pm\ell_{2}, \qquad f(x)\,g(x)\to\ell_{1}\ell_{2},\\
  \frac{f(x)}{g(x)}\to\frac{\ell_{1}}{\ell_{2}} \quad(\text{se } \ell_{2}\neq0).
\end{gathered}
\end{equation}
Se $\ell_{1}$ o $\ell_{2}$ sono infiniti valgono le regole \eqref{eq:inf1}--\eqref{eq:inf7}
del \cref{teo:operazioni-infinite}, con le funzioni al posto delle successioni; restano
escluse le forme indeterminate $\infty-\infty$, $0\cdot\infty$, $\frac\infty\infty$,
$\frac00$.
\end{teorema}

\begin{schema}
\Ipotesi{$f,g:A\to\R$; \ $x_{0}$ di accumulazione per $A$; \ $f(x)\to\ell_{1}$,
$g(x)\to\ell_{2}$ per $x\to x_{0}$.}
\Tesi{le \eqref{eq:operazioni-funzioni} (per il quoziente, $\ell_{2}\neq0$); le regole
\eqref{eq:inf1}--\eqref{eq:inf7} se un limite è infinito.}
\end{schema}

\begin{strategia}
Si segue una generica successione $x_{n}\to x_{0}$: lungo di essa $f(x_{n})$ e
$g(x_{n})$ sono successioni, e per esse le regole di calcolo sono già note dal
\cref{cap:successioni}. L'unico punto che richiede attenzione è il quoziente, che è
definito solo dove $g\neq0$: con un argomento di permanenza del segno si verifica che vicino a
$x_{0}$ il denominatore non si annulla, così che $x_{0}$ resta di accumulazione per il
dominio di $f/g$.
\end{strategia}

\begin{dimostrazione}
Dimostriamo ad esempio che il limite di un prodotto è uguale al prodotto dei limiti.
Sia $x_{n}$ una qualunque successione di punti di $A\setminus\{x_{0}\}$ che tende a
$x_{0}$. Per ipotesi, e per la definizione di limite, $f(x_{n})\to\ell_{1}$ e
$g(x_{n})\to\ell_{2}$. Per la \eqref{eq:lim-prodotto}, che esprime il modo di calcolare il
limite del prodotto di due successioni, risulta $f(x_{n})\cdot g(x_{n})\to\ell_{1}\cdot\ell_{2}$;
data l'arbitrarietà della successione, ciò completa la prova. Somma, differenza e casi
infiniti si trattano allo stesso modo.

Il quoziente richiede una precisazione: $f/g$ è definita sull'insieme
$A'=\{x\in A : g(x)\neq0\}$, e bisogna verificare che $x_{0}$ sia di accumulazione per
$A'$. Supponiamo $x_{0}\in\R$. Applicando la \eqref{eq:eps-delta} a $g$ con
$\epsilon=\abs{\ell_{2}}/2>0$, troviamo $\delta>0$ tale che, per $x\in A$ con
$0<\abs{x-x_{0}}<\delta$, si ha $\abs{g(x)-\ell_{2}}<\abs{\ell_{2}}/2$ e quindi, per la
disuguaglianza triangolare inversa (\cref{cor:triangolare-inversa}),
$\abs{g(x)}>\abs{\ell_{2}}-\abs{\ell_{2}}/2=\abs{\ell_{2}}/2>0$. Dunque tutti i punti di
$A\setminus\{x_{0}\}$ con $\abs{x-x_{0}}<\delta$ stanno in $A'$, e $x_{0}$, che è di
accumulazione per $A$, lo è anche per $A'$. Lungo ogni successione $x_{n}\in A'\setminus\{x_{0}\}$
convergente a $x_{0}$ si ha poi $f(x_{n})/g(x_{n})\to\ell_{1}/\ell_{2}$ per la
\eqref{eq:lim-quoziente}. Il caso $x_{0}=\pm\infty$ è identico, con una semiretta al posto
dell'intorno di raggio $\delta$.
\end{dimostrazione}

\begin{esempio}\label{es:limite-polinomi}
Dall'\cref{es:limiti-elementari}(iii) e dal teorema precedente, applicato un numero
finito di volte, segue che per ogni polinomio $P(x)=a_{m}x^{m}+\dots+a_{1}x+a_{0}$ si ha
$\lim_{x\to x_{0}}P(x)=P(x_{0})$ per ogni $x_{0}\in\R$, e che per ogni funzione razionale
$P(x)/Q(x)$ con $Q(x_{0})\neq0$ si ha $\lim_{x\to x_{0}}P(x)/Q(x)=P(x_{0})/Q(x_{0})$. Ad
esempio $\lim_{x\to2}\frac{x^{3}-1}{x+1}=\frac73$.
\end{esempio}

Come applicazione calcoliamo il limite, per $x\to0$, del rapporto $\frac{1-\cos x}{x^{2}}$.
È una forma indeterminata $\frac00$. Moltiplicando numeratore e denominatore per
$1+\cos x$, che è positivo per $\abs x<\frac\pi2$, e usando la relazione fondamentale
$\sen^{2}x+\cos^{2}x=1$, otteniamo per $0<\abs x<\frac\pi2$
\begin{equation}\label{eq:lim-1-cos}
\begin{split}
  \frac{1-\cos x}{x^{2}}&=\frac{1-\cos^{2}x}{x^{2}(1+\cos x)}
  =\frac{\sen^{2}x}{x^{2}(1+\cos x)}\\
  &=\Bigl(\frac{\sen x}{x}\Bigr)^{2}\cdot\frac1{1+\cos x}
  \ \xrightarrow[x\to0]{}\ 1^{2}\cdot\frac1{1+1}=\frac12 ,
\end{split}
\end{equation}
dove abbiamo usato la \eqref{eq:lim-senx-x}, la \eqref{eq:lim-sen-cos-0} e il
\cref{teo:operazioni-funzioni}. Con lo stesso metodo, oppure scrivendo
$\frac{1-\cos x}{x}=x\cdot\frac{1-\cos x}{x^{2}}$, si ottiene
\begin{equation}\label{eq:lim-1-cos-x}
  \lim_{x\to0}\frac{1-\cos x}{x}=0\cdot\frac12=0 .
\end{equation}

\begin{idea}[Il secondo ordine dei piccoli angoli]
Nel \cref{cap:successioni} abbiamo giustificato l'approssimazione $\sen\alpha\approx\alpha$.
La \eqref{eq:lim-1-cos} giustifica la successiva, $\cos\alpha\approx1-\frac{\alpha^{2}}2$:
l'errore $\cos\alpha-\bigl(1-\frac{\alpha^{2}}2\bigr)$, diviso per $\alpha^{2}$, tende a
zero. Per questo in meccanica del volo il coseno di un angolo piccolo si confonde con
$1$ «al primo ordine»: in una virata corretta con angolo di sbandamento $\phi$ il fattore
di carico è $n=\frac1{\cos\phi}$, e per piccoli $\phi$ si ha
$n\approx1+\frac{\phi^{2}}2$; con $\phi=10^{\circ}=0{,}1745$~rad si ottiene $n\approx1{,}0152$,
contro il valore esatto $1{,}0154$. L'aumento di carico è dunque del secondo ordine
nell'angolo: piccolo, ma non nullo.
\end{idea}

Il calcolo numerico di limiti come \eqref{eq:lim-1-cos} nasconde però un'insidia, che il
programma seguente mette in evidenza.

\begin{python}[Una formula giusta, un risultato sbagliato]
import math

# (1 - cos x)/x^2 -> 1/2. Due formule uguali in aritmetica esatta:
ingenua = lambda x: (1 - math.cos(x)) / x**2
# 1 - cos x = 2 sen^2(x/2): nessuna sottrazione tra numeri vicini
stabile = lambda x: 2 * math.sin(x / 2)**2 / x**2

for k in range(1, 10):
    x = 10.0**(-k)
    print(f"x = 1e-{k}:   ingenua = {ingenua(x):.10f}   "
          f"stabile = {stabile(x):.10f}")
\end{python}
\begin{risultato}
x = 1e-1:   ingenua = 0.4995834722   stabile = 0.4995834722
x = 1e-2:   ingenua = 0.4999958333   stabile = 0.4999958333
x = 1e-3:   ingenua = 0.4999999583   stabile = 0.4999999583
x = 1e-4:   ingenua = 0.4999999970   stabile = 0.4999999996
x = 1e-5:   ingenua = 0.5000000414   stabile = 0.5000000000
x = 1e-6:   ingenua = 0.5000444503   stabile = 0.5000000000
x = 1e-7:   ingenua = 0.4996003611   stabile = 0.5000000000
x = 1e-8:   ingenua = 0.0000000000   stabile = 0.5000000000
x = 1e-9:   ingenua = 0.0000000000   stabile = 0.5000000000
\end{risultato}

\begin{attenzione}[La cancellazione numerica]
Le due formule sono identiche in aritmetica esatta, ma la prima peggiora a partire da
$x=10^{-4}$ e da $x=10^{-8}$ in poi restituisce $0$. Il motivo è che per $x$ piccolo
$\cos x$ è vicinissimo a $1$: in doppia precisione, con circa $16$ cifre significative,
per $x=10^{-8}$ il valore $\cos x=1-5\cdot10^{-17}+\dots$ viene arrotondato a $1$, e la
differenza $1-\cos x$ si annulla. Sottraendo due numeri quasi uguali si perdono le cifre
significative che hanno in comune (\emph{cancellazione}). La formula
$1-\cos x=2\sen^{2}\frac x2$ evita la sottrazione e dà il risultato corretto. È lo
stesso fenomeno osservato nel \cref{cap:successioni} per $\bigl(1+\frac1n\bigr)^{n}$: un
limite si dimostra con il ragionamento, e anche il calcolo numerico «di controllo» va
impostato tenendo conto della teoria.
\end{attenzione}

\subsection{Teoremi di confronto}

I teoremi di confronto del \cref{sec:confronto} si trasferiscono alle funzioni. Nel loro
enunciato è comodo supporre che le disuguaglianze tra le funzioni valgano soltanto
\emph{vicino} a $x_{0}$: diremo che una proprietà vale \emph{in un intorno di $x_{0}$}
se esiste un intorno $U$ di $x_{0}$ tale che la proprietà sia vera per ogni
$x\in(U\cap A)\setminus\{x_{0}\}$.

\begin{teorema}[Confronto]\label{teo:confronto-funzioni}\index{teorema!della permanenza del segno}\index{teorema!dei carabinieri}
Siano $f,g,h:A\to\R$ e $x_{0}\in\overline{\R}$ di accumulazione per $A$.
\begin{enumerate}[label=\textup{(\roman*)}]
  \item \emph{(Permanenza del segno)} Se $\lim_{x\to x_{0}}f(x)=\ell$ con $\ell>0$
  oppure $\ell=+\infty$, allora $f(x)>0$ in un intorno di $x_{0}$.
  \item \emph{(Confronto dei limiti)} Se $f(x)\le g(x)$ in un intorno di $x_{0}$, e
  $f(x)\to\ell_{1}$, $g(x)\to\ell_{2}$, allora $\ell_{1}\le\ell_{2}$.
  \item \emph{(Carabinieri)} Se $f(x)\le h(x)\le g(x)$ in un intorno di $x_{0}$, e
  $f(x)\to\ell$, $g(x)\to\ell$ con $\ell\in\R$, allora anche $h(x)\to\ell$.
  \item \emph{(Confronto con limiti infiniti)} Se $f(x)\le g(x)$ in un intorno di
  $x_{0}$ e $f(x)\to+\infty$, allora $g(x)\to+\infty$.
  \item \emph{(Limitata per infinitesima)} Se $f$ è limitata in un intorno di $x_{0}$ e
  $g(x)\to0$, allora $f(x)g(x)\to0$.
\end{enumerate}
\end{teorema}

\begin{schema}
\Ipotesi{(i) $f\to\ell>0$; \ (ii) $f\le g$, $f\to\ell_{1}$, $g\to\ell_{2}$; \
(iii) $f\le h\le g$, $f,g\to\ell$; \ (iv) $f\le g$, $f\to+\infty$; \ (v) $\abs f\le L$,
$g\to0$ (le disuguaglianze in un intorno di $x_{0}$).}
\Tesi{(i) $f>0$ vicino a $x_{0}$; \ (ii) $\ell_{1}\le\ell_{2}$; \ (iii) $h\to\ell$; \
(iv) $g\to+\infty$; \ (v) $fg\to0$.}
\end{schema}

\begin{strategia}
Per la (i) usiamo la definizione con le disuguaglianze, scegliendo la tolleranza
$\epsilon=\ell/2$ come per le successioni. Per le altre basta osservare che, lungo
qualsiasi successione $x_{n}\to x_{0}$, i termini finiscono per stare nell'intorno in
cui valgono le ipotesi; a quel punto i teoremi del \cref{sec:confronto} si applicano
alle successioni $f(x_{n})$, $g(x_{n})$, $h(x_{n})$, eventualmente traslate.
\end{strategia}

\begin{dimostrazione}
(i) Sia $x_{0}\in\R$ e $\ell\in\R$. Con $\epsilon=\ell/2$ la \eqref{eq:eps-delta} fornisce
$\delta>0$ tale che $\abs{f(x)-\ell}<\ell/2$, e quindi $f(x)>\ell-\ell/2=\ell/2>0$, per
ogni $x\in A$ con $0<\abs{x-x_{0}}<\delta$. Se $\ell=+\infty$ si usa la definizione con
$M=1$; se $x_{0}=\pm\infty$ si sostituisce l'intorno di raggio $\delta$ con una semiretta.

(ii)--(v) Sia $U$ l'intorno in cui valgono le disuguaglianze (o la limitatezza), e sia
$x_{n}\in A\setminus\{x_{0}\}$ una successione tendente a $x_{0}$. Esiste $\nu$ tale che
$x_{n}\in U$ per $n>\nu$ (\cref{oss:ponte-infinito}). Le successioni traslate
$a_{k}=f(x_{\nu+k})$, $b_{k}=g(x_{\nu+k})$, $c_{k}=h(x_{\nu+k})$ verificano le
disuguaglianze per \emph{ogni} $k\in\N$, e hanno gli stessi limiti di $f(x_{n})$,
$g(x_{n})$, $h(x_{n})$. Basta allora applicare, rispettivamente, il
\cref{cor:confronto-limiti}, il \cref{teo:carabinieri}, il \cref{teo:confronto-infinito}
e il \cref{teo:limitata-infinitesima}.
\end{dimostrazione}

\begin{osservazione}
Come per le successioni, nella (ii) le disuguaglianze strette non passano al limite:
$x^{2}>0$ per ogni $x\neq0$, ma $\lim_{x\to0}x^{2}=0$. Dalla (i) applicata a $f-L$ segue
inoltre che, se $f(x)\to\ell>L$, allora $f(x)>L$ in un intorno di $x_{0}$.
\end{osservazione}

\begin{esempio}\label{es:x-sen-1-x}
\begin{enumerate}[label=\textup{(\roman*)}]
  \item $\displaystyle\lim_{x\to0}x\sen\frac1x=0$. La funzione $\sen\frac1x$ non ha
  limite, ma è limitata ($\abs{\sen\frac1x}\le1$), e $x\to0$: si applica la (v). Si può
  anche usare la (iii), perché $-\abs x\le x\sen\frac1x\le\abs x$ e $\pm\abs x\to0$:
  è ciò che si vede nella parte destra della \cref{fig:sen-1-x}.
  \item $\displaystyle\lim_{x\to+\infty}\frac{\sen x}{x}=0$: il seno è limitato e
  $\frac1x\to0$ per $x\to+\infty$. È l'informazione contenuta nella stima
  \eqref{eq:senx-stima} e visibile nella \cref{fig:senx-iperboli}.
\end{enumerate}
\end{esempio}

\subsection{Limiti di funzioni composte}

Una ulteriore proprietà, molto utile nelle applicazioni, riguarda la composizione di
funzioni: è la giustificazione rigorosa del «cambiamento di variabile» nei limiti.

\begin{teorema}[Limite di una funzione composta]\label{teo:limite-composta}\index{limite!di una funzione composta}
Siano $g:X\to Y$ e $f:Y\to\R$ due funzioni, $x_{0}\in\overline{\R}$ di accumulazione per
$X$ e $y_{0}\in\overline{\R}$ di accumulazione per $Y$, tali che
\begin{equation}\label{eq:composta-ipotesi}
  \lim_{x\to x_{0}}g(x)=y_{0}, \qquad \lim_{y\to y_{0}}f(y)=\ell,
\end{equation}
ed esista un intorno $U$ di $x_{0}$ tale che risulti $g(x)\neq y_{0}$ per ogni
$x\in(U\cap X)\setminus\{x_{0}\}$. Allora è anche
\begin{equation}\label{eq:composta-tesi}
  \lim_{x\to x_{0}}f(g(x))=\ell .
\end{equation}
\end{teorema}

\begin{schema}
\Ipotesi{$g(x)\to y_{0}$ per $x\to x_{0}$; \ $f(y)\to\ell$ per $y\to y_{0}$; \
$g(x)\neq y_{0}$ per $x\neq x_{0}$ vicino a $x_{0}$.}
\Tesi{$f(g(x))\to\ell$ per $x\to x_{0}$.}
\end{schema}

\begin{strategia}
Si segue una successione $x_{n}\to x_{0}$ attraverso le due funzioni: i valori
$y_{n}=g(x_{n})$ formano una successione che tende a $y_{0}$, e quindi $f(y_{n})\to\ell$
per la definizione di limite di $f$. C'è però un punto delicato: la definizione di
limite di $f$ riguarda solo le successioni con $y_{n}\neq y_{0}$, e nulla garantisce,
in generale, che $g(x_{n})$ sia diverso da $y_{0}$. A questo serve la terza ipotesi.
\end{strategia}

\begin{dimostrazione}
Consideriamo una generica successione $x_{n}$ convergente a $x_{0}$, con $x_{n}\in X$ e
$x_{n}\neq x_{0}$ per ogni $n$. Esiste $\nu$ tale che $x_{n}\in U$ per ogni $n>\nu$, e
perciò per ogni $n>\nu$ risulta anche $g(x_{n})\neq y_{0}$. La successione
$y_{n}=g(x_{n})$ è contenuta in $Y$, converge a $y_{0}$ (per la prima delle
\eqref{eq:composta-ipotesi}) ed è tale che $y_{n}\neq y_{0}$ per $n>\nu$; allora, per la
seconda delle \eqref{eq:composta-ipotesi} applicata alla successione traslata
$y_{\nu+k}$, si ha $f(y_{\nu+k})\to\ell$, e quindi $f(y_{n})\to\ell$. Cioè
$f(g(x_{n}))\to\ell$, come si voleva dimostrare.
\end{dimostrazione}

\begin{osservazione}[L'ipotesi $g(x)\neq y_{0}$ non si può omettere]\label{oss:composta-controesempio}
Siano $g(x)=0$ per ogni $x$ e $f(y)=1$ per $y\neq0$, $f(0)=0$. Allora $g(x)\to0$ per
$x\to0$ e $f(y)\to1$ per $y\to0$ (il valore $f(0)$ non conta), ma $f(g(x))=f(0)=0$ per
ogni $x$, e quindi $f(g(x))\to0\neq1$. L'ipotesi è invece automaticamente soddisfatta
se $y_{0}=\pm\infty$ (i valori di $g$ sono numeri reali), oppure se $y_{0}$ non
appartiene al dominio $Y$ di $f$. Vedremo nel \cref{sec:continue} che essa può essere
sostituita dalla continuità di $f$ in $y_{0}$.
\end{osservazione}

\begin{esempio}\label{es:composte}
\begin{enumerate}[label=\textup{(\roman*)}]
  \item Per ogni $k\neq0$, $\displaystyle\lim_{x\to0}\frac{\sen kx}{x}=k$. Infatti
  $\frac{\sen kx}{x}=k\cdot\frac{\sen kx}{kx}$, e con $g(x)=kx$ (che tende a $0$ ed è
  diversa da $0$ per $x\neq0$) e $f(y)=\frac{\sen y}{y}$ il \cref{teo:limite-composta}
  dà $\frac{\sen kx}{kx}\to1$.
  \item Di nuovo la \eqref{eq:lim-1-cos}, per un'altra via: poiché
  $1-\cos x=2\sen^{2}\frac x2$ (dalla \eqref{eq:duplicazione} con $\frac x2$ al posto di
  $x$, e dalla relazione fondamentale \eqref{eq:relazione-fondamentale}), si ha
  $\frac{1-\cos x}{x^{2}}=\frac12\Bigl(\frac{\sen(x/2)}{x/2}\Bigr)^{2}\to\frac12$.
  \item \emph{Cambio di variabile $y=\frac1x$.} Con $g(x)=\frac1x$, che tende a $+\infty$
  per $x\to0^{+}$ e a $-\infty$ per $x\to0^{-}$, e con $f(y)=\bigl(1+\frac1y\bigr)^{y}$,
  si ha $f(g(x))=(1+x)^{1/x}$; dalla \eqref{eq:e-funzioni} e dalla \cref{prop:dx-sx}
  segue
  \begin{equation}\label{eq:lim-1-piu-x}
    \lim_{x\to0}(1+x)^{1/x}=e .
  \end{equation}
  \item Ancora con $y=\frac1x$: $\displaystyle\lim_{x\to+\infty}x\sen\frac1x
  =\lim_{y\to0^{+}}\frac{\sen y}{y}=1$. Si confronti con l'\cref{es:x-sen-1-x}(i): la
  stessa espressione ha limite $0$ per $x\to0$ e limite $1$ per $x\to+\infty$.
\end{enumerate}
\end{esempio}

%% ----------------------------------------------------------------------------
%%  §4.5  Funzioni continue
%% ----------------------------------------------------------------------------
\section{Funzioni continue}
\label{sec:continue}

Come nel paragrafo precedente, consideriamo funzioni $f$ definite in un dominio $A$
costituito da un intervallo, o dall'unione finita di intervalli, e punti $x_{0}$ di $A$.
Abbiamo introdotto i limiti di funzioni soprattutto per descrivere il comportamento di
una funzione nelle vicinanze di punti «singolari», come $x_{0}=0$ per $\frac{\sen x}{x}$.
Naturalmente possiamo calcolare il limite per $x\to x_{0}$ anche se la funzione non
presenta alcuna singolarità in $x_{0}$. Ad esempio abbiamo già calcolato in
\eqref{eq:lim-sen-cos-0} i limiti
\begin{equation}\label{eq:sen-cos-in-0}
  \lim_{x\to0}\sen x=0=\sen0, \qquad \lim_{x\to0}\cos x=1=\cos0:
\end{equation}
il valore limite, per $x\to0$, è uguale al valore che si ottiene calcolando la funzione
per $x=0$. Si dice che le funzioni $\sen x$ e $\cos x$ sono \emph{continue} per $x=0$ (e
in realtà, come vedremo, sono continue in ogni $x_{0}\in\R$), in accordo con la seguente
definizione.

\begin{definizione}[Funzione continua]\label{def:continua}\index{funzione!continua}\index{continuità}
Siano $f:A\to\R$ e $x_{0}\in A$ un punto di accumulazione per $A$. La funzione $f$ si
dice \emph{continua nel punto $x_{0}$} se
\begin{equation}\label{eq:def-continua}
  \lim_{x\to x_{0}}f(x)=f(x_{0}) .
\end{equation}
La funzione si dice \emph{continua in un insieme} $B\subseteq A$ se è continua in ogni
punto $x_{0}\in B$; si dice semplicemente \emph{continua} se è continua in tutto il suo
dominio.
\end{definizione}

Nella definizione sono contenute tre richieste: che $f$ sia definita in $x_{0}$, che
esista il limite di $f$ per $x\to x_{0}$, e che il limite coincida con il valore. Se
$A=[a,b]$ e $x_{0}=a$, nella \eqref{eq:def-continua} intervengono soltanto punti $x>a$
(non ci sono punti del dominio a sinistra di $a$), e dunque il limite è di fatto il solo
limite destro $x\to a^{+}$; analogamente, se $x_{0}=b$, si considera il limite sinistro
$x\to b^{-}$. Una funzione è quindi continua nell'intervallo $[a,b]$ se è continua in ogni
punto interno, e se agli estremi il limite destro (in $a$) e il limite sinistro (in $b$)
coincidono con i valori $f(a)$ e $f(b)$.

\begin{osservazione}
Se $x_{0}$ fosse un punto isolato di $A$ (cioè un punto di $A$ che non è di
accumulazione), la \eqref{eq:def-continua} non avrebbe senso. Per convenzione si dice
che ogni funzione è continua nei punti isolati del suo dominio; con i domini che
consideriamo, intervalli non degeneri o loro unioni finite, il caso non si presenta.
\end{osservazione}

La continuità ammette due riformulazioni, che useremo continuamente: una con le
successioni, senza la restrizione $x_{n}\neq x_{0}$, e una con $\epsilon$ e $\delta$,
senza la restrizione $x\neq x_{0}$.

\begin{proposizione}[Caratterizzazioni della continuità]\label{prop:caratt-continuita}
Siano $f:A\to\R$ e $x_{0}\in A$ di accumulazione per $A$. Sono equivalenti:
\begin{enumerate}[label=\textup{(\alph*)}]
  \item $f$ è continua in $x_{0}$;
  \item per ogni successione $x_{n}$ di punti di $A$ convergente a $x_{0}$ (anche con
  $x_{n}=x_{0}$ per qualche $n$) risulta $f(x_{n})\to f(x_{0})$;
  \item per ogni $\epsilon>0$ esiste $\delta>0$ tale che $\abs{f(x)-f(x_{0})}<\epsilon$ per
  ogni $x\in A$ con $\abs{x-x_{0}}<\delta$.
\end{enumerate}
\end{proposizione}

\begin{schema}
\Ipotesi{$f:A\to\R$; \ $x_{0}\in A$ di accumulazione per $A$.}
\Tesi{(a) $\sse$ (b) $\sse$ (c).}
\end{schema}

\begin{strategia}
Si dimostra il ciclo di implicazioni (a)$\implica$(c)$\implica$(b)$\implica$(a), sfruttando
il \cref{teo:eps-delta}. L'osservazione chiave è che, quando il limite coincide con il
valore $f(x_{0})$, la condizione $x\neq x_{0}$ (o $x_{n}\neq x_{0}$) diventa superflua:
nel punto $x_{0}$ la distanza $\abs{f(x)-f(x_{0})}$ è nulla, e quindi minore di ogni
$\epsilon$.
\end{strategia}

\begin{dimostrazione}
(a)$\implica$(c): per il \cref{teo:eps-delta} la disuguaglianza
$\abs{f(x)-f(x_{0})}<\epsilon$ vale per $x\in A$ con $0<\abs{x-x_{0}}<\delta$; per
$x=x_{0}$ vale banalmente, perché il primo membro è zero.

(c)$\implica$(b): è la prima parte della dimostrazione del \cref{teo:eps-delta}, in cui
non serve più l'ipotesi $x_{n}\neq x_{0}$: se $x_{n}\to x_{0}$, definitivamente
$\abs{x_{n}-x_{0}}<\delta$, e quindi $\abs{f(x_{n})-f(x_{0})}<\epsilon$.

(b)$\implica$(a): le successioni con $x_{n}\neq x_{0}$ sono un caso particolare di quelle
considerate in (b); dunque $f(x)\to f(x_{0})$ per la \cref{def:limite-funzione}.
\end{dimostrazione}

La (b) si legge così: \emph{per una funzione continua il limite «passa dentro» la
funzione}, $\lim_{n}f(x_{n})=f\bigl(\lim_{n}x_{n}\bigr)$. La (c), scritta con
$x=x_{0}+h$, dice che $f(x_{0}+h)-f(x_{0})\to0$ per $h\to0$: a un incremento piccolo
della variabile corrisponde un incremento piccolo della funzione, che è esattamente la
definizione di Cauchy riportata nell'epigrafe del capitolo. In particolare
\begin{equation}\label{eq:continua-h}
  f \text{ è continua in } x_{0} \quad\sse\quad \lim_{h\to0}f(x_{0}+h)=f(x_{0}).
\end{equation}

\subsection{Operazioni con le funzioni continue}

\begin{teorema}[Operazioni con funzioni continue]\label{teo:operazioni-continue}
Siano $f,g:A\to\R$ e $x_{0}\in A$ di accumulazione per $A$. Se $f$ e $g$ sono continue
in $x_{0}$, sono continue in $x_{0}$ anche $f+g$, $f-g$, $f\cdot g$ e, se
$g(x_{0})\neq0$, $f/g$.
\end{teorema}

\begin{schema}
\Ipotesi{$f$, $g$ continue in $x_{0}$ \ (per il quoziente, $g(x_{0})\neq0$).}
\Tesi{$f\pm g$, $fg$, $f/g$ continue in $x_{0}$.}
\end{schema}

\begin{strategia}
Basta tradurre la continuità in termini di limiti e applicare il
\cref{teo:operazioni-funzioni} con $\ell_{1}=f(x_{0})$ e $\ell_{2}=g(x_{0})$: il limite
dell'operazione è l'operazione sui valori, cioè il valore della nuova funzione in
$x_{0}$.
\end{strategia}

\begin{dimostrazione}
Dato che il limite di somma, differenza, prodotto, quoziente è uguale rispettivamente
alla somma, differenza, prodotto, quoziente dei limiti (\cref{teo:operazioni-funzioni}),
si ha ad esempio $\lim_{x\to x_{0}}f(x)g(x)=f(x_{0})g(x_{0})$. Per il quoziente
occorre al solito fare attenzione ai punti dove il denominatore si annulla: se
$g(x_{0})\neq0$, la dimostrazione del \cref{teo:operazioni-funzioni} mostra che $g$ non
si annulla vicino a $x_{0}$, e $f(x)/g(x)\to f(x_{0})/g(x_{0})$.
\end{dimostrazione}

\begin{teorema}[Continuità della funzione composta]\label{teo:composta-continua}\index{funzione!composta!continuità}
Siano $g:X\to Y$ continua in $x_{0}\in X$ e $f:Y\to\R$ continua in $y_{0}=g(x_{0})$.
Allora la funzione composta $f\circ g$ è continua in $x_{0}$.
\end{teorema}

\begin{schema}
\Ipotesi{$g$ continua in $x_{0}$; \ $f$ continua in $g(x_{0})$.}
\Tesi{$f(g(x))\to f(g(x_{0}))$ per $x\to x_{0}$.}
\end{schema}

\begin{strategia}
Con la caratterizzazione per successioni, senza la restrizione $x_{n}\neq x_{0}$, la
dimostrazione è una catena: $x_{n}\to x_{0}$ implica $g(x_{n})\to g(x_{0})$, che a sua
volta implica $f(g(x_{n}))\to f(g(x_{0}))$. Non occorre preoccuparsi dei termini con
$g(x_{n})=g(x_{0})$, che invece creavano difficoltà nel \cref{teo:limite-composta}.
\end{strategia}

\begin{dimostrazione}
Usiamo la caratterizzazione (b) della \cref{prop:caratt-continuita}. Se $x_{n}\in X$ e
$x_{n}\to x_{0}$, per la continuità di $g$ si ha $y_{n}=g(x_{n})\to g(x_{0})=y_{0}$, con
$y_{n}\in Y$; per la continuità di $f$ in $y_{0}$ si ha allora $f(y_{n})\to f(y_{0})$,
cioè $f(g(x_{n}))\to f(g(x_{0}))$.
\end{dimostrazione}

Si noti che qui non serve l'ipotesi $g(x)\neq y_{0}$ del \cref{teo:limite-composta}:
proprio perché la (b) ammette successioni con $y_{n}=y_{0}$. Lo stesso argomento prova
la seguente versione, che useremo spesso nel calcolo dei limiti.

\begin{corollario}\label{cor:limite-composta-continua}
Siano $g:X\to Y$, $f:Y\to\R$ e $x_{0}\in\overline{\R}$ di accumulazione per $X$. Se
$\lim_{x\to x_{0}}g(x)=y_{0}$, con $y_{0}\in Y$, e $f$ è continua in $y_{0}$, allora
\begin{equation}\label{eq:limite-composta-continua}
  \lim_{x\to x_{0}}f(g(x))=f\Bigl(\lim_{x\to x_{0}}g(x)\Bigr)=f(y_{0}) .
\end{equation}
\end{corollario}

\begin{schema}
\Ipotesi{$g(x)\to y_{0}$ per $x\to x_{0}$, con $y_{0}\in Y$; \ $f$ continua in $y_{0}$.}
\Tesi{$f(g(x))\to f(y_{0})$ per $x\to x_{0}$.}
\end{schema}

\begin{strategia}
È la stessa catena del teorema precedente, con una sola differenza: ora $x_{0}$ può
non appartenere al dominio di $g$, e si considerano soltanto successioni con
$x_{n}\neq x_{0}$. Dove arrivano i valori $g(x_{n})$ non importa, perché la
caratterizzazione (b) della continuità di $f$ ammette anche $g(x_{n})=y_{0}$.
\end{strategia}

\begin{dimostrazione}
Se $x_{n}\in X\setminus\{x_{0}\}$ e $x_{n}\to x_{0}$, allora $g(x_{n})\to y_{0}$, e per
la \cref{prop:caratt-continuita}(b) applicata a $f$ (che non richiede $g(x_{n})\neq
y_{0}$) si ha $f(g(x_{n}))\to f(y_{0})$.
\end{dimostrazione}

\subsection{Continuità delle funzioni elementari}

L'importanza delle funzioni continue sta anche nel fatto che \emph{le funzioni
elementari sono continue nel loro insieme di definizione}. Esaminiamole una per volta.

\paragraph{Polinomi e funzioni razionali.} Per l'\cref{es:limite-polinomi}, ogni polinomio
è continuo su tutto $\R$, e ogni funzione razionale $P(x)/Q(x)$ è continua in ogni punto
in cui $Q(x)\neq0$.

\paragraph{Valore assoluto.} La continuità della funzione $f(x)=\abs x$ su $\R$ segue
dalla disuguaglianza triangolare inversa (\cref{cor:triangolare-inversa})
\begin{equation}\label{eq:modulo-lipschitz}
  \bigl\lvert\abs x-\abs{x_{0}}\bigr\rvert\le\abs{x-x_{0}} \qquad\perogni x,x_{0}\in\R:
\end{equation}
basta infatti porre $\delta=\epsilon$ nella \cref{prop:caratt-continuita}(c): se
$\abs{x-x_{0}}<\epsilon$, allora anche $\abs{f(x)-f(x_{0})}=\bigl\lvert\abs x-\abs{x_{0}}\bigr\rvert<\epsilon$.

\begin{proposizione}[Continuità di seno, coseno, tangente]\label{prop:sen-cos-continue}\index{funzione!seno!continuità}
Le funzioni $\sen x$ e $\cos x$ sono continue su $\R$; la funzione $\tg x$ è continua
in ogni punto $x\neq\frac\pi2+k\pi$, $k\in\Z$.
\end{proposizione}

\begin{schema}
\Ipotesi{$\sen h\to0$, $\cos h\to1$ per $h\to0$ \eqref{eq:sen-cos-in-0}; \ formule di
addizione.}
\Tesi{$\sen$, $\cos$ continue in ogni $x_{0}\in\R$; \ $\tg$ continua dove $\cos x\neq0$.}
\end{schema}

\begin{strategia}
Conosciamo la continuità di seno e coseno nel solo punto $0$ \eqref{eq:sen-cos-in-0}. Le
formule di addizione permettono di «trasportare» questa informazione in un punto
qualsiasi: $\sen(x_{0}+h)$ si esprime per mezzo di $\sen h$ e $\cos h$, che per $h\to0$
tendono a $0$ e a $1$.
\end{strategia}

\begin{dimostrazione}
Sia $x_{0}\in\R$ e sia $x_{n}\to x_{0}$. Posto $h_{n}=x_{n}-x_{0}$, si ha $h_{n}\to0$, e
quindi $\sen h_{n}\to0$, $\cos h_{n}\to1$ per il \cref{teo:lim-sen-cos}. Per la formula
di addizione del seno \eqref{eq:addizione-sen}
\begin{equation}\label{eq:sen-continuo}
\begin{split}
  \sen x_{n}=\sen(x_{0}+h_{n})&=\sen x_{0}\cos h_{n}+\cos x_{0}\sen h_{n}\\
  &\to\ \sen x_{0}\cdot1+\cos x_{0}\cdot0=\sen x_{0},
\end{split}
\end{equation}
per le operazioni con i limiti di successioni. Analogamente, per la
\eqref{eq:addizione-cos}, $\cos x_{n}=\cos x_{0}\cos h_{n}-\sen x_{0}\sen h_{n}\to\cos x_{0}$.
Per la \cref{prop:caratt-continuita}(b) seno e coseno sono continui in $x_{0}$. Infine la
continuità della funzione $\tg x=\frac{\sen x}{\cos x}$ nei punti in cui $\cos x\neq0$,
cioè per $x\neq\frac\pi2+k\pi$, discende dal \cref{teo:operazioni-continue}.
\end{dimostrazione}

\paragraph{Esponenziali, logaritmi, potenze.} Anche le funzioni esponenziali $a^{x}$
($a>0$), i logaritmi $\log_{a}x$ ($a>0$, $a\neq1$) e le potenze $x^{b}$ (per $x>0$, con
$b\in\R$) sono continue nel loro insieme di definizione. La dimostrazione richiede il
teorema sul limite delle funzioni monotòne, ed è svolta nel \cref{sec:monotone-inverse}
(\cref{teo:continuita-esponenziale}); essa non usa alcun risultato di questo paragrafo
che dipenda dalla continuità di tali funzioni, per cui possiamo servircene fin d'ora
negli esempi. Nello stesso paragrafo ricaveremo, come applicazione, alcuni limiti
notevoli che coinvolgono $e^{x}$ e $\log x$.

\begin{table}[htb]
\centering
\small
\renewcommand{\arraystretch}{1.3}
\begin{tabular}{l l l}
\toprule
funzione & continua in & dimostrazione\\
\midrule
polinomi & $\R$ & \cref{es:limite-polinomi}\\
$P(x)/Q(x)$ & $Q(x)\neq0$ & \\
$\abs x$ & $\R$ & \eqref{eq:modulo-lipschitz}\\
$\sen x$, $\cos x$ & $\R$ & \cref{prop:sen-cos-continue}\\
$\tg x$ & $x\neq\frac\pi2+k\pi$ & \\
$a^{x}$ & $\R$ & \cref{teo:continuita-esponenziale}\\
$\log_{a}x$ & $(0,+\infty)$ & \\
$x^{b}$ & $(0,+\infty)$ & \\
\bottomrule
\end{tabular}
\caption{La continuità delle funzioni elementari. Somme, prodotti, quozienti (dove il
denominatore non si annulla) e composizioni di funzioni continue sono continue
(\cref{teo:operazioni-continue,teo:composta-continua}).}
\label{tab:elementari-continue}
\end{table}

Grazie alla \cref{tab:elementari-continue} e ai \cref{teo:operazioni-continue,%
teo:composta-continua}, la continuità di funzioni come
$x\mapsto\sen(x^{2}+1)$, $x\mapsto e^{\cos x}$ o $x\mapsto\abs{\log x}$ è immediata, e i
corrispondenti limiti nei punti del dominio si calcolano per sostituzione: ad esempio
$\lim_{x\to\pi}e^{\cos x}=e^{\cos\pi}=e^{-1}$.

\begin{idea}[L'atmosfera standard: continua, ma con uno spigolo]
Nel modello dell'atmosfera standard internazionale (ISA) la temperatura, in funzione
della quota $h$ espressa in metri, è
\[
  T(h)=
  \begin{cases}
    288{,}15-0{,}0065\,h & \text{per } 0\le h\le11\,000 \ \ \text{(troposfera)},\\
    216{,}65 & \text{per } 11\,000<h\le20\,000 \ \ \text{(stratosfera)},
  \end{cases}
\]
in kelvin. Ciascun tratto è un polinomio, quindi continuo; nel punto di raccordo, la
tropopausa, il limite sinistro vale $288{,}15-0{,}0065\cdot11\,000=216{,}65$, uguale al
limite destro e al valore: $T$ è continua anche lì. Il modello non poteva essere
altrimenti: una temperatura discontinua in quota non avrebbe senso fisico, e la
pressione e la densità, che si ricavano da $T(h)$, risultano anch'esse continue. Il
grafico di $T$ ha però uno \emph{spigolo} a $11\,000$~m: la pendenza passa bruscamente da
$-6{,}5$~K/km a zero. La continuità non esclude cambi di pendenza improvvisi; li
misureremo con le derivate.
\end{idea}

\begin{notastorica}[Che cosa vuol dire «continua»?]
Per Eulero, nell'\emph{Introductio in analysin infinitorum} (1748), una funzione era
«continua» se era data da un'unica espressione analitica, e «discontinua» (o mista) se
era definita da espressioni diverse in parti diverse del dominio: secondo questo uso la
funzione $T(h)$ dell'atmosfera standard sarebbe stata discontinua, e $\frac1x$
continua! Discontinua sarebbe stata anche $\abs x$, che vale $x$ per $x\ge0$ e $-x$ per
$x<0$, ed è per noi continua ovunque. A mettere in crisi questo linguaggio fu il problema
della corda vibrante. Nel 1747 d'Alembert ne scrisse la soluzione generale come somma di
due onde che viaggiano in versi opposti; Eulero obiettò che la forma iniziale della
corda può essere qualsiasi, ad esempio triangolare, con uno spigolo, come quella di una
corda pizzicata, e che bisognava ammettere anche curve «discontinue», perfino tracciate a
mano libera. Daniel Bernoulli (1753) sostenne invece che ogni vibrazione è una
sovrapposizione di oscillazioni sinusoidali, cioè la somma di una serie trigonometrica.
Eulero e d'Alembert respinsero questa tesi: una serie di seni era, ai loro occhi, un'unica
espressione, e non poteva rappresentare una curva «discontinua». La risposta venne da
Joseph Fourier, che nella \emph{Théorie analytique de la chaleur} (1822) rappresentò con
serie trigonometriche funzioni definite a tratti, anche con salti: la stessa «formula»
poteva descrivere oggetti molto diversi, e il concetto di funzione andava separato da
quello di formula. La definizione moderna di continuità, che guarda al comportamento della
funzione vicino a ogni punto e non alla forma della formula, fu data indipendentemente da
Bernard Bolzano (1817) e da Augustin-Louis Cauchy, nel \emph{Cours d'analyse} (1821).
\end{notastorica}

%% ----------------------------------------------------------------------------
%%  §4.6  Discontinuità e prolungamento per continuità
%% ----------------------------------------------------------------------------
\section{Discontinuità}
\label{sec:discontinuita}

La funzione
\begin{equation}\label{eq:segno}
  f(x)=\frac{\abs x}{x}=
  \begin{cases}
    1 & \text{se } x>0,\\
    -1 & \text{se } x<0
  \end{cases}
\end{equation}
è continua per $x\neq0$ (in un intorno di ogni $x_{0}\neq0$ è costante, e il limite è
una proprietà locale, \cref{lem:localita}), ma non è continua se $x=0$. Il grafico di
questa funzione presenta per $x=0$ un salto, appunto una \emph{discontinuità}
(\cref{fig:modulo-su-x}). In particolare $f$ non è continua nel punto $x_{0}=0$ perché
non è definita in tal punto, cioè perché non esiste il valore $f(x_{0})=f(0)$.

Estendiamo allora $f$ anche a $x_{0}=0$ con un valore $\ell\in\R$, considerando cioè la
nuova funzione
\begin{equation}\label{eq:segno-esteso}
  \tilde f(x)=
  \begin{cases}
    \dfrac{\abs x}{x} & \text{se } x\neq0,\\[1ex]
    \ell & \text{se } x=0 .
  \end{cases}
\end{equation}
La funzione $\tilde f$ è definita anche nel punto $x_{0}=0$, ma non è continua in tale
punto, perché non esiste il limite per $x\to0$ di $\tilde f(x)$: precisamente, il limite
destro è diverso dal limite sinistro \eqref{eq:modulo-su-x}. Tutto ciò accade
\emph{qualunque sia} il valore $\ell$ scelto in \eqref{eq:segno-esteso}: $\tilde f$ è
un'estensione non continua (o \emph{prolungamento non continuo}) della funzione $f$, e
la discontinuità di $f$ nel punto $x_{0}=0$ si dice \emph{non eliminabile}.

Viceversa, la funzione studiata nel \cref{sec:premessa-funzioni},
$f(x)=\frac{\sen x}{x}$, non è continua in $x_{0}=0$ (perché non è definita), ma è
possibile \emph{prolungarla per continuità} mediante la funzione
\begin{equation}\label{eq:sinc}
  \tilde f(x)=
  \begin{cases}
    \dfrac{\sen x}{x} & \text{se } x\neq0,\\[1ex]
    1 & \text{se } x=0 .
  \end{cases}
\end{equation}
A causa del limite notevole \eqref{eq:lim-senx-x}, $\tilde f$ è continua anche nel punto
$x_{0}=0$. Il grafico di $\tilde f$ si ottiene «completando» il disegno della
\cref{fig:senx-x} con l'ulteriore punto di coordinate $(0,1)$, che riempie il tondino
vuoto.

\subsection{Classificazione delle discontinuità}

Sia $f$ una funzione definita in $A$ e $x_{0}$ un punto di $A$, di accumulazione per $A$.
Se $f$ non è continua in $x_{0}$ si dice che $x_{0}$ è un punto di \emph{discontinuità}
per $f$. Le discontinuità si classificano nel modo seguente
(\cref{fig:discontinuita}).

\begin{definizione}[Tipi di discontinuità]\label{def:discontinuita}\index{discontinuità}
\begin{enumerate}[label=\textup{(\alph*)}]
  \item La funzione $f$ presenta in $x_{0}$ una \emph{discontinuità eliminabile} se
  esiste finito il limite di $f(x)$ per $x\to x_{0}$ e risulta
  \begin{equation}\label{eq:disc-eliminabile}
    \lim_{x\to x_{0}}f(x)\neq f(x_{0}) .
  \end{equation}
  \item La funzione $f$ presenta in $x_{0}$ una \emph{discontinuità di prima specie} (o
  \emph{di salto}) se esistono finiti i limiti destro e sinistro di $f(x)$ in $x_{0}$ e
  si ha
  \begin{equation}\label{eq:disc-prima-specie}
    \lim_{x\to x_{0}^{-}}f(x)\neq\lim_{x\to x_{0}^{+}}f(x);
  \end{equation}
  la differenza tra il limite destro e il limite sinistro si chiama \emph{salto} di $f$
  in $x_{0}$.
  \item La funzione $f$ presenta in $x_{0}$ una \emph{discontinuità di seconda specie}
  se almeno uno dei due limiti
  \begin{equation}\label{eq:disc-seconda-specie}
    \lim_{x\to x_{0}^{-}}f(x), \qquad \lim_{x\to x_{0}^{+}}f(x)
  \end{equation}
  non esiste oppure è infinito.
\end{enumerate}
\end{definizione}

La classificazione è completa: se $f$ è discontinua in $x_{0}$ (interno al dominio) e
non siamo nel caso (c), i due limiti unilaterali esistono finiti; se sono diversi siamo
nel caso (b), se sono uguali il limite esiste (\cref{prop:dx-sx}) ed è diverso da
$f(x_{0})$, e siamo nel caso (a). Per abuso di linguaggio, le stesse parole si usano
anche per un punto $x_{0}$ di accumulazione che \emph{non} appartiene al dominio, come
$x_{0}=0$ per $\frac{\sen x}{x}$ (singolarità eliminabile) o per $\frac{\abs x}{x}$
(salto).

Nel caso (a), posto $\ell=\lim_{x\to x_{0}}f(x)$, la funzione
\begin{equation}\label{eq:elimina}
  \tilde f(x)=
  \begin{cases}
    f(x) & \text{se } x\in A\setminus\{x_{0}\},\\
    \ell & \text{se } x=x_{0}
  \end{cases}
\end{equation}
risulta continua nel punto $x_{0}$: la discontinuità si «elimina» correggendo un solo
valore. Nei casi (b) e (c) nessuna modifica del valore in $x_{0}$ rende $f$ continua.

\begin{figure}[htb]
\centering
\includegraphics[width=\textwidth]{discontinuita}
\caption{I tre tipi di discontinuità. (a) La funzione vale $\frac{\sen x}{x}$ per
$x\neq0$ e $\frac12$ in $0$: il limite esiste, ma è diverso dal valore. (b) I limiti
destro e sinistro esistono finiti ma sono diversi. (c) La funzione $e^{1/x}$ ha limite
sinistro $0$ e limite destro $+\infty$ (asse $y$ in scala logaritmica per $y>1$)
(\texttt{python/fig\_discontinuita.py}).}
\label{fig:discontinuita}
\end{figure}

\begin{esempio}\label{es:parte-intera-continuita}
La funzione $f(x)=\lfloor x\rfloor$, \emph{parte intera} di $x$
(\cref{prop:parte-intera}, rappresentata nella \cref{fig:parte-intera}), presenta
discontinuità di prima specie in corrispondenza a ogni $k\in\Z$, mentre è continua per
ogni $x\in\R\setminus\Z$. Infatti, se $x_{0}\notin\Z$, esiste $k\in\Z$ con
$k<x_{0}<k+1$, e nell'intorno $(k,k+1)$ di $x_{0}$ la funzione è costante, uguale a $k$.
Se invece $x_{0}=k\in\Z$, si ha $\lfloor x\rfloor=k-1$ per $k-1\le x<k$ e
$\lfloor x\rfloor=k$ per $k\le x<k+1$; quindi
\[
  \lim_{x\to k^{-}}\lfloor x\rfloor=k-1, \qquad \lim_{x\to k^{+}}\lfloor x\rfloor=k=\lfloor k\rfloor:
\]
il salto vale $1$, e la funzione è continua da destra ma non da sinistra.
\end{esempio}

\begin{esempio}\label{es:seconda-specie}
\begin{enumerate}[label=\textup{(\roman*)}]
  \item La funzione $f(x)=\sen\frac1x$ per $x\neq0$, prolungata con un qualsiasi valore
  in $0$, ha in $0$ una discontinuità di seconda specie: il limite destro non esiste
  (lo si vede con le successioni dell'\cref{es:sen-1-x}, che sono positive).
  \item La funzione $f(x)=\frac1x$ per $x\neq0$, $f(0)=0$, ha in $0$ una discontinuità di
  seconda specie: i limiti unilaterali sono infiniti (\cref{es:modulo-su-x}).
  \item La funzione $e^{1/x}$ ha in $0$ una singolarità di seconda specie: per $x\to0^{+}$
  si ha $\frac1x\to+\infty$ ed $e^{y}\to+\infty$ per $y\to+\infty$, quindi, per il
  \cref{teo:limite-composta} (con $y_{0}=+\infty$, caso in cui l'ipotesi $g(x)\neq y_{0}$
  è automatica), $e^{1/x}\to+\infty$; per $x\to0^{-}$, invece, $\frac1x\to-\infty$ ed
  $e^{1/x}\to0$ (\cref{fig:discontinuita}(c)).
\end{enumerate}
\end{esempio}

\begin{osservazione}[Una funzione discontinua ovunque]\label{oss:dirichlet}\index{funzione!di Dirichlet}
La \emph{funzione di Dirichlet}
\[
  D(x)=
  \begin{cases}
    1 & \text{se } x\in\Q,\\
    0 & \text{se } x\in\R\setminus\Q
  \end{cases}
\]
è discontinua in ogni punto, con discontinuità di seconda specie. Fissato $x_{0}\in\R$,
per la densità dei razionali e degli irrazionali (\cref{teo:densita-Q},
\cref{cor:densita-irrazionali}) in ogni intervallo $(x_{0},x_{0}+\frac1n)$ esistono un
razionale $q_{n}$ e un irrazionale $r_{n}$; le successioni $q_{n}$ ed $r_{n}$ tendono a
$x_{0}$ da destra, ma $D(q_{n})=1\to1$ e $D(r_{n})=0\to0$, quindi il limite destro non
esiste (\cref{cor:non-esistenza}). Peter Gustav Lejeune Dirichlet la propose nel 1829,
nel suo lavoro sulle serie di Fourier, come esempio di funzione per cui l'integrale, e
con esso la sua analisi, perde significato; nel 1837 formulò esplicitamente la nozione di
funzione come corrispondenza arbitraria, priva di qualunque regolarità. Questi esempi
contribuirono a imporre la definizione moderna di funzione.
\end{osservazione}

\subsection{Il prolungamento per continuità}

\begin{definizione}[Prolungamento per continuità]\label{def:prolungamento}\index{prolungamento per continuità}
Sia $A$ un intervallo (o un'unione finita di intervalli), $x_{0}\in A$, e sia $f$ una
funzione definita in $A\setminus\{x_{0}\}$. Se esiste finito il limite
\begin{equation}\label{eq:prolungamento-limite}
  \lim_{x\to x_{0}}f(x)=\ell,
\end{equation}
la funzione $\tilde f$ definita in $A$ da
\begin{equation}\label{eq:prolungamento}
  \tilde f(x)=
  \begin{cases}
    f(x) & \text{se } x\in A\setminus\{x_{0}\},\\
    \ell & \text{se } x=x_{0}
  \end{cases}
\end{equation}
è detta \emph{prolungamento per continuità di $f$ in $x_{0}$}; per costruzione $\tilde f$
risulta continua in $x_{0}$. Se poi $f$ è continua in $A\setminus\{x_{0}\}$, allora
$\tilde f$, continua su tutto l'insieme $A$, è detta \emph{prolungamento per continuità
di $f$ su $A$}.
\end{definizione}

Il prolungamento per continuità, quando esiste, è unico, perché il limite è unico.
La \eqref{eq:sinc} è il prolungamento per continuità di $\frac{\sen x}{x}$ su $\R$; la
funzione che si ottiene, detta \emph{seno cardinale} e indicata con $\operatorname{sinc}x$,
è di uso quotidiano nell'elaborazione dei segnali: è, a meno di costanti, la trasformata
di Fourier di un impulso rettangolare, e descrive ad esempio il diagramma di
diffrazione di un'apertura o di un'antenna. Nei programmi il prolungamento va scritto
esplicitamente.

\begin{python}[Il prolungamento per continuità nei programmi]
import math
import numpy as np

try:
    print(math.sin(0.0) / 0.0)
except ZeroDivisionError as err:
    print("sen(0)/0 ->", type(err).__name__, ":", err)

def sinc(x):
    """Prolungamento per continuità di sen(x)/x: vale 1 in x = 0."""
    return 1.0 if x == 0 else math.sin(x) / x

print(sinc(0.0), sinc(1e-9), sinc(math.pi / 2))
# NumPy usa la versione normalizzata sen(pi x)/(pi x), estesa in 0
print(np.sinc(0.0), np.sinc(0.5), np.sinc(1.0))
\end{python}
\begin{risultato}
sen(0)/0 -> ZeroDivisionError : float division by zero
1.0 1.0 0.6366197723675814
1.0 0.6366197723675814 3.8981718325193755e-17
\end{risultato}

L'ultimo valore dovrebbe essere $\sen\pi/\pi=0$: è il solito residuo di arrotondamento
dovuto alla rappresentazione approssimata di $\pi$.

\begin{idea}[L'onda d'urto: una discontinuità di prima specie]
In un flusso supersonico può formarsi un'onda d'urto normale: attraversandola, la
pressione, la densità e la temperatura aumentano bruscamente e la velocità diminuisce.
Nel modello di fluido non viscoso (equazioni di Eulero) l'urto è una vera discontinuità
di prima specie: lungo l'asse della corrente la pressione vale $p_{1}$ a monte e $p_{2}$
a valle, con
\[
  \frac{p_{2}}{p_{1}}=1+\frac{2\gamma}{\gamma+1}\bigl(M_{1}^{2}-1\bigr),
\]
dove $M_{1}$ è il numero di Mach a monte e $\gamma=1{,}4$ per l'aria; con $M_{1}=2$ il
salto porta la pressione a $4{,}5$ volte il valore di monte. Nella realtà la viscosità
«smussa» l'urto in uno strato continuo spesso pochi cammini liberi medi delle molecole,
frazioni di micrometro al livello del mare: alla scala di un'ala, una discontinuità a
tutti gli effetti. Nei codici CFD con tecniche di \emph{shock capturing} lo stesso salto
viene invece distribuito su poche celle di calcolo: la soluzione numerica è continua e
approssima, raffinando la griglia, una soluzione discontinua.
\end{idea}

%% ----------------------------------------------------------------------------
%%  §4.7  Alcuni teoremi sulle funzioni continue
%% ----------------------------------------------------------------------------
\section{Alcuni teoremi sulle funzioni continue}
\label{sec:teoremi-continue}

Il teorema seguente è analogo al teorema della permanenza del segno per le successioni
(\cref{teo:permanenza-segno}) e alla sua versione per i limiti di funzioni
(\cref{teo:confronto-funzioni}(i)); la differenza è che ora la conclusione vale anche
\emph{nel} punto $x_{0}$.

\begin{teorema}[Permanenza del segno]\label{teo:permanenza-segno-continue}\index{teorema!della permanenza del segno}
Sia $f$ una funzione definita in un intervallo $A$ e continua in $x_{0}\in A$. Se
$f(x_{0})>0$, esiste un numero $\delta>0$ con la proprietà che $f(x)>0$ per ogni
$x\in A\cap(x_{0}-\delta,x_{0}+\delta)$.
\end{teorema}

\begin{schema}
\Ipotesi{$f$ continua in $x_{0}$; \ $f(x_{0})>0$.}
\Tesi{$\esiste\delta>0$: \ $f(x)>0$ \ $\perogni x\in A$, $\abs{x-x_{0}}<\delta$.}
\end{schema}

\begin{strategia}
Si sceglie una tolleranza che non permetta a $f(x)$ di raggiungere lo zero: con
$\epsilon=f(x_{0})/2$, la caratterizzazione (c) della continuità tiene $f(x)$ sopra
$f(x_{0})/2>0$ in tutto un intorno di $x_{0}$, punto $x_{0}$ compreso.
\end{strategia}

\begin{dimostrazione}
La dimostrazione si fa come per le successioni: dato che $f(x_{0})>0$, possiamo scegliere
$\epsilon=f(x_{0})/2$. Per la \cref{prop:caratt-continuita}(c) esiste quindi un numero
$\delta>0$ per cui $\abs{f(x)-f(x_{0})}<f(x_{0})/2$ per ogni $x\in A$ con
$\abs{x-x_{0}}<\delta$. Ciò equivale a $-f(x_{0})/2<f(x)-f(x_{0})<f(x_{0})/2$; in
particolare
\begin{equation}\label{eq:permanenza-continue}
  f(x)>f(x_{0})-\frac{f(x_{0})}2=\frac{f(x_{0})}2>0 . \qedhere
\end{equation}
\end{dimostrazione}

Molto importanti sono i tre teoremi che seguono: il teorema dell'esistenza degli zeri,
il teorema dell'esistenza dei valori intermedi e il teorema di Weierstrass sull'esistenza
del massimo e del minimo. Sono, nello schema della \cref{fig:schema-teoremi}, i primi
teoremi di esistenza che riguardano le funzioni.

\subsection{Il teorema dell'esistenza degli zeri}

\begin{teorema}[Esistenza degli zeri]\label{teo:zeri}\index{teorema!dell'esistenza degli zeri}
Sia $f$ una funzione continua in un intervallo $[a,b]$. Se $f(a)<0$ e $f(b)>0$, allora
esiste $x_{0}\in(a,b)$ tale che $f(x_{0})=0$.
\end{teorema}

\begin{schema}
\Ipotesi{$f$ continua in $[a,b]$; \ $f(a)<0<f(b)$.}
\Tesi{$\esiste x_{0}\in(a,b)$: \ $f(x_{0})=0$.}
\end{schema}

Naturalmente la tesi vale anche se $f(a)>0$ e $f(b)<0$ (basta applicare il teorema a
$-f$): il teorema dell'esistenza degli zeri vale supponendo che i valori $f(a)$, $f(b)$
siano di segno discorde. L'enunciato garantisce l'esistenza di \emph{almeno} uno zero,
non la sua unicità: gli zeri possono essere molti (\cref{fig:zeri}). La dimostrazione
del teorema, con il metodo di bisezione, è riportata nel \cref{sec:bisezione}.

\begin{figure}[htb]
\centering
\begin{tikzpicture}
\begin{axis}[width=0.92\textwidth, height=5cm, xmin=0, xmax=6.4, ymin=-1.9, ymax=2.6,
  axis lines=middle, xtick={0.5,6}, xticklabels={$a$,$b$}, ytick=\empty,
  xlabel={$x$}, ylabel={$y$},
  every axis x label/.style={at={(axis description cs:1,0.46)},anchor=west}]
  \addplot[curva, domain=0.5:6, samples=150] {0.15*(x-1.4)*(x-3.1)*(x-4.9)};
  \addplot[only marks, mark=*, mark size=1.6pt, blu] coordinates {(0.5,-1.5444) (6,2.2011)};
  \addplot[only marks, mark=*, mark size=1.6pt, rame] coordinates {(1.4,0) (3.1,0) (4.9,0)};
  \addplot[tratteggio] coordinates {(0.5,0) (0.5,-1.5444)};
  \addplot[tratteggio] coordinates {(6,0) (6,2.2011)};
  \node[font=\scriptsize, right] at (axis cs:0.55,-1.5444) {$f(a)<0$};
  \node[font=\scriptsize, left] at (axis cs:5.95,2.2011) {$f(b)>0$};
\end{axis}
\end{tikzpicture}
\caption{Una funzione continua in $[a,b]$ che assume agli estremi valori di segno
discorde attraversa l'asse $x$ almeno una volta: qui gli zeri sono tre. Il teorema
garantisce che ce n'è almeno uno, non dice quanti sono né dove si trovano.}
\label{fig:zeri}
\end{figure}

Tutte le ipotesi sono essenziali. Senza la continuità la tesi è falsa: la funzione che
vale $-1$ in $[-1,0)$ e $1$ in $[0,1]$ cambia segno senza annullarsi. È essenziale anche
che il dominio sia un intervallo: $\frac1x$ assume valori di segno discorde in $-1$ e in
$1$, è continua nel suo dominio $\R\setminus\{0\}$, ma non si annulla mai. Infine, come
vedremo nel \cref{sec:bisezione}, è essenziale l'assioma di completezza.

Per mostrare la portata del teorema, consideriamo come esempio le due equazioni nella
incognita $x$
\begin{align}
  &x^{3}+x-1=0, \label{eq:cubica}\\
  &e^{x}+x=0, \label{eq:trascendente}
\end{align}
che non rientrano tra le equazioni algebriche di primo e secondo grado di cui è facile
ricordare la formula risolutiva. La prima delle due equazioni è un'\emph{equazione
algebrica di terzo grado}, mentre la seconda è un'\emph{equazione trascendente}.
Procediamo per tentativi, assegnando a $x$ alcuni valori:
\begin{center}
\renewcommand{\arraystretch}{1.4}
\begin{tabular}{c|ccccc}
  $x$ & $-2$ & $-1$ & $0$ & $1$ & $2$\\ \hline
  $x^{3}+x-1$ & $-11$ & $-3$ & $-1$ & $1$ & $9$\\
  $e^{x}+x$ & $e^{-2}-2$ & $e^{-1}-1$ & $1$ & $e+1$ & $e^{2}+2$
\end{tabular}
\end{center}
Nel caso $f(x)=x^{3}+x-1$, abbiamo $f(0)<0$, $f(1)>0$. Poiché $f$ è un polinomio, e
quindi è continua, in base al teorema dell'esistenza degli zeri esiste un numero $x_{0}$
nell'intervallo $(0,1)$ tale che $f(x_{0})=0$; cioè $x_{0}$ è una soluzione
dell'equazione \eqref{eq:cubica}. Nel \cref{sec:bisezione} vedremo come calcolare
numericamente tale radice, e troveremo che $x_{0}=0{,}682327\ldots$ Nel secondo caso,
$f(x)=e^{x}+x$ è continua (\cref{teo:continuita-esponenziale}) e risulta
$f(-1)=\frac1e-1<0$, $f(0)=1>0$. Quindi esiste nell'intervallo $(-1,0)$ una radice
$x_{0}$ dell'equazione \eqref{eq:trascendente}; troveremo che $x_{0}=-0{,}567143\ldots$
In entrambi i casi la radice è unica, perché le due funzioni sono strettamente crescenti
(somme di funzioni strettamente crescenti): due zeri distinti $x_{1}<x_{2}$ darebbero
$0=f(x_{1})<f(x_{2})=0$.

Esiste una formula risolutiva per le equazioni di terzo grado (si ricordi la nota storica
su Cardano e Bombelli del \cref{sec:complessi}), che dà come soluzione reale
dell'equazione \eqref{eq:cubica} il numero
\begin{equation}\label{eq:cardano}
  x_{0}=\sqrt[3]{\frac12+\sqrt{\frac{31}{108}}}-\sqrt[3]{\sqrt{\frac{31}{108}}-\frac12}
  =0{,}682327\ldots
\end{equation}
Viceversa, non è nota alcuna formula risolutiva, in termini di funzioni elementari, per
l'equazione \eqref{eq:trascendente}.

\begin{notastorica}[Le equazioni che non si risolvono con una formula]
Dopo le formule per le equazioni di terzo e quarto grado, trovate nel Cinquecento da
Scipione del Ferro, Niccolò Tartaglia, Gerolamo Cardano e Ludovico Ferrari, i matematici
cercarono per quasi tre secoli una formula «per radicali» per l'equazione generale di
quinto grado. Paolo Ruffini (1799) e, in forma completa, Niels Henrik Abel (1824)
dimostrarono che una tale formula \emph{non esiste}; Évariste Galois chiarì poco dopo
quali equazioni particolari siano risolubili per radicali. Per le equazioni
trascendenti come la \eqref{eq:trascendente} la situazione è ancora più netta: la sua
soluzione si esprime con la \emph{funzione $W$ di Lambert}, $x_{0}=-W(1)$, ma $W$ è
definita proprio come inversa di una funzione, $w\mapsto we^{w}$, e non è una funzione
elementare. Il numero $W(1)=0{,}567143\ldots$ è detto \emph{costante omega}. Il teorema
degli zeri, insieme al metodo di bisezione, cambia la prospettiva: anche quando una
formula non esiste, l'esistenza della soluzione è garantita, e la soluzione si può
calcolare con la precisione che si desidera.
\end{notastorica}

Applichiamo ancora una volta il teorema dell'esistenza degli zeri per ritrovare, per
un'altra via, un risultato del \cref{cap:complementi}: l'esistenza della radice $n$-esima
(\cref{teo:radice-n}). \emph{Per ogni
$y_{0}\ge0$ esiste un numero reale $x_{0}\ge0$ soluzione dell'equazione}
\begin{equation}\label{eq:radice-zeri}
  x^{n}=y_{0} .
\end{equation}
Ricordiamo che, essendo la funzione $x^{n}$ strettamente crescente per $x\ge0$
(\cref{prop:potenza-crescente}), un tale numero $x_{0}$ è unico. Dimostriamo che
l'equazione ha soluzione: se $n=1$ basta $x_{0}=y_{0}$, e quindi possiamo supporre
$n\ge2$; se $y_{0}=0$ naturalmente è $x_{0}=0$, e se $y_{0}=1$ basta prendere
$x_{0}=1$. Negli altri casi ($y_{0}>0$, $y_{0}\neq1$) poniamo $f(x)=x^{n}-y_{0}$, che è
continua; risulta
$f(0)=-y_{0}<0$, e rimane da trovare un punto dove $f$ è positiva. Se $y_{0}<1$, allora
$f(1)=1-y_{0}>0$, e quindi esiste una radice $x_{0}$ nell'intervallo $(0,1)$. Se invece
$y_{0}>1$, allora $f(y_{0})=y_{0}^{n}-y_{0}=y_{0}\bigl(y_{0}^{n-1}-1\bigr)>0$, e quindi in
questo caso esiste una radice nell'intervallo $(0,y_{0})$. Nel \cref{cap:complementi}
avevamo ottenuto lo stesso risultato direttamente con l'estremo superiore; il teorema
degli zeri, che a sua volta poggia sull'assioma di completezza, incorpora una volta per
tutte quel ragionamento.

\subsection{Il teorema dei valori intermedi}

\begin{teorema}[Primo teorema dei valori intermedi]\label{teo:valori-intermedi-1}\index{teorema!dei valori intermedi}
Una funzione continua in un intervallo $[a,b]$ assume tutti i valori compresi tra $f(a)$
e $f(b)$.
\end{teorema}

\begin{schema}
\Ipotesi{$f$ continua in $[a,b]$; \ $y_{0}$ compreso tra $f(a)$ e $f(b)$.}
\Tesi{$\esiste x_{0}\in[a,b]$: \ $f(x_{0})=y_{0}$.}
\end{schema}

\begin{strategia}
Si «abbassa» il grafico di $y_{0}$: la funzione $g(x)=f(x)-y_{0}$ è continua e cambia
segno tra $a$ e $b$, e uno zero di $g$ è esattamente un punto in cui $f$ vale $y_{0}$
(\cref{fig:valori-intermedi}). Il teorema dei valori intermedi è dunque il teorema degli
zeri applicato a una funzione traslata.
\end{strategia}

\begin{dimostrazione}
Per semplificare le notazioni consideriamo il caso in cui $f(a)\le f(b)$ (l'altro caso
si ottiene applicando questo a $-f$). La tesi consiste nel provare che, qualunque sia
$y_{0}\in[f(a),f(b)]$, esiste $x_{0}\in[a,b]$ tale che $f(x_{0})=y_{0}$. Se $y_{0}=f(a)$ si
può porre $x_{0}=a$; analogamente, se $y_{0}=f(b)$, basta prendere $x_{0}=b$. Per trattare
il caso $y_{0}\in(f(a),f(b))$ consideriamo la funzione
\begin{equation}\label{eq:vi-g}
  g(x)=f(x)-y_{0} \qquad\perogni x\in[a,b],
\end{equation}
che è continua per il \cref{teo:operazioni-continue}; essendo $f(a)<y_{0}<f(b)$, risulta
\begin{equation}\label{eq:vi-segni}
  g(a)=f(a)-y_{0}<0, \qquad g(b)=f(b)-y_{0}>0 .
\end{equation}
Per il teorema dell'esistenza degli zeri esiste un numero $x_{0}\in(a,b)$ tale che
$g(x_{0})=0$, cioè $f(x_{0})=y_{0}$.
\end{dimostrazione}

\begin{figure}[htb]
\centering
\begin{tikzpicture}
\begin{axis}[width=0.85\textwidth, height=5cm, xmin=0, xmax=5.4, ymin=0, ymax=3.3,
  axis lines=left, xtick={0.6,5}, xticklabels={$a$,$b$}, ytick={0.7,1.552,2.404},
  yticklabels={$f(a)$,$y_{0}$,$f(b)$}, xlabel={$x$}, ylabel={$y$}]
  \addplot[curva, domain=0.6:5, samples=120] {0.7+0.32*(x-0.6)+0.95*sin(deg(1.5*(x-0.6)))};
  \addplot[curva2, domain=0:5.4] {1.552};
  \addplot[tratteggio] coordinates {(0,0.7) (0.6,0.7) (0.6,0)};
  \addplot[tratteggio] coordinates {(0,2.404) (5,2.404) (5,0)};
  \addplot[only marks, mark=*, mark size=1.6pt, rame] coordinates {(1.1332,1.552) (2.5277,1.552) (4.5017,1.552)};
\end{axis}
\end{tikzpicture}
\caption{Il teorema dei valori intermedi: ogni retta orizzontale $y=y_{0}$ con $y_{0}$
compreso tra $f(a)$ e $f(b)$ incontra il grafico di una funzione continua su $[a,b]$
almeno una volta (qui tre volte).}
\label{fig:valori-intermedi}
\end{figure}

\begin{idea}[L'equilibrio portanza--peso]
In volo rettilineo orizzontale uniforme la portanza deve uguagliare il peso:
$\frac12\rho V^{2}S\,C_{L}(\alpha)=W$. Per una data velocità $V$ e quota (cioè densità
$\rho$), esiste un angolo d'attacco che realizza l'equilibrio? Il coefficiente di
portanza $C_{L}(\alpha)$ è una funzione continua dell'incidenza (tra l'incidenza di
portanza nulla e lo stallo), e così la portanza $L(\alpha)=\frac12\rho V^{2}S\,C_{L}(\alpha)$.
Se a un'incidenza $\alpha_{1}$ la portanza è minore del peso e a un'incidenza $\alpha_{2}$
è maggiore, il teorema dei valori intermedi garantisce un'incidenza $\bar\alpha$ compresa
tra le due per cui $L(\bar\alpha)=W$. Se invece anche al $C_{L,\max}$ la portanza resta
inferiore al peso, l'equilibrio è impossibile: la velocità è minore della velocità di
stallo $V_{S}=\sqrt{2W/(\rho S\,C_{L,\max})}$.
\end{idea}

\subsection{Il teorema di Weierstrass}

\begin{teorema}[di Weierstrass]\label{teo:weierstrass}\index{teorema!di Weierstrass}
Sia $f$ una funzione continua in un intervallo chiuso e limitato $[a,b]$. Allora $f$
assume massimo e minimo in $[a,b]$, cioè esistono in $[a,b]$ due punti $x_{1}$, $x_{2}$
tali che
\begin{equation}\label{eq:weierstrass}
  f(x_{1})\le f(x)\le f(x_{2}) \qquad\perogni x\in[a,b].
\end{equation}
\end{teorema}

\begin{schema}
\Ipotesi{$f$ continua in $[a,b]$, intervallo chiuso e limitato.}
\Tesi{$\esiste x_{1},x_{2}\in[a,b]$: \ $f(x_{1})\le f(x)\le f(x_{2})$ \ $\perogni x\in[a,b]$.}
\end{schema}

I numeri $x_{1}$, $x_{2}$ sono detti rispettivamente \emph{punti di minimo} e \emph{di
massimo} per $f$ nell'intervallo $[a,b]$; i corrispondenti valori $m=f(x_{1})$ e
$M=f(x_{2})$ sono detti \emph{minimo} e \emph{massimo} di $f$ in $[a,b]$
(\cref{fig:weierstrass}). Il massimo e il minimo sono unici (\cref{prop:unicita-max}),
i punti di massimo e di minimo in generale no.

\begin{figure}[htb]
\centering
\begin{tikzpicture}
\begin{axis}[width=0.85\textwidth, height=5.2cm, xmin=0, xmax=6.2, ymin=0, ymax=3.4,
  axis lines=left, xtick={0.6,1.62,4.6,5.6}, xticklabels={$a$,$x_{1}$,$x_{2}$,$b$},
  ytick={0.9,3}, yticklabels={$m$,$M$}, xlabel={$x$}, ylabel={$f(x)$}]
  \addplot[curva, domain=0.6:5.6, samples=150, smooth]
    coordinates {(0.6,1.5) (1.0,1.05) (1.62,0.9) (2.4,1.15) (3.2,1.9) (4.0,2.75) (4.6,3.0) (5.2,2.75) (5.6,2.1)};
  \addplot[tratteggio] coordinates {(0,0.9) (1.62,0.9) (1.62,0)};
  \addplot[tratteggio] coordinates {(0,3) (4.6,3) (4.6,0)};
  \addplot[tratteggio] coordinates {(0.6,0) (0.6,1.5)};
  \addplot[tratteggio] coordinates {(5.6,0) (5.6,2.1)};
  \addplot[only marks, mark=*, mark size=1.6pt, rame] coordinates {(1.62,0.9) (4.6,3)};
\end{axis}
\end{tikzpicture}
\caption{Il teorema di Weierstrass: una funzione continua su un intervallo chiuso e
limitato assume un valore minimo $m=f(x_{1})$ e un valore massimo $M=f(x_{2})$.}
\label{fig:weierstrass}
\end{figure}

Il teorema di Weierstrass è dimostrato nel \cref{sec:dim-weierstrass}; in questa sede ci
limitiamo a mettere in luce con degli esempi l'importanza delle ipotesi (funzione
\emph{continua} definita in un intervallo \emph{chiuso} e \emph{limitato}) che
garantiscono l'esistenza del massimo e del minimo (\cref{fig:controesempi-weierstrass}).

\begin{enumerate}[label=\textup{(\alph*)}]
  \item Consideriamo per $x>0$ la funzione $f(x)=\frac1x$. Essa \emph{non assume
  massimo} nell'intervallo aperto a sinistra $(0,1]$; infatti non è limitata
  superiormente in tale intervallo: per ogni $M>0$ risulta $f(x)>M$ se
  $x\in\bigl(0,\frac1M\bigr)$. Qui manca la chiusura dell'intervallo.
  \item La stessa funzione \emph{non assume minimo} nell'intervallo illimitato
  $[1,+\infty)$. La funzione è limitata in tale intervallo, perché $0<\frac1x\le1$ per
  ogni $x\ge1$; però $0$ non è il minimo di $f$ in $[1,+\infty)$, perché non esiste alcun
  $x_{1}\ge1$ per cui $f(x_{1})=\frac1{x_{1}}=0$ (una frazione è nulla se e soltanto se
  il numeratore è nullo), e nessun numero reale positivo $m$ è minimo per $f$, perché
  risulta $f(x)<m$ per ogni $x>\frac1m$. L'estremo inferiore è $0$, ma non è raggiunto:
  qui manca la limitatezza dell'intervallo.
  \item La funzione $f(x)=x^{2}$ assume massimo e minimo in ogni intervallo chiuso e
  limitato $[a,b]$, in particolare nell'intervallo $[-1,1]$. Invece la funzione
  \begin{equation}\label{eq:g-senza-minimo}
    g(x)=
    \begin{cases}
      x^{2} & \text{se } x\in[-1,1]\setminus\{0\},\\
      1 & \text{se } x=0
    \end{cases}
  \end{equation}
  non assume minimo nell'intervallo chiuso e limitato $[-1,1]$, perché assume valori
  positivi arbitrariamente vicini allo zero ($g(x)\to0$ per $x\to0$) ma non è uguale a
  zero per alcun valore di $x\in[-1,1]$. In questo caso la mancanza del minimo è causata
  dalla discontinuità (eliminabile) di $g$ nel punto $x_{0}=0$.
  \item Per la sua discontinuità in corrispondenza dei numeri interi, la funzione
  \emph{parte frazionaria} $f(x)=x-\lfloor x\rfloor$, $x\in\R$, non ha massimo in un
  qualunque intervallo $[a,b]$ che contenga un numero intero $k$ con $a<k\le b$: nei
  punti immediatamente a sinistra di $k$ essa assume valori arbitrariamente vicini a $1$
  (si ha $f(x)=x-k+1$ per $k-1\le x<k$), ma il valore $1$ non è mai assunto, essendo
  $0\le f(x)<1$ per ogni $x$. Si osservi invece che in $[0,\frac12]$, che contiene l'intero
  $0$ ma solo come estremo sinistro, $f(x)=x$ ha massimo $\frac12$.
\end{enumerate}

\begin{figure}[htb]
\centering
\includegraphics[width=0.95\textwidth]{controesempi_weierstrass}
\caption{Quattro funzioni per cui la tesi del teorema di Weierstrass non vale, perché
manca una delle ipotesi: (a) l'intervallo non è chiuso; (b) l'intervallo non è limitato;
(c) e (d) la funzione non è continua (\texttt{python/fig\_controesempi\_weierstrass.py}).}
\label{fig:controesempi-weierstrass}
\end{figure}

Le ipotesi del teorema sono sufficienti, non necessarie: la funzione $\sen x$ ha massimo
$1$ e minimo $-1$ su tutto $\R$, che non è limitato; la funzione di Dirichlet ha massimo
e minimo su ogni intervallo pur essendo discontinua ovunque. Il teorema di Weierstrass
garantisce l'esistenza del massimo, ma non dice dove si trovi: per trovarlo occorrono
altri strumenti (il calcolo differenziale, oppure, come nell'esempio che segue,
un'osservazione algebrica, oppure una ricerca numerica).

\begin{idea}[L'efficienza massima di un velivolo]
Con la \emph{polare parabolica} $C_{D}=C_{D0}+kC_{L}^{2}$, dove $k=\frac1{\pi e\,AR}$
dipende dall'allungamento alare $AR$ e dal fattore di Oswald $e$, l'efficienza
aerodinamica
\[
  E(C_{L})=\frac{C_{L}}{C_{D}}=\frac{C_{L}}{C_{D0}+kC_{L}^{2}}, \qquad 0\le C_{L}\le C_{L,\max},
\]
è una funzione continua su un intervallo chiuso e limitato (il denominatore è positivo):
per il teorema di Weierstrass \emph{esiste} un coefficiente di portanza $C_{L}^{*}$ di
efficienza massima. Per trovarlo, osserviamo che per $C_{L}>0$
\[
  \frac1{E}=\frac{C_{D0}}{C_{L}}+kC_{L}\ \ge\ 2\sqrt{\frac{C_{D0}}{C_{L}}\cdot kC_{L}}=2\sqrt{kC_{D0}},
\]
per la disuguaglianza $u+v\ge2\sqrt{uv}$, valida per $u,v>0$ perché equivale a
$\bigl(\sqrt u-\sqrt v\bigr)^{2}\ge0$, con l'uguaglianza se e solo se $u=v$. Dunque
$E\le E_{\max}=\frac1{2\sqrt{kC_{D0}}}$, con l'uguaglianza se e solo se
$\frac{C_{D0}}{C_{L}}=kC_{L}$, cioè per $C_{L}^{*}=\sqrt{C_{D0}/k}$, la condizione in cui
la resistenza parassita uguaglia quella indotta (purché $C_{L}^{*}\le C_{L,\max}$).
Il programma che segue confronta la formula con una ricerca del massimo su una griglia
fitta, per un velivolo con $C_{D0}=0{,}025$, $AR=8$, $e=0{,}8$ (\cref{fig:efficienza}).
\end{idea}

\begin{python}[Efficienza massima: Weierstrass, poi la formula]
import math

# Polare parabolica: C_D = C_D0 + k C_L^2, con k = 1/(pi e AR)
CD0, AR, e_osw, CLmax = 0.025, 8.0, 0.80, 1.4
k = 1 / (math.pi * e_osw * AR)
E = lambda CL: CL / (CD0 + k * CL**2)       # efficienza L/D

# Weierstrass: E è continua su [0, CLmax], quindi ha massimo.
# Lo cerchiamo su una griglia fitta e lo confrontiamo con la formula.
N = 140000
CL_griglia = max((CLmax * i / N for i in range(N + 1)), key=E)
CL_star = math.sqrt(CD0 / k)
print(f"k = {k:.5f}")
print(f"griglia:  C_L = {CL_griglia:.4f}   E = {E(CL_griglia):.4f}")
E_max = 1 / (2 * math.sqrt(CD0 * k))
print(f"formula:  C_L = {CL_star:.4f}   E = {E_max:.4f}")
\end{python}
\begin{risultato}
k = 0.04974
griglia:  C_L = 0.7090   E = 14.1796
formula:  C_L = 0.7090   E = 14.1796
\end{risultato}

\begin{figure}[htb]
\centering
\includegraphics[width=0.92\textwidth]{efficienza}
\caption{L'efficienza $E=C_{L}/C_{D}$ in funzione di $C_{L}$ per la polare parabolica
dell'esempio: il massimo, che esiste per il teorema di Weierstrass, è
$E_{\max}\approx14{,}18$ per $C_{L}^{*}\approx0{,}709$ (\texttt{python/fig\_efficienza.py}).}
\label{fig:efficienza}
\end{figure}

\subsection{Il secondo teorema dei valori intermedi}

Siamo ora in grado di provare una nuova formulazione del teorema di esistenza dei valori
intermedi.

\begin{teorema}[Secondo teorema dei valori intermedi]\label{teo:valori-intermedi-2}
Una funzione continua in un intervallo $[a,b]$ assume tutti i valori compresi tra il
minimo e il massimo.
\end{teorema}

\begin{schema}
\Ipotesi{$f$ continua in $[a,b]$; \ $m=\min f$, $M=\max f$ (che esistono per Weierstrass).}
\Tesi{per ogni $y_{0}\in[m,M]$ \ $\esiste x_{0}\in[a,b]$: \ $f(x_{0})=y_{0}$.}
\end{schema}

\begin{strategia}
Si ripete la dimostrazione del primo teorema dei valori intermedi, ma al posto degli
estremi $a$, $b$ si usano un punto di minimo $x_{1}$ e un punto di massimo $x_{2}$, che
esistono per il teorema di Weierstrass: la funzione $f(x)-y_{0}$ cambia segno tra
$x_{1}$ e $x_{2}$, e il teorema degli zeri, applicato sull'intervallo di estremi $x_{1}$
e $x_{2}$, fornisce il punto cercato.
\end{strategia}

\begin{dimostrazione}
I valori di massimo $M$ e di minimo $m$ sono assunti in base al teorema di Weierstrass;
rimane da provare che, qualunque sia $y_{0}\in(m,M)$, esiste $x_{0}\in[a,b]$ tale che
$f(x_{0})=y_{0}$. Indichiamo con $x_{1}$, $x_{2}$ i punti di minimo e di massimo di $f$,
cioè tali che $f(x_{1})=m$, $f(x_{2})=M$, e consideriamo la funzione
\begin{equation}\label{eq:vi2-g}
  g(x)=f(x)-y_{0} \qquad\perogni x\in[a,b].
\end{equation}
Essendo $f(x_{1})=m<y_{0}<M=f(x_{2})$, risulta
\begin{equation}\label{eq:vi2-segni}
  g(x_{1})=f(x_{1})-y_{0}<0, \qquad g(x_{2})=f(x_{2})-y_{0}>0;
\end{equation}
per il teorema dell'esistenza degli zeri, applicato a $g$ nell'intervallo di estremi
$x_{1}$, $x_{2}$ (che è $[x_{1},x_{2}]$ oppure $[x_{2},x_{1}]$, e su cui $g$ è continua),
esiste un numero $x_{0}$, appartenente all'intervallo aperto di estremi $x_{1}$, $x_{2}$,
tale che $g(x_{0})=0$, cioè tale che $f(x_{0})=y_{0}$.
\end{dimostrazione}

\begin{corollario}\label{cor:immagine-intervallo}
L'immagine di un intervallo chiuso e limitato mediante una funzione continua è un
intervallo chiuso e limitato: se $f$ è continua in $[a,b]$, allora
$f([a,b])=[m,M]$, dove $m$ ed $M$ sono il minimo e il massimo di $f$.
\end{corollario}

\begin{schema}
\Ipotesi{$f$ continua in $[a,b]$; \ $m=\min f$, $M=\max f$.}
\Tesi{$f([a,b])=[m,M]$.}
\end{schema}

\begin{strategia}
Si dimostrano le due inclusioni: $f([a,b])\subseteq[m,M]$ è la definizione stessa di
minimo e massimo, $[m,M]\subseteq f([a,b])$ è il secondo teorema dei valori intermedi.
\end{strategia}

\begin{dimostrazione}
Ogni valore $f(x)$ sta in $[m,M]$ per definizione di minimo e massimo; viceversa ogni
$y_{0}\in[m,M]$ è un valore di $f$ per il \cref{teo:valori-intermedi-2}.
\end{dimostrazione}

\begin{notastorica}[Una curva che riempie un quadrato]
Una funzione continua trasforma dunque un segmento in un segmento, eventualmente ridotto
a un punto. Che cosa accade se i valori stanno nel piano, cioè se si considerano due
funzioni continue $x(t)$, $y(t)$, con $t\in[0,1]$? L'intuizione suggerisce che il punto
$\bigl(x(t),y(t)\bigr)$ descriva una linea, priva di area. Nel 1878 Georg Cantor aveva
sorpreso i contemporanei dimostrando che tra il segmento e il quadrato esiste una
corrispondenza biunivoca (lo abbiamo raccontato nella nota storica del
\S\ref{sec:invertibili}); l'anno seguente Eugen Netto provò che una tale corrispondenza
non può essere continua. Nel 1890 Giuseppe Peano pubblicò sui \emph{Mathematische
Annalen} una nota di quattro pagine, \emph{Sur une courbe, qui remplit toute une aire
plane}: vi costruiva due funzioni continue $x(t)$, $y(t)$ tali che la coppia
$\bigl(x(t),y(t)\bigr)$, al variare di $t$ in $[0,1]$, assume \emph{tutti} i punti del
quadrato $[0,1]\times[0,1]$. Peano le definì per via aritmetica, con lo sviluppo di $t$
in base tre, senza alcuna figura; l'anno dopo David Hilbert ne pubblicò una variante
geometrica, e fu il primo a illustrarla con un disegno (\cref{fig:curva-hilbert}). Per
il teorema di Netto una curva siffatta non può essere iniettiva: passa più volte per
alcuni punti. La lezione è che la continuità, da sola, non conserva la dimensione; la
conservano le corrispondenze biunivoche e continue insieme alla loro inversa, come
dimostrò Brouwer nel 1911.
\end{notastorica}

\begin{figure}[htb]
\centering
\includegraphics[width=\textwidth]{curva_hilbert}
\caption{Le prime quattro iterate della curva di Hilbert. La $n$-esima è una poligonale
che percorre, uno dopo l'altro, i centri dei $4^{n}$ quadratini di lato $2^{-n}$; ogni
quadratino viene poi diviso in quattro, che l'iterata successiva percorre
consecutivamente. Al crescere di $n$ le poligonali, opportunamente parametrizzate su
$[0,1]$, convergono a una funzione continua che assume tutti i punti del quadrato
(\texttt{python/fig\_curva\_hilbert.py}).}
\label{fig:curva-hilbert}
\end{figure}

\subsection{Un criterio di invertibilità}

Chiudiamo il paragrafo precisando un criterio, introdotto nel \cref{sec:invertibili},
per riconoscere se una data funzione è invertibile. Là avevamo osservato che una
funzione strettamente monotòna è iniettiva (\cref{teo:monotona-invertibile}); la
continuità permette ora di individuarne esattamente l'immagine. La continuità della
funzione inversa è invece studiata nel \cref{sec:monotone-inverse}.

\begin{teorema}[Criterio di invertibilità]\label{teo:invertibilita}\index{funzione!invertibile}
Una funzione continua e strettamente monotòna in un intervallo $[a,b]$ è invertibile in
tale intervallo: se è strettamente crescente, $f:[a,b]\to[f(a),f(b)]$ è biiettiva.
\end{teorema}

\begin{schema}
\Ipotesi{$f$ continua e strettamente crescente (o decrescente) in $[a,b]$.}
\Tesi{$f:[a,b]\to[f(a),f(b)]$ (rispettivamente $[f(b),f(a)]$) è biiettiva.}
\end{schema}

\begin{strategia}
La biiettività si spezza in due: la suriettività sull'intervallo di estremi $f(a)$,
$f(b)$ viene dalla continuità, attraverso il primo teorema dei valori intermedi;
l'iniettività viene dalla stretta monotonia, perché punti distinti hanno immagini
ordinate allo stesso modo, e quindi distinte.
\end{strategia}

\begin{dimostrazione}
Proponiamo la dimostrazione nel caso in cui la funzione $f$ sia strettamente crescente
in $[a,b]$ (il caso decrescente è analogo, con $[f(b),f(a)]$). Risulta
\begin{equation}\label{eq:invertibilita-estremi}
  f(a)<f(x)<f(b) \qquad\perogni x\in(a,b);
\end{equation}
quindi $f(a)$ è il minimo di $f$ in $[a,b]$, mentre $f(b)$ è il massimo, e
$f([a,b])\subseteq[f(a),f(b)]$. Inoltre, per il \cref{teo:valori-intermedi-1}, $f$
assume tutti i valori compresi tra $f(a)$ e $f(b)$: per ogni $y\in[f(a),f(b)]$ esiste
almeno un $x\in[a,b]$ per cui $f(x)=y$. Tale $x$ è unico; infatti, se esistessero due
valori $x_{1}$, $x_{2}$ distinti tra loro, diciamo $x_{1}<x_{2}$, per cui
$y=f(x_{1})=f(x_{2})$, allora dovrebbe risultare anche $f(x_{1})<f(x_{2})$, dato che $f$ è
strettamente crescente. Quindi $f:[a,b]\to[f(a),f(b)]$ è suriettiva e iniettiva, cioè
invertibile.
\end{dimostrazione}

%% ----------------------------------------------------------------------------
%%  §4.8  Appendice: il metodo di bisezione
%% ----------------------------------------------------------------------------
\section{Appendice: il metodo di bisezione per il calcolo delle radici}
\label{sec:bisezione}

In questo paragrafo dimostriamo il teorema dell'esistenza degli zeri (\cref{teo:zeri}).
Utilizziamo nella dimostrazione il \emph{metodo di bisezione}, già incontrato nel
\cref{cap:numeri} per approssimare $\sqrt2$ e nel \cref{cap:successioni} per dimostrare
il teorema di Bolzano--Weierstrass: si tratta di un procedimento costruttivo che, oltre a
dimostrare l'esistenza di una soluzione di un'equazione data, fornisce anche un metodo
per calcolarla.

Prendiamo in considerazione \emph{equazioni} del tipo
\begin{equation}\label{eq:equazione}
  f(x)=0,
\end{equation}
con $f$ funzione definita in un intervallo $[a,b]$. Risolvere l'equazione significa
determinare tutti i numeri reali $x_{0}\in[a,b]$ per cui $f(x_{0})=0$; tali numeri si
dicono \emph{soluzioni} dell'equazione \eqref{eq:equazione}, o anche \emph{zeri} della
funzione $f$. Se $f$ è un polinomio, si dice che \eqref{eq:equazione} è un'\emph{equazione
algebrica}; se $f$ è una funzione trascendente (ad esempio composta tramite le funzioni
$e^{x}$, $\log x$, $\sen x$, $\cos x$), la \eqref{eq:equazione} prende il nome di
\emph{equazione trascendente}. Una soluzione di un'equazione algebrica si dice anche
\emph{radice} dell'equazione; per estensione, si usa il termine di radici anche per le
soluzioni di equazioni trascendenti.

\begin{strategia}[Piano d'attacco per il \cref{teo:zeri}]
Ricordiamo le ipotesi: $f$ è continua in $[a,b]$ e $f(a)<0<f(b)$. L'idea è di
\emph{intrappolare} uno zero in intervalli sempre più piccoli, mantenendo l'invariante
«$f$ è negativa nell'estremo sinistro e positiva nell'estremo destro». Si dimezza
l'intervallo e si tiene la metà agli estremi della quale $f$ continua a cambiare segno.
Gli estremi sinistri formano una successione crescente e limitata, che converge per il
teorema sulle successioni monotòne (qui entra l'assioma di completezza); gli estremi
destri convergono allo stesso limite $x_{0}$, perché le lunghezze si dimezzano. Infine
la continuità di $f$ permette di passare al limite nelle disuguaglianze $f(a_{n})<0$ e
$f(b_{n})>0$, che diventano $f(x_{0})\le0$ e $f(x_{0})\ge0$.
\end{strategia}

\begin{dimostrazione}[Dimostrazione del \cref{teo:zeri}]
\emph{Costruzione.} Consideriamo il numero $c$, punto di mezzo dell'intervallo $[a,b]$,
cioè $c=\frac{a+b}2$. Se $f(c)=0$ abbiamo trovato una radice. Altrimenti consideriamo i
due casi $f(c)>0$, $f(c)<0$. Se $f(c)>0$, la funzione $f$ assume valori di segno discorde
agli estremi dell'intervallo $[a,c]$, mentre se $f(c)<0$ è $[c,b]$ l'intervallo dove $f$
cambia segno. Indichiamo con $[a_{1},b_{1}]$ l'intervallo da considerare, cioè definiamo
\begin{equation}\label{eq:bisezione-primo}
  \begin{cases}
    \text{se } f(c)>0 \ \implica\ a_{1}=a,\ b_{1}=c;\\
    \text{se } f(c)<0 \ \implica\ a_{1}=c,\ b_{1}=b.
  \end{cases}
\end{equation}
Abbiamo così trovato un intervallo $[a_{1},b_{1}]$, di ampiezza metà del precedente
$[a,b]$, per cui risulta $f(a_{1})<0$, $f(b_{1})>0$. Definiamo $c_{1}=\frac{a_{1}+b_{1}}2$ e
ripetiamo il ragionamento. Posto $a_{0}=a$, $b_{0}=b$, otteniamo tre successioni $a_{n}$,
$b_{n}$, $c_{n}$ definite per ricorrenza da $c_{n}=\frac{a_{n}+b_{n}}2$ e
\begin{equation}\label{eq:bisezione-ricorrenza}
  \begin{cases}
    \text{se } f(c_{n})>0 \ \implica\ a_{n+1}=a_{n},\ b_{n+1}=c_{n};\\
    \text{se } f(c_{n})<0 \ \implica\ a_{n+1}=c_{n},\ b_{n+1}=b_{n}.
  \end{cases}
\end{equation}
Se per qualche $n$ risulta $f(c_{n})=0$, ci si ferma perché si è trovata una radice, e
il teorema è dimostrato. Altrimenti le successioni sono definite per ogni $n$ e, per
costruzione, risulta
\begin{equation}\label{eq:bisezione-segni}
  f(a_{n})<0, \qquad f(b_{n})>0 \qquad\perogni n\ge0 .
\end{equation}

\emph{Lunghezze.} È semplice scrivere la relazione che lega $a_{n}$ con $b_{n}$: infatti,
a ogni iterazione, la lunghezza dell'intervallo $[a_{n},b_{n}]$ si dimezza. Quindi
$b_{1}-a_{1}=\frac{b-a}2$, $b_{2}-a_{2}=\frac{b-a}{2^{2}}$ e, dopo $n$ passi (per
induzione),
\begin{equation}\label{eq:bisezione-lunghezza}
  b_{n}-a_{n}=\frac{b-a}{2^{n}} \qquad\perogni n\ge0 .
\end{equation}

\emph{Convergenza.} Per costruzione, la successione $a_{n}$ è crescente
($a_{0}\le a_{1}\le a_{2}\le\cdots$), la successione $b_{n}$ è decrescente, e
$a\le a_{n}<b_{n}\le b$ per ogni $n$: in particolare $a_{n}$ è limitata. Per il teorema
sulle successioni monotòne (\cref{teo:monotone}) $a_{n}$ ammette limite finito; sia
$x_{0}$ tale limite. Anche la successione $b_{n}$, espressa mediante la
\eqref{eq:bisezione-lunghezza} da
\begin{equation}\label{eq:bisezione-bn}
  b_{n}=a_{n}+\frac{b-a}{2^{n}},
\end{equation}
converge a $x_{0}$ per $n\to+\infty$, perché $\frac1{2^{n}}\to0$ (\cref{teo:lim-geometrica}).
Inoltre $x_{0}\in[a,b]$, per il \cref{cor:confronto-limiti} applicato alle
disuguaglianze $a\le a_{n}\le b$.

\emph{Passaggio al limite.} Quindi, ricordando la \eqref{eq:bisezione-segni}, dalla
continuità di $f$ nel punto $x_{0}$ (\cref{prop:caratt-continuita}(b)) e dal
\cref{cor:limite-non-negativo} si ottiene
\begin{equation}\label{eq:bisezione-limite}
  f(x_{0})=\lim_{n\to+\infty}f(a_{n})\le0, \qquad f(x_{0})=\lim_{n\to+\infty}f(b_{n})\ge0 .
\end{equation}
Perciò $f(x_{0})=0$. Infine $x_{0}\neq a$ e $x_{0}\neq b$, perché $f(a)<0$ e $f(b)>0$;
dunque $x_{0}\in(a,b)$, e il teorema dell'esistenza degli zeri è provato.
\end{dimostrazione}

\begin{strategia}[Un'altra via]
Invece di costruire lo zero, lo si individua con l'estremo superiore: si considera
l'insieme $E$ dei punti di $[a,b]$ in cui $f$ è negativa, e il suo estremo superiore $y$.
Se fosse $f(y)<0$, oppure $f(y)>0$, la permanenza del segno darebbe un intorno di $y$ in
cui $f$ conserva il segno, e in entrambi i casi $y$ non potrebbe essere il minimo dei
maggioranti di $E$. L'assioma di completezza interviene direttamente, nella forma
dell'esistenza dell'estremo superiore; la continuità, soltanto attraverso la permanenza
del segno.
\end{strategia}

\begin{dimostrazione}[Seconda dimostrazione]
Sia
\begin{equation}\label{eq:zeri-insieme-E}
  E=\{x\in[a,b] : f(x)<0\} .
\end{equation}
L'insieme $E$ non è vuoto, perché $a\in E$, ed è limitato superiormente da $b$; per il
\cref{teo:esistenza-sup} esiste $y=\sup E$, e $a\le y\le b$. Dimostriamo che $f(y)=0$,
escludendo gli altri due casi.

\emph{Non può essere $f(y)<0$.} In tal caso $y\neq b$, perché $f(b)>0$; dunque $y<b$. La
funzione $-f$ è continua in $y$ (\cref{teo:operazioni-continue}) e $-f(y)>0$: per la
permanenza del segno (\cref{teo:permanenza-segno-continue}) esiste $\delta>0$ tale che
$f(x)<0$ per ogni $x\in[a,b]$ con $\abs{x-y}<\delta$. Il punto
$x_{1}=\min\bigl\{y+\frac\delta2,b\bigr\}$ appartiene a $[a,b]$ e verifica
$y<x_{1}<y+\delta$; quindi $f(x_{1})<0$, cioè $x_{1}\in E$. Ma $x_{1}>y$, contro il fatto
che $y$ è un maggiorante di $E$.

\emph{Non può essere $f(y)>0$.} Per la permanenza del segno esiste $\delta>0$ tale che
$f(x)>0$ per ogni $x\in[a,b]$ con $\abs{x-y}<\delta$. Per la caratterizzazione
dell'estremo superiore (\cref{prop:caratt-sup}), con $\epsilon=\delta$, esiste $x_{2}\in E$
tale che $y-\delta<x_{2}\le y$. Allora $\abs{x_{2}-y}<\delta$, e quindi $f(x_{2})>0$;
d'altra parte $f(x_{2})<0$, perché $x_{2}\in E$: è una contraddizione.

Resta $f(y)=0$. Infine $y\neq a$ e $y\neq b$, perché $f(a)<0<f(b)$; quindi $y\in(a,b)$.
\end{dimostrazione}

\begin{osservazione}[Due dimostrazioni a confronto]\label{oss:zeri-due-vie}
Le due dimostrazioni poggiano entrambe sull'assioma di completezza, ma lo usano in modo
diverso. La bisezione passa attraverso il teorema sulle successioni monotòne ed è
\emph{costruttiva}: produce le approssimazioni $c_{n}$ e ne controlla l'errore. La
seconda usa l'esistenza dell'estremo superiore e non indica alcun modo per calcolare
$\sup E$: garantisce lo zero senza dire come avvicinarlo. In compenso ne individua uno
ben preciso: $y=\sup E$ è uno zero a destra del quale $f$ non assume più valori
negativi, perché i punti di $(y,b]$ non stanno in $E$. Infine, la seconda dimostrazione
usa la continuità solo «localmente», attraverso la permanenza del segno nel punto $y$:
la stessa struttura ritorna in molte dimostrazioni di esistenza dell'Analisi.
\end{osservazione}

\begin{notastorica}[Bolzano, Cauchy e il teorema degli zeri]
Il saggio di Bolzano del 1817, ricordato nel \cref{sec:sup} per il principio
dell'estremo superiore, ha per titolo l'enunciato stesso del nostro teorema:
\emph{Dimostrazione puramente analitica del teorema che fra due valori che danno
risultati di segno opposto giace almeno una radice reale dell'equazione}. Bolzano vi
definì la continuità in modo sostanzialmente identico al nostro, e ricavò il teorema
dall'esistenza dell'estremo superiore: è l'idea della seconda dimostrazione. Quattro
anni dopo Cauchy, nel \emph{Cours d'analyse}, si accontentò nel testo di un argomento geometrico
(il grafico deve incontrare la retta), ma in una nota dedicata alla risoluzione numerica
delle equazioni diede una dimostrazione diretta: divideva l'intervallo in $m$ parti
uguali, ne sceglieva una ai cui estremi la funzione ha segni opposti e ripeteva
l'operazione, ottenendo due successioni monotòne di estremi con lo stesso limite. È la
nostra prima dimostrazione, che corrisponde a $m=2$; con $m=10$ si guadagna una cifra
decimale a ogni passo. Né l'uno né l'altro, però, poteva giustificare fino in fondo il
passo decisivo, l'esistenza dell'estremo superiore o del limite di una successione
monotòna e limitata: mancava una teoria dei numeri reali, che arrivò mezzo secolo più
tardi (si veda la nota storica del \cref{sec:bw-cauchy}).
\end{notastorica}

\begin{figure}[htb]
\centering
\includegraphics[width=0.92\textwidth]{bisezione}
\caption{Il metodo di bisezione per $f(x)=x^{3}+x-1$ su $[0,1]$: i primi intervalli
incapsulati $[a_{n},b_{n}]$ si stringono attorno alla radice $x_{0}=0{,}6823\ldots$, e a
ogni passo la loro lunghezza si dimezza (\texttt{python/fig\_bisezione.py}).}
\label{fig:bisezione}
\end{figure}

\subsection{Il calcolo numerico della radice}

Dalla dimostrazione proposta risulta chiaro come calcolare numericamente la soluzione
$x_{0}$. Infatti le tre successioni $a_{n}$, $b_{n}$, $c_{n}$ convergono a $x_{0}$ (la terza
perché è compresa tra le prime due). I termini delle tre successioni sono valori
approssimati di $x_{0}$; in particolare, i valori di $a_{n}$ sono approssimazioni per
difetto, quelli di $b_{n}$ sono approssimazioni per eccesso, cioè
\begin{equation}\label{eq:bisezione-stima}
  a_{n}\le x_{0}\le b_{n} \qquad\perogni n\ge0
\end{equation}
(la successione crescente $a_{n}$ non supera il suo limite, la decrescente $b_{n}$ non
scende sotto il suo, per il \cref{teo:monotone}). Dalle \eqref{eq:bisezione-lunghezza},
\eqref{eq:bisezione-stima} si deduce che l'errore di approssimazione che si commette
sostituendo $x_{0}$ con $a_{n}$ (oppure con $b_{n}$) non supera $\frac{b-a}{2^{n}}$. Dato
che $c_{n}$ è il punto di mezzo dell'intervallo $[a_{n},b_{n}]$, che contiene $x_{0}$,
l'errore che si commette nell'approssimare $x_{0}$ con $c_{n}$ non supera
$\frac{b-a}{2^{n+1}}$:
\begin{equation}\label{eq:bisezione-errore}
  \abs{c_{n}-x_{0}}\le\frac{b_{n}-a_{n}}2=\frac{b-a}{2^{n+1}} .
\end{equation}

Riprendiamo in considerazione le equazioni \eqref{eq:cubica}, \eqref{eq:trascendente}.
Ci proponiamo il calcolo delle rispettive radici con un errore inferiore a $10^{-3}$. In
entrambi i casi abbiamo un intervallo di ampiezza $b-a=1$; infatti nel primo caso
$[a,b]=[0,1]$, nel secondo $[a,b]=[-1,0]$. L'errore di approssimazione che si commette
sostituendo la soluzione $x_{0}$ con $c_{n}$ non supera $\frac1{2^{n+1}}$; in particolare,
per $n=9$, risulta
\begin{equation}\label{eq:bisezione-c9}
  \abs{c_{9}-x_{0}}\le\frac1{2^{10}}=\frac1{1024}<10^{-3} .
\end{equation}
Il programma seguente esegue il calcolo, e alla fine prosegue fino a $50$ dimezzamenti
per confrontare il risultato con la formula di Cardano \eqref{eq:cardano}.

\begin{python}[Il metodo di bisezione]
import math

def bisezione(f, a, b, passi):
    """Se f(a) < 0 < f(b), restituisce i punti medi c_0, ..., c_n."""
    assert f(a) < 0 < f(b)
    medi = []
    for n in range(passi + 1):
        c = (a + b) / 2
        medi.append(c)
        if f(c) == 0:          # radice esatta: ci si ferma
            break
        if f(c) > 0:
            b = c              # il cambio di segno è in [a, c]
        else:
            a = c              # il cambio di segno è in [c, b]
    return medi

f1 = lambda x: x**3 + x - 1           # f1(0) = -1 < 0 < 1 = f1(1)
f2 = lambda x: math.exp(x) + x        # f2(-1) = 1/e - 1 < 0 < f2(0)

c1 = bisezione(f1, 0, 1, 9)
c2 = bisezione(f2, -1, 0, 9)
for n in range(10):
    print(f"c_{n} = {c1[n]:.6f}     c_{n} = {c2[n]:+.6f}")

# 50 passi: errore < 2^-51; confronto con la formula di Cardano
r1 = bisezione(f1, 0, 1, 50)[-1]
s = math.sqrt(31 / 108)
cardano = (0.5 + s) ** (1/3) - (s - 0.5) ** (1/3)
print(f"x^3 + x - 1 = 0:  {r1:.15f}   (Cardano: {cardano:.15f})")
print(f"e^x + x = 0:     {bisezione(f2, -1, 0, 50)[-1]:.15f}")
\end{python}
\begin{risultato}
c_0 = 0.500000     c_0 = -0.500000
c_1 = 0.750000     c_1 = -0.750000
c_2 = 0.625000     c_2 = -0.625000
c_3 = 0.687500     c_3 = -0.562500
c_4 = 0.656250     c_4 = -0.593750
c_5 = 0.671875     c_5 = -0.578125
c_6 = 0.679688     c_6 = -0.570312
c_7 = 0.683594     c_7 = -0.566406
c_8 = 0.681641     c_8 = -0.568359
c_9 = 0.682617     c_9 = -0.567383
x^3 + x - 1 = 0:  0.682327803828020   (Cardano: 0.682327803828019)
e^x + x = 0:     -0.567143290409784
\end{risultato}

Quindi la radice dell'equazione $x^{3}+x-1=0$ soddisfa $\abs{x_{0}-0{,}682617}<0{,}001$,
cioè $0{,}6816<x_{0}<0{,}6837$; la radice dell'equazione $e^{x}+x=0$ soddisfa
$\abs{x_{0}+0{,}567383}<0{,}001$, cioè $-0{,}5684<x_{0}<-0{,}5663$. Il numero $\pm0{,}001$ è
una \emph{stima} dell'errore, non l'errore effettivo: quest'ultimo, come mostrano le
ultime righe dell'output, è di circa $2{,}9\cdot10^{-4}$ nel primo caso e di circa
$2{,}4\cdot10^{-4}$ nel secondo. Si osservi anche
che la successione $c_{n}$ non è monotòna: si avvicina a $x_{0}$ oscillando, ma la
\eqref{eq:bisezione-errore} ne controlla l'errore a ogni passo.

\begin{osservazione}[Quanti passi servono?]\label{oss:passi-bisezione}
Per ottenere $\abs{c_{n}-x_{0}}<\tau$ basta, per la \eqref{eq:bisezione-errore}, che
$\frac{b-a}{2^{n+1}}<\tau$, cioè $2^{n+1}>\frac{b-a}\tau$. Poiché $2^{10}=1024\approx10^{3}$,
ogni dieci dimezzamenti si guadagnano circa tre cifre decimali esatte: il metodo è
lento ma \emph{garantito}, e il numero di passi necessari si conosce in anticipo, prima
di cominciare il calcolo. È questa robustezza a renderlo ancora oggi il metodo di
riferimento, spesso combinato con metodi più veloci che da soli potrebbero non
convergere.
\end{osservazione}

\begin{python}[Gli zeri con la libreria \texttt{scipy}]
import math
from scipy.optimize import bisect, brentq

f = lambda x: math.exp(x) + x

# entrambe vogliono [a, b] con f(a), f(b) di segno discorde
x_b, info_b = bisect(f, -1, 0, xtol=1e-14, full_output=True)
x_r, info_r = brentq(f, -1, 0, xtol=1e-14, full_output=True)
print(f"bisect: {x_b:.14f}  ({info_b.iterations} iterazioni)")
print(f"brentq: {x_r:.14f}  ({info_r.iterations} iterazioni)")
try:
    bisect(f, 0, 1)            # f(0) = 1, f(1) = e + 1: stesso segno
except ValueError as err:
    print("errore:", err)
\end{python}
\begin{risultato}
bisect: -0.56714329040978  (47 iterazioni)
brentq: -0.56714329040978  (6 iterazioni)
errore: f(a) and f(b) must have different signs
\end{risultato}

La funzione \texttt{bisect} è esattamente il metodo della dimostrazione: con
$b-a=1$ e tolleranza $10^{-14}$ servono $47$ dimezzamenti, perché $2^{47}\approx1{,}4\cdot10^{14}$.
La funzione \texttt{brentq} (metodo di Brent) combina la bisezione con interpolazioni
più veloci e arriva allo stesso risultato in pochissime iterazioni, conservando la
garanzia di convergenza. Entrambe, come il teorema, \emph{pretendono} che $f$ cambi
segno agli estremi: senza questa ipotesi non garantiscono nulla e si rifiutano di
partire.

\subsection{Il ruolo dell'assioma di completezza}

Chiudiamo il paragrafo con un'osservazione sull'assioma di completezza
(\cref{ass:completezza}). Abbiamo utilizzato tale assioma nella dimostrazione del teorema
dell'esistenza degli zeri, in particolare nell'affermazione che la successione $a_{n}$,
essendo monotòna e limitata, risulta convergente. Ciò è essenziale; infatti,
nell'ambito dei numeri razionali $\Q$, dove non è verificato l'assioma di completezza,
non vale nemmeno il teorema dell'esistenza degli zeri. Ad esempio, la funzione
$f(x)=x^{2}-2$, considerata sull'insieme dei razionali $\{x\in\Q : 0\le x\le2\}$, è
continua e verifica $f(0)<0$, $f(2)>0$, ma non si annulla mai, perché $\sqrt2$ non è
razionale (\cref{teo:radice2}). Le successioni $a_{n}$, $b_{n}$ della bisezione sarebbero
ancora ben definite e razionali, si stringerebbero attorno al «buco» lasciato da
$\sqrt2$, ma non avrebbero un limite in $\Q$. Allo stesso modo, nella seconda
dimostrazione, l'insieme $E=\{x\in\Q : 0\le x\le2,\ x^{2}<2\}$ è limitato, ma non ha
estremo superiore in $\Q$ (\cref{teo:Q-non-completo}). L'assioma di completezza è essenziale anche
negli altri teoremi di esistenza: nel teorema di Weierstrass, o, come già detto, nel
teorema sull'esistenza del limite per le successioni monotòne.

%% ----------------------------------------------------------------------------
%%  §4.9  Appendice: dimostrazione del teorema di Weierstrass
%% ----------------------------------------------------------------------------
\section{Appendice: dimostrazione del teorema di Weierstrass}
\label{sec:dim-weierstrass}

Dimostriamo ora il teorema enunciato nel \cref{sec:teoremi-continue}: \emph{se $f$ è una
funzione continua in un intervallo chiuso e limitato $[a,b]$, allora $f$ assume minimo e
massimo in $[a,b]$}, cioè esistono $x_{1},x_{2}\in[a,b]$ per cui vale la
\eqref{eq:weierstrass}.

\begin{strategia}[Piano d'attacco per il \cref{teo:weierstrass}]
Consideriamo l'estremo superiore $M$ dei valori di $f$, che a priori potrebbe anche
essere $+\infty$. Per le proprietà dell'estremo superiore esiste una successione di punti
$x_{n}$ in cui i valori $f(x_{n})$ si avvicinano a $M$ (una successione
\emph{massimizzante}). I punti $x_{n}$ potrebbero non convergere, ma stanno tutti in
$[a,b]$, che è limitato: per il teorema di Bolzano--Weierstrass se ne può estrarre una
successione convergente a un punto $x_{0}$, che sta ancora in $[a,b]$ perché
l'intervallo è chiuso. Per la continuità, i valori di $f$ lungo l'estratta tendono a
$f(x_{0})$; ma tendono anche a $M$. Dunque $M=f(x_{0})$: l'estremo superiore è finito ed
è un valore assunto, cioè un massimo. Ogni ipotesi interviene in un punto preciso: la
limitatezza dell'intervallo per applicare Bolzano--Weierstrass, la chiusura per non
perdere il punto limite, la continuità per passare al limite dentro $f$.
\end{strategia}

\begin{dimostrazione}[Dimostrazione del \cref{teo:weierstrass}]
Posto $M=\sup\{f(x) : x\in[a,b]\}$ (con la convenzione $M=+\infty$ se l'insieme dei
valori non è limitato superiormente, \eqref{eq:sup-infinito}), verifichiamo che esiste
una successione $x_{n}$ di punti di $[a,b]$ tale che
\begin{equation}\label{eq:weier-massimizzante}
  \lim_{n\to+\infty}f(x_{n})=M .
\end{equation}
Infatti, se $M=+\infty$, nessun numero $n\in\N$ è un maggiorante dei valori di $f$, e
quindi per ogni $n\in\N$ esiste $x_{n}\in[a,b]$ tale che $f(x_{n})>n$; perciò
$f(x_{n})\to M=+\infty$ (\cref{teo:confronto-infinito}). Se invece risulta $M<+\infty$,
per la caratterizzazione dell'estremo superiore (\cref{prop:caratt-sup}) con
$\epsilon=\frac1n$, per ogni $n\in\N$ esiste $x_{n}$ in $[a,b]$ tale che
\begin{equation}\label{eq:weier-stima}
  M-\frac1n<f(x_{n})\le M,
\end{equation}
e perciò $f(x_{n})\to M$ per il teorema dei carabinieri.

La successione $x_{n}$ è limitata, perché $a\le x_{n}\le b$ per ogni $n$. Per il teorema
di Bolzano--Weierstrass (\cref{teo:bolzano-weierstrass}) esistono un'estratta $x_{n_{k}}$
da $x_{n}$ e un punto $x_{0}$ tali che
\begin{equation}\label{eq:weier-estratta}
  x_{n_{k}}\to x_{0} \qquad\text{per } k\to+\infty ;
\end{equation}
inoltre $x_{0}\in[a,b]$, per il \cref{cor:confronto-limiti} applicato alle disuguaglianze
$a\le x_{n_{k}}\le b$. Poiché $f$ è continua in $x_{0}$, ne segue
(\cref{prop:caratt-continuita}(b), che ammette anche eventuali termini $x_{n_{k}}=x_{0}$)
\begin{equation}\label{eq:weier-continuita}
  f(x_{n_{k}})\to f(x_{0}) .
\end{equation}
D'altra parte $f(x_{n_{k}})$ è un'estratta di $f(x_{n})$, e quindi, per la
\eqref{eq:weier-massimizzante} e la \cref{prop:estratta-limite}, tende a $M$. Per
l'unicità del limite
\begin{equation}\label{eq:weier-uguaglianza}
  M=\lim_{n\to+\infty}f(x_{n})=\lim_{k\to+\infty}f(x_{n_{k}})=f(x_{0}) .
\end{equation}
Abbiamo così dimostrato che
\begin{equation}\label{eq:weier-tesi}
  f(x_{0})=M=\sup\{f(x) : x\in[a,b]\};
\end{equation}
ciò implica allo stesso tempo che $M<+\infty$ (perché $f(x_{0})$ è un numero reale) e
che l'estremo superiore è, in effetti, un massimo: posto $x_{2}=x_{0}$, si ha
$f(x)\le f(x_{2})$ per ogni $x\in[a,b]$.

Analogamente si ragiona per determinare un punto di minimo, partendo dall'estremo
inferiore di $f$ in $[a,b]$; oppure si applica quanto appena dimostrato alla funzione
continua $-f$: un punto di massimo $x_{1}$ di $-f$ verifica $-f(x)\le-f(x_{1})$, cioè
$f(x_{1})\le f(x)$, per ogni $x\in[a,b]$.
\end{dimostrazione}

\begin{corollario}\label{cor:continua-limitata}
Una funzione continua in un intervallo chiuso e limitato è limitata.
\end{corollario}

\begin{schema}
\Ipotesi{$f$ continua in $[a,b]$.}
\Tesi{$\esiste L>0$: \ $\abs{f(x)}\le L$ \ $\perogni x\in[a,b]$.}
\end{schema}

\begin{strategia}
Il teorema di Weierstrass dà molto di più della limitatezza: un minimo e un massimo
\emph{reali}. Basta allora racchiudere l'intervallo $[m,M]$ in un intervallo simmetrico
$[-L,L]$.
\end{strategia}

\begin{dimostrazione}
Per il teorema di Weierstrass $m\le f(x)\le M$ per ogni $x\in[a,b]$, con $m$, $M$ numeri
reali; quindi $\abs{f(x)}\le L=\max\{\abs m,\abs M,1\}$ per ogni $x\in[a,b]$
(\cref{prop:limitato-modulo}).
\end{dimostrazione}

Rileggiamo alla luce della dimostrazione gli esempi del \cref{sec:teoremi-continue}
(\cref{fig:controesempi-weierstrass}). Per $\frac1x$ su $(0,1]$ e per $g$ su $[-1,1]$
l'estremo da raggiungere (rispettivamente $+\infty$ e $0$) è avvicinato lungo una
successione che tende a $0$: nel primo caso il punto limite $0$ esce dall'intervallo, che
non è chiuso; nel secondo il punto limite c'è, ma la discontinuità di $g$ impedisce di
passare al limite dentro la funzione, e $g(0)=1$ non ha nulla a che fare con $0$. Per
$\frac1x$ su $[1,+\infty)$ la successione minimizzante $x_{n}=n$ diverge, e nessuna sua
estratta converge: è la limitatezza dell'intervallo che viene a mancare.

\begin{notastorica}[Bolzano, Weierstrass e la compattezza]
Karl Weierstrass dimostrò il teorema che porta il suo nome nelle lezioni berlinesi dei
primi anni Sessanta dell'Ottocento, con un argomento basato proprio su quello che oggi chiamiamo teorema di
Bolzano--Weierstrass. Ancora una volta, però, Bolzano lo aveva preceduto: nella sua
\emph{Functionenlehre}, scritta intorno al 1830, si trova una dimostrazione del fatto che
una funzione continua su un intervallo chiuso e limitato è limitata e assume il massimo;
il manoscritto rimase inedito fino al 1930. La proprietà dell'intervallo $[a,b]$ usata
nella dimostrazione, che \emph{ogni successione di suoi punti ammette un'estratta
convergente a un punto dell'intervallo}, fu isolata all'inizio del Novecento da Maurice
Fréchet, che tra il 1904 e il 1906 introdusse per gli insiemi con questa proprietà il
nome di \emph{compatti} (nella sua terminologia originaria il punto limite non doveva
necessariamente appartenere all'insieme). Con questo linguaggio il teorema di Weierstrass si enuncia in grande
generalità: una funzione continua su un insieme compatto ha massimo e minimo.
\end{notastorica}

%% ----------------------------------------------------------------------------
%%  §4.10  Appendice: funzioni monotòne, funzioni inverse, funzioni elementari
%% ----------------------------------------------------------------------------
\section{Appendice: continuità delle funzioni monotòne e delle funzioni inverse}
\label{sec:monotone-inverse}

Con lo stesso metodo utilizzato per la dimostrazione del teorema sulle successioni
monotòne (\cref{teo:monotone}) si prova il seguente risultato, che riguarda i limiti
unilaterali di una funzione monotòna.

\begin{teorema}[Limiti delle funzioni monotòne]\label{teo:limite-monotone-funzioni}\index{funzione!monotòna!limiti}
Sia $f$ monotòna in $[a,b]$. Allora esistono finiti i limiti
\begin{align}
  &\lim_{x\to a^{+}}f(x), \qquad \lim_{x\to b^{-}}f(x); \label{eq:monotone-estremi}\\
  &\lim_{x\to x_{0}^{-}}f(x), \qquad \lim_{x\to x_{0}^{+}}f(x) \qquad\perogni x_{0}\in(a,b).
  \label{eq:monotone-interni}
\end{align}
Se $f$ è crescente, si ha inoltre (per $x_{0}\in(a,b]$ nella prima uguaglianza, per
$x_{0}\in[a,b)$ nella seconda)
\begin{equation}\label{eq:monotone-sup-inf}
\begin{gathered}
  \lim_{x\to x_{0}^{-}}f(x)=\sup\{f(x) : a\le x<x_{0}\}, \\
  \lim_{x\to x_{0}^{+}}f(x)=\inf\{f(x) : x_{0}<x\le b\} .
\end{gathered}
\end{equation}
\end{teorema}

\begin{schema}
\Ipotesi{$f$ monotòna (ad esempio crescente) in $[a,b]$.}
\Tesi{esistono finiti i limiti unilaterali in ogni punto, dati dalle \eqref{eq:monotone-sup-inf}.}
\end{schema}

\begin{strategia}
È la versione «continua» del teorema sulle successioni monotòne: là il limite di una
successione crescente era l'estremo superiore dei suoi termini, qui il limite sinistro
in $x_{0}$ di una funzione crescente è l'estremo superiore dei valori assunti a sinistra
di $x_{0}$. La dimostrazione è la stessa: se un valore $f(x_{1})$ supera $\ell-\epsilon$,
la monotonia costringe tutti i valori $f(x)$ con $x_{1}<x<x_{0}$ a fare altrettanto, e
l'intervallo $(x_{1},x_{0})$ fornisce il $\delta$ cercato.
\end{strategia}

\begin{dimostrazione}
Consideriamo il caso di una funzione $f$ crescente in $[a,b]$ (se $f$ è decrescente, si
applica il ragionamento a $-f$). Osserviamo subito che $f$ è limitata in $[a,b]$:
\begin{equation}\label{eq:monotone-limitata}
  f(a)\le f(x)\le f(b) \qquad\perogni x\in[a,b];
\end{equation}
cioè $f(a)$ è il minimo di $f$ nell'intervallo $[a,b]$, mentre $f(b)$ è il massimo.
Fissato $x_{0}\in(a,b]$, poniamo
\begin{equation}\label{eq:monotone-ell}
  \ell=\sup\{f(x) : x\in[a,x_{0})\}.
\end{equation}
Per la \eqref{eq:monotone-limitata} l'estremo superiore $\ell$ è finito. Per le proprietà
dell'estremo superiore (\cref{prop:caratt-sup}), per ogni $\epsilon>0$ esiste
$x_{1}\in[a,x_{0})$ tale che
\begin{equation}\label{eq:monotone-x1}
  \ell-\epsilon<f(x_{1}).
\end{equation}
Per $x_{1}<x<x_{0}$ risulta $f(x)\ge f(x_{1})$, perché $f$ è crescente, e $f(x)\le\ell$,
perché $\ell$ è un maggiorante; dunque
\begin{equation}\label{eq:monotone-catena}
  \ell-\epsilon<f(x_{1})\le f(x)\le\ell<\ell+\epsilon \qquad\perogni x\in(x_{1},x_{0}),
\end{equation}
cioè, posto $\delta=x_{0}-x_{1}>0$, si ha $\abs{f(x)-\ell}<\epsilon$ per ogni $x$ con
$x_{0}-\delta<x<x_{0}$. Per la \eqref{eq:def-limite-sx} ne segue
\begin{equation}\label{eq:monotone-limite}
  \lim_{x\to x_{0}^{-}}f(x)=\ell .
\end{equation}
Si procede in modo analogo per il limite per $x\to x_{0}^{+}$, con $x_{0}\in[a,b)$,
usando l'estremo inferiore $\inf\{f(x) : x_{0}<x\le b\}$.
\end{dimostrazione}

Osserviamo che, se $f$ è crescente in $[a,b]$, i limiti \eqref{eq:monotone-estremi},
\eqref{eq:monotone-interni} si possono ordinare nel modo seguente:
\begin{equation}\label{eq:monotone-ordine}
\begin{split}
  f(a)&\le\lim_{x\to a^{+}}f(x)\le\lim_{x\to x_{0}^{-}}f(x)\le f(x_{0})\\
  &\le\lim_{x\to x_{0}^{+}}f(x)\le\lim_{x\to b^{-}}f(x)\le f(b)
\end{split}
\end{equation}
per ogni $x_{0}\in(a,b)$. Ad esempio, la disuguaglianza centrale
$\lim_{x\to x_{0}^{-}}f(x)\le f(x_{0})$ vale perché $f(x_{0})$ è un maggiorante di
$\{f(x) : a\le x<x_{0}\}$, e il limite sinistro ne è l'estremo superiore; le altre si
provano allo stesso modo, o con il confronto dei limiti (\cref{teo:confronto-funzioni}(ii)).

\begin{corollario}\label{cor:monotone-salti}
Le eventuali discontinuità di una funzione monotòna in $[a,b]$ nei punti interni
$x_{0}\in(a,b)$ sono tutte di prima specie (salti). Negli estremi, se $f$ è crescente,
una discontinuità significa $f(a)<\lim_{x\to a^{+}}f(x)$ oppure
$\lim_{x\to b^{-}}f(x)<f(b)$.
\end{corollario}

\begin{schema}
\Ipotesi{$f$ monotòna (ad esempio crescente) in $[a,b]$; \ $f$ discontinua in $x_{0}$.}
\Tesi{se $x_{0}\in(a,b)$, la discontinuità è di prima specie; \ se $x_{0}=a$ (o $b$),
vale la disuguaglianza stretta indicata.}
\end{schema}

\begin{strategia}
Si scorre la classificazione della \cref{def:discontinuita}: il caso (c) è escluso perché
i limiti unilaterali esistono finiti, il caso (a) perché la catena di disuguaglianze
\eqref{eq:monotone-ordine} costringe il valore $f(x_{0})$ a stare tra il limite sinistro
e il limite destro. Resta solo il salto.
\end{strategia}

\begin{dimostrazione}
Per il \cref{teo:limite-monotone-funzioni} i limiti destro e sinistro esistono sempre
finiti, quindi il caso (c) della \cref{def:discontinuita} è escluso; per la
\eqref{eq:monotone-ordine}, se i due limiti coincidono coincidono anche con $f(x_{0})$,
e quindi è escluso anche il caso (a). Negli estremi esiste un solo limite unilaterale,
finito, che per la \eqref{eq:monotone-ordine} è maggiore o uguale a $f(a)$ (rispettivamente
minore o uguale a $f(b)$): la discontinuità equivale alla disuguaglianza stretta.
\end{dimostrazione}

\subsection{Un criterio di continuità per le funzioni monotòne}

Una funzione monotòna può dunque essere discontinua soltanto «saltando», e un salto
lascia necessariamente fuori dall'immagine un intervallo di valori (\cref{fig:monotona-salto}).
Questa osservazione porta a un criterio molto semplice.

\begin{figure}[htb]
\centering
\begin{tikzpicture}
\begin{axis}[width=0.85\textwidth, height=5.4cm, xmin=0, xmax=6.2, ymin=0, ymax=3.6,
  axis lines=left, xtick={0.6,3.2,5.6}, xticklabels={$a$,$x_{0}$,$b$},
  ytick={0.5,1.45,2.35,3.2}, yticklabels={$f(a)$,$\ell_{1}$,$\ell_{2}$,$f(b)$},
  xlabel={$x$}, ylabel={$y$}]
  \fill[rosso!12] (axis cs:0,1.45) rectangle (axis cs:6.2,2.35);
  \addplot[curva, domain=0.6:3.2, samples=60] {0.5+0.95*((x-0.6)/2.6)^0.8};
  \addplot[curva, domain=3.2:5.6, samples=60] {2.35+0.85*((x-3.2)/2.4)^1.3};
  \draw[blu, thick, fill=white] (axis cs:3.2,1.45) circle (2pt);
  \fill[blu] (axis cs:3.2,2.35) circle (2pt);
  \addplot[tratteggio] coordinates {(0,1.45) (3.2,1.45)};
  \addplot[tratteggio] coordinates {(0,2.35) (3.2,2.35) (3.2,0)};
  \addplot[tratteggio] coordinates {(0,3.2) (5.6,3.2) (5.6,0)};
  \addplot[tratteggio] coordinates {(0.6,0) (0.6,0.5)};
  \node[font=\scriptsize, rosso, anchor=west] at (axis cs:3.6,1.9) {valori non assunti};
\end{axis}
\end{tikzpicture}
\caption{Una funzione crescente con un salto in $x_{0}$: limite sinistro $\ell_{1}$,
limite destro $\ell_{2}=f(x_{0})$. I valori dell'intervallo $(\ell_{1},\ell_{2})$, che
pure stanno tra $f(a)$ e $f(b)$, non sono assunti: l'immagine di $f$ non è un intervallo.}
\label{fig:monotona-salto}
\end{figure}

\begin{teorema}[Criterio di continuità per le funzioni monotòne]\label{teo:criterio-monotone}\index{funzione!monotòna!continuità}
Sia $f$ una funzione monotòna nell'intervallo chiuso e limitato $[a,b]$. Allora $f$ è
continua in $[a,b]$ se e solo se l'immagine di $f$ è tutto l'intervallo di estremi
$f(a)$, $f(b)$.
\end{teorema}

\begin{schema}
\Ipotesi{$f$ monotòna (ad esempio crescente) in $[a,b]$.}
\Tesi{$f$ continua in $[a,b]$ \ $\sse$ \ $f([a,b])=[f(a),f(b)]$.}
\end{schema}

\begin{strategia}
Una implicazione è il teorema dei valori intermedi e non richiede la monotonia. Per
l'altra procediamo per assurdo: se $f$ non fosse continua in un punto, per il
\cref{cor:monotone-salti} avrebbe lì un salto, e i valori del salto, pur compresi tra
$f(a)$ e $f(b)$, non sarebbero assunti; ciò contraddice l'ipotesi sull'immagine.
\end{strategia}

\begin{dimostrazione}
Supponiamo $f$ crescente (altrimenti si considera $-f$). Se $f$ è continua in $[a,b]$,
allora, indipendentemente dalla monotonia, assume tutti i valori compresi tra $f(a)$ e
$f(b)$ (\cref{teo:valori-intermedi-1}); e per la \eqref{eq:monotone-limitata} non ne
assume altri. Quindi $f([a,b])=[f(a),f(b)]$.

Viceversa, sia $f([a,b])=[f(a),f(b)]$ e supponiamo per assurdo che $f$ non sia continua
in un punto $x_{0}\in(a,b)$. Posto
\begin{equation}\label{eq:criterio-salto}
  \ell_{1}=\lim_{x\to x_{0}^{-}}f(x), \qquad \ell_{2}=\lim_{x\to x_{0}^{+}}f(x),
\end{equation}
per la \eqref{eq:monotone-ordine} si ha $\ell_{1}\le f(x_{0})\le\ell_{2}$, e almeno una
delle due disuguaglianze è stretta (se fossero entrambe uguaglianze, $f$ sarebbe
continua in $x_{0}$ per la \cref{prop:dx-sx}). Supponiamo ad esempio $\ell_{1}<f(x_{0})$
(l'altro caso è analogo). Allora $f$ non assume alcun valore nell'intervallo
$(\ell_{1},f(x_{0}))$: per $x<x_{0}$ si ha $f(x)\le\ell_{1}$, per la
\eqref{eq:monotone-sup-inf}, e per $x\ge x_{0}$ si ha $f(x)\ge f(x_{0})$. D'altra parte
$f(a)\le\ell_{1}<f(x_{0})\le f(b)$, quindi $(\ell_{1},f(x_{0}))\subseteq[f(a),f(b)]$:
questi valori dovrebbero essere assunti, e si ha una contraddizione. Si procede in modo
analogo se $x_{0}=a$ oppure $x_{0}=b$: ad esempio, se $f$ non è continua per $x=a$,
allora $f$ non assume alcun valore nell'intervallo $(f(a),\ell)$, essendo
\begin{equation}\label{eq:criterio-estremo}
  f(a)<\ell=\lim_{x\to a^{+}}f(x)=\inf\{f(x) : a<x\le b\}. \qedhere
\end{equation}
\end{dimostrazione}

\subsection{Continuità della funzione inversa}

\begin{teorema}[Continuità delle funzioni inverse]\label{teo:continuita-inversa}\index{funzione!inversa!continuità}
Sia $f$ una funzione strettamente monotòna in $[a,b]$. Se $f$ è continua, anche la
funzione inversa $f^{-1}$ è continua.
\end{teorema}

\begin{schema}
\Ipotesi{$f$ continua e strettamente crescente (o decrescente) in $[a,b]$.}
\Tesi{$f^{-1}:[f(a),f(b)]\to[a,b]$ (rispettivamente $[f(b),f(a)]\to[a,b]$) è continua.}
\end{schema}

\begin{strategia}
Non serve alcuna stima: si applica il criterio appena dimostrato alla funzione $f^{-1}$.
Basta osservare che anche $f^{-1}$ è monotòna, e che la sua immagine è tutto
l'intervallo $[a,b]$.
\end{strategia}

\begin{dimostrazione}
Osserviamo che la stretta monotonia di $f$ su $[a,b]$ implica la sua invertibilità (si
veda il criterio di invertibilità, \cref{teo:invertibilita}). Supponiamo, per fissare le
notazioni, che $f$ sia strettamente crescente in $[a,b]$; allora
\begin{equation}\label{eq:inversa-domini}
  f:[a,b]\to[f(a),f(b)], \qquad f^{-1}:[f(a),f(b)]\to[a,b]
\end{equation}
sono biiettive. La funzione $f^{-1}$ è strettamente crescente: se $y_{1}<y_{2}$ e fosse
$f^{-1}(y_{1})\ge f^{-1}(y_{2})$, applicando la funzione crescente $f$ si otterrebbe
$y_{1}\ge y_{2}$. Inoltre $f^{-1}$ assume tutti i valori dell'intervallo $[a,b]$, che è
l'intervallo di estremi $f^{-1}(f(a))=a$ e $f^{-1}(f(b))=b$. Per il criterio precedente
(di continuità delle funzioni monotòne), $f^{-1}$ è continua.
\end{dimostrazione}

\begin{strategia}[Un'altra via]
Si può fare a meno del criterio per le funzioni monotòne e ragionare per successioni,
come nella dimostrazione del teorema di Weierstrass. Se $f^{-1}$ non fosse continua in
un punto $y_{0}$, si troverebbero valori $y_{n}\to y_{0}$ le cui controimmagini $x_{n}$
restano lontane da $x_{0}=f^{-1}(y_{0})$; per il teorema di Bolzano--Weierstrass le
$x_{n}$ hanno un'estratta convergente a un punto $x'\neq x_{0}$, e la continuità di $f$
dà $f(x')=y_{0}=f(x_{0})$, contro l'iniettività. Della stretta monotonia si usa soltanto
l'iniettività; sono essenziali, invece, la continuità di $f$ e il fatto che $[a,b]$ sia
chiuso e limitato.
\end{strategia}

\begin{dimostrazione}[Seconda dimostrazione]
Sia $J=f([a,b])$ il dominio di $f^{-1}$, che è un intervallo chiuso e limitato
(\cref{cor:immagine-intervallo}); siano $y_{0}\in J$ e $x_{0}=f^{-1}(y_{0})$. Supponiamo per
assurdo che $f^{-1}$ non sia continua in $y_{0}$. Negando la caratterizzazione (c) della
\cref{prop:caratt-continuita} e scegliendo $\delta=\frac1n$, come nella dimostrazione del
\cref{teo:eps-delta}, si trova $\epsilon_{0}>0$ tale che per ogni $n\in\N$ esiste
$y_{n}\in J$ con
\begin{equation}\label{eq:inversa-succ-cattiva}
  \abs{y_{n}-y_{0}}<\frac1n, \qquad \abs{f^{-1}(y_{n})-x_{0}}\ge\epsilon_{0} .
\end{equation}
Posto $x_{n}=f^{-1}(y_{n})$, si ha dunque $y_{n}\to y_{0}$, $x_{n}\in[a,b]$,
$f(x_{n})=y_{n}$ e $\abs{x_{n}-x_{0}}\ge\epsilon_{0}$ per ogni $n$. La successione
$x_{n}$ è limitata: per il teorema di Bolzano--Weierstrass (\cref{teo:bolzano-weierstrass})
ha un'estratta $x_{n_{k}}$ convergente a un punto $x'$, che appartiene ad $[a,b]$ per il
\cref{cor:confronto-limiti}. Non può essere $x'=x_{0}$: altrimenti si avrebbe
definitivamente $\abs{x_{n_{k}}-x_{0}}<\epsilon_{0}$, contro la
\eqref{eq:inversa-succ-cattiva}. Per la continuità di $f$ in $x'$
(\cref{prop:caratt-continuita}(b)) si ha $f(x_{n_{k}})\to f(x')$; d'altra parte
$f(x_{n_{k}})=y_{n_{k}}$ è un'estratta di $y_{n}$, e quindi tende a $y_{0}$
(\cref{prop:estratta-limite}). Per l'unicità del limite $f(x')=y_{0}=f(x_{0})$, con
$x'\neq x_{0}$: ciò contraddice l'iniettività di $f$.
\end{dimostrazione}

\begin{osservazione}[Che cosa insegna la seconda via]\label{oss:inversa-compattezza}
La seconda dimostrazione non usa la monotonia di $f$, ma solo la sua iniettività: prova
dunque che \emph{l'inversa di una funzione continua e iniettiva su $[a,b]$ è continua}.
Usa invece in modo essenziale la proprietà delle estratte convergenti, che nella nota
storica del \cref{sec:dim-weierstrass} abbiamo chiamato compattezza: vale quindi, con
la stessa dimostrazione, per le funzioni continue e iniettive su un qualunque insieme $K$
in cui ogni successione ha un'estratta convergente a un punto di $K$, anche se $K$ non è
un intervallo. La prima dimostrazione sfrutta invece l'ordine dei numeri reali.
Nell'approfondimento del \S\ref{subsec:continue-iniettive} vedremo che le due ipotesi
non sono poi così diverse: su un intervallo, una funzione continua e iniettiva è
necessariamente strettamente monotòna.
\end{osservazione}

\begin{figure}[htb]
\centering
\begin{tikzpicture}
\begin{axis}[width=0.7\textwidth, height=0.7\textwidth, xmin=-3, xmax=3.2, ymin=-3, ymax=3.2,
  axis lines=middle, xtick={-2,-1,1,2,3}, ytick={-2,-1,1,2,3}, xlabel={$x$}, ylabel={$y$},
  tick label style={font=\scriptsize}, axis equal image]
  \addplot[tratteggio, domain=-3:3.2] {x};
  \addplot[curva, domain=-3:1.16, samples=80] {exp(x)} node[pos=0.93, left, font=\scriptsize] {$e^{x}$};
  \addplot[curva2, domain=0.05:3.2, samples=120] {ln(x)} node[pos=0.9, below right, font=\scriptsize] {$\log x$};
\end{axis}
\end{tikzpicture}
\caption{La funzione esponenziale e la sua inversa, il logaritmo: i grafici sono
simmetrici rispetto alla bisettrice $y=x$. La continuità dell'una si trasmette all'altra
(\cref{teo:continuita-inversa}).}
\label{fig:inversa-continua}
\end{figure}

\begin{idea}[Invertire una curva di taratura]
Molti strumenti di bordo misurano una grandezza per ricavarne un'altra attraverso una
relazione monotòna. L'anemometro misura la differenza di pressione
$\Delta p=\frac12\rho V^{2}$ tra presa totale e presa statica, e ne ricava la velocità
$V=\sqrt{2\Delta p/\rho}$: la funzione $V\mapsto\Delta p$ è continua e strettamente
crescente per $V\ge0$, e la sua inversa è continua per il \cref{teo:continuita-inversa}.
Il significato pratico è rassicurante: un piccolo errore nella misura di $\Delta p$
produce un piccolo errore su $V$. La continuità, però, non dice \emph{quanto} piccolo:
a basse velocità la curva $\Delta p(V)$ è quasi piatta, e un errore fissato su $\Delta p$
si traduce in un errore su $V$ molto più grande che ad alta velocità. È il motivo per cui
gli anemometri sono poco precisi alle basse velocità; per quantificare questa
sensibilità servirà la derivata.
\end{idea}

\subsection{Continuità delle funzioni elementari}

Siamo finalmente in grado di dimostrare la continuità di esponenziali, logaritmi e
potenze, annunciata nel \cref{sec:continue}, e di chiudere un conto rimasto aperto nei
\cref{cap:complementi,cap:successioni}: le regole di calcolo delle potenze con
esponente reale.

\begin{teorema}[Continuità di esponenziali, logaritmi e potenze]\label{teo:continuita-esponenziale}\index{funzione!esponenziale!continuità}\index{funzione!logaritmo!continuità}
\begin{enumerate}[label=\textup{(\roman*)}]
  \item Per ogni $a>0$ la funzione $x\mapsto a^{x}$ è continua su $\R$.
  \item Per ogni $a>0$, $a\neq1$, la funzione $y\mapsto\log_{a}y$ è continua su $(0,+\infty)$.
  \item Per ogni $b\in\R$ la funzione $x\mapsto x^{b}$ è continua su $(0,+\infty)$.
\end{enumerate}
\end{teorema}

\begin{schema}
\Ipotesi{la \cref{def:potenza-reale} di $a^{x}$ e il \cref{teo:esponenziale}
(monotonia stretta e suriettività su $(0,+\infty)$).}
\Tesi{continuità di $a^{x}$ su $\R$, di $\log_{a}y$ su $(0,+\infty)$, di $x^{b}$ su
$(0,+\infty)$.}
\end{schema}

\begin{strategia}
Non c'è bisogno di stime con $\epsilon$ e $\delta$: tutto il lavoro è già stato fatto
nel \cref{teo:esponenziale}, che fornisce una funzione strettamente crescente la cui
immagine è un intervallo. Il criterio di continuità per le funzioni monotòne trasforma
questa informazione in continuità, su ogni intervallo chiuso e limitato; la località del
limite la estende a tutto $\R$. Il logaritmo è l'inversa dell'esponenziale, e per esso
si usa il \cref{teo:continuita-inversa}. Per le potenze, infine, si scrive
$x^{b}=e^{b\log x}$: ma per giustificare questa uguaglianza servono le regole delle
potenze con esponente reale, che dimostreremo subito dopo usando la (i).
\end{strategia}

\begin{dimostrazione}[Dimostrazione di (i) e (ii)]
(i) Sia dapprima $a>1$, e fissiamo $x_{0}\in\R$. Consideriamo l'intervallo
$[\alpha,\beta]=[x_{0}-1,x_{0}+1]$. La funzione $f(x)=a^{x}$ è strettamente crescente
(\cref{teo:esponenziale}(ii)), e l'immagine di $[\alpha,\beta]$ è l'intervallo
$[a^{\alpha},a^{\beta}]$: infatti ogni valore $a^{x}$ con $\alpha\le x\le\beta$ sta in
$[a^{\alpha},a^{\beta}]$ per la monotonia, e viceversa, se $c\in[a^{\alpha},a^{\beta}]$,
per il \cref{teo:esponenziale}(iii) esiste $x\in\R$ con $a^{x}=c$, e deve essere
$\alpha\le x\le\beta$, perché $x<\alpha$ darebbe $a^{x}<a^{\alpha}\le c$ e $x>\beta$
darebbe $a^{x}>a^{\beta}\ge c$. Per il \cref{teo:criterio-monotone} la restrizione di
$f$ a $[\alpha,\beta]$ è continua, in particolare nel punto interno $x_{0}$; poiché
$[\alpha,\beta]$ contiene l'intorno $(x_{0}-1,x_{0}+1)$ di $x_{0}$, per la località del
limite (\cref{lem:localita}(ii)) $f$ è continua in $x_{0}$. Se $a=1$ la funzione è
costante. Se $0<a<1$, si ha $a^{x}=(1/a)^{-x}$ per definizione, e la funzione è
composta delle funzioni continue $x\mapsto-x$ e $t\mapsto(1/a)^{t}$, con $1/a>1$; è
dunque continua per il \cref{teo:composta-continua}.

(ii) Sia $a>1$ e sia $y_{0}>0$; posto $x_{0}=\log_{a}y_{0}$ e
$[\alpha,\beta]=[x_{0}-1,x_{0}+1]$, la funzione $a^{x}$ è continua e strettamente
crescente su $[\alpha,\beta]$, con immagine $[a^{\alpha},a^{\beta}]$. Per il
\cref{teo:continuita-inversa} la sua inversa, che è la restrizione di $\log_{a}$ a
$[a^{\alpha},a^{\beta}]$, è continua; poiché $a^{\alpha}<y_{0}<a^{\beta}$, l'intervallo
$[a^{\alpha},a^{\beta}]$ contiene un intorno di $y_{0}$, e per la località del limite
$\log_{a}$ è continua in $y_{0}$. Se $0<a<1$, dalla definizione $a^{x}=(1/a)^{-x}$ segue che
$a^{x}=y$ se e solo se $(1/a)^{-x}=y$, cioè $-x=\log_{1/a}y$: dunque
$\log_{a}y=-\log_{1/a}y$, che è continua per il caso già dimostrato, essendo $1/a>1$.
(Si noti che non abbiamo usato la formula del cambiamento di base \eqref{eq:log4}, che
con esponenti reali richiede le regole delle potenze, dimostrate soltanto più avanti.)
\end{dimostrazione}

\begin{teorema}[Potenze con esponente reale]\label{teo:potenze-reali}\index{potenza a esponente reale!proprietà}
Per ogni $a,b>0$ e per ogni $x,z,c\in\R$ si ha
\begin{align}
  &a^{x}\,a^{z}=a^{x+z}, \qquad (a^{x})^{z}=a^{xz}, \qquad (ab)^{x}=a^{x}\,b^{x};
  \label{eq:potenze-reali-algebra}\\
  &a<b,\ c>0 \ \implica\ a^{c}<b^{c}; \qquad a<b,\ c<0 \ \implica\ a^{c}>b^{c} .
  \label{eq:potenze-reali-ordine}
\end{align}
Cioè le proprietà \eqref{eq:pot1}, \eqref{eq:pot3}, \eqref{eq:pot4}, note per esponenti
razionali, valgono per esponenti reali qualsiasi.
\end{teorema}

\begin{schema}
\Ipotesi{$a,b>0$; \ $x,z,c\in\R$; \ le proprietà per esponenti razionali (\cref{sec:potenze});
la continuità delle funzioni esponenziali (\cref{teo:continuita-esponenziale}(i)).}
\Tesi{le uguaglianze \eqref{eq:potenze-reali-algebra} e le disuguaglianze
\eqref{eq:potenze-reali-ordine}.}
\end{schema}

\begin{strategia}
Un numero reale è limite di una successione di razionali; per i razionali le regole sono
note; la continuità delle funzioni esponenziali permette di passare al limite nei due
membri di ciascuna uguaglianza, e l'unicità del limite conclude. È un procedimento tipico,
che si chiama \emph{estensione per continuità}: una relazione tra funzioni continue che
vale su un insieme denso vale ovunque. Per $(a^{x})^{z}$ occorre un passo intermedio,
perché la base $a^{x}$ non è più un numero «noto»: si tratta prima il caso di $z$
razionale, usando l'unicità della radice $q$-esima.
\end{strategia}

\begin{dimostrazione}
Ricordiamo che per ogni $x\in\R$ esiste una successione di razionali $q_{n}\to x$: per
la densità di $\Q$ (\cref{teo:densita-Q}) esiste $q_{n}\in\Q$ con $x-\frac1n<q_{n}<x$, e
$q_{n}\to x$ per il teorema dei carabinieri.

\emph{Prima uguaglianza.} Siano $q_{n}\to x$ ed $r_{n}\to z$ successioni di razionali.
Per la \eqref{eq:pot1} con esponenti razionali, $a^{q_{n}}a^{r_{n}}=a^{q_{n}+r_{n}}$ per
ogni $n$. Per la continuità di $t\mapsto a^{t}$ nei punti $x$, $z$ e $x+z$
(\cref{prop:caratt-continuita}(b)), il primo membro tende a $a^{x}a^{z}$ e il secondo,
poiché $q_{n}+r_{n}\to x+z$, tende a $a^{x+z}$. Per l'unicità del limite
$a^{x}a^{z}=a^{x+z}$. In particolare $a^{x}a^{-x}=a^{0}=1$, cioè $a^{-x}=1/a^{x}$.

\emph{Terza uguaglianza.} Allo stesso modo $(ab)^{q_{n}}=a^{q_{n}}b^{q_{n}}$ e, passando al
limite con la continuità delle tre funzioni esponenziali di base $ab$, $a$, $b$, si
ottiene $(ab)^{x}=a^{x}b^{x}$.

\emph{Seconda uguaglianza.} Sia dapprima $z=\frac pq$ razionale, con $q\in\N$ e
$p\in\Z$. Per la prima uguaglianza, applicata $q-1$ volte,
\[
  \bigl(a^{x/q}\bigr)^{q}=\underbrace{a^{x/q}\cdots a^{x/q}}_{q\ \text{volte}}=a^{x};
\]
dunque $a^{x/q}$ è un numero positivo la cui potenza $q$-esima è $a^{x}$, e per l'unicità
della radice $q$-esima (\cref{teo:radice-n}) è $a^{x/q}=(a^{x})^{1/q}$. Se $p>0$ si ha
allora, per le proprietà delle potenze con esponente razionale (che danno
$y^{p/q}=\bigl(y^{1/q}\bigr)^{p}$ per $y>0$) e ancora per la prima uguaglianza,
$(a^{x})^{p/q}=\bigl((a^{x})^{1/q}\bigr)^{p}=\bigl(a^{x/q}\bigr)^{p}=a^{xp/q}$; se $p<0$ si
passa al reciproco, usando $a^{-y}=1/a^{y}$; se $p=0$ entrambi i membri valgono $1$.
Sia ora $z$ reale, e sia $r_{n}\to z$ una successione di razionali. Per quanto appena
visto $(a^{x})^{r_{n}}=a^{xr_{n}}$. Il primo membro tende a $(a^{x})^{z}$, per la
continuità della funzione esponenziale di base $a^{x}>0$; il secondo tende a $a^{xz}$,
perché $xr_{n}\to xz$ e $t\mapsto a^{t}$ è continua. Quindi $(a^{x})^{z}=a^{xz}$.

\emph{Disuguaglianze.} Siano $0<a<b$ e $c>0$. Poiché $\frac ba>1$, la stretta monotonia
dell'esponenziale di base $\frac ba$ (\cref{teo:esponenziale}(ii)) dà
$\bigl(\frac ba\bigr)^{c}>\bigl(\frac ba\bigr)^{0}=1$; per la terza uguaglianza
\[
  b^{c}=\Bigl(\frac ba\cdot a\Bigr)^{c}=\Bigl(\frac ba\Bigr)^{c}a^{c}>a^{c} .
\]
Se $c<0$, si ha $a^{-c}<b^{-c}$ per quanto appena visto, e passando ai reciproci
$a^{c}=\frac1{a^{-c}}>\frac1{b^{-c}}=b^{c}$.
\end{dimostrazione}

\begin{dimostrazione}[Dimostrazione del \cref{teo:continuita-esponenziale}(iii)]
Per $x>0$ si ha $x=e^{\log x}$ \eqref{eq:log-inversa}, e quindi, per la seconda delle
\eqref{eq:potenze-reali-algebra},
\begin{equation}\label{eq:potenza-esponenziale}
  x^{b}=\bigl(e^{\log x}\bigr)^{b}=e^{b\log x} \qquad\perogni x>0 .
\end{equation}
La funzione $x\mapsto b\log x$ è continua su $(0,+\infty)$ per la (ii), e $t\mapsto e^{t}$
è continua su $\R$ per la (i): per il \cref{teo:composta-continua} la funzione
$x\mapsto x^{b}$ è continua su $(0,+\infty)$.
\end{dimostrazione}

\begin{osservazione}
Se $b>0$, ponendo $0^{b}=0$ la funzione $x^{b}$ risulta continua anche in $0$ (da destra):
per $x\to0^{+}$ si ha $b\log x\to-\infty$ (\cref{prop:lim-logaritmo}) ed $e^{t}\to0$ per
$t\to-\infty$ (\cref{prop:lim-esponenziale}), quindi $x^{b}=e^{b\log x}\to0$ per il
\cref{teo:limite-composta} (con $y_{0}=-\infty$). Dalla \eqref{eq:potenza-esponenziale}
segue anche la regola $\log x^{b}=b\log x$ per ogni $b\in\R$, cioè la \eqref{eq:log3}
con esponente reale. Ne segue anche la formula del cambiamento di base \eqref{eq:log4} per
esponenti reali: se $y=a^{x}$, allora $\log y=\log\bigl(e^{\log a}\bigr)^{x}=x\log a$. Le
osservazioni sulle potenze a esponente reale fatte nel \cref{sec:esponenziale-reale} e nel
\cref{sec:numero-e} sono così completamente giustificate; e la dimostrazione non ha usato
alcun risultato che dipendesse da esse.
\end{osservazione}

\subsection{Limiti notevoli}

La continuità di esponenziali e logaritmi, insieme al limite \eqref{eq:lim-1-piu-x},
permette di ricavare alcuni limiti notevoli di uso continuo.

\begin{proposizione}[Limiti notevoli di esponenziali e logaritmi]\label{prop:limiti-notevoli-exp}\index{limite!notevole}
Per ogni $b\in\R$ e ogni $a>0$ si ha
\begin{align}
  &\lim_{x\to0}(1+bx)^{1/x}=e^{b}; \label{eq:ln-e-b}\\
  &\lim_{x\to0}\frac{\log(1+x)}{x}=1; \label{eq:ln-log}\\
  &\lim_{x\to0}\frac{e^{x}-1}{x}=1, \qquad \lim_{x\to0}\frac{a^{x}-1}{x}=\log a; \label{eq:ln-exp}\\
  &\lim_{x\to0}\frac{(1+x)^{b}-1}{x}=b . \label{eq:ln-pot}
\end{align}
\end{proposizione}

\begin{schema}
\Ipotesi{$b\in\R$, \ $a>0$; \ $(1+x)^{1/x}\to e$ per $x\to0$ \eqref{eq:lim-1-piu-x}; \
continuità di $e^{t}$, $\log t$, $t^{b}$ (\cref{teo:continuita-esponenziale}).}
\Tesi{i limiti \eqref{eq:ln-e-b}--\eqref{eq:ln-pot}.}
\end{schema}

\begin{strategia}
Tutti i limiti sono forme indeterminate ($1^{\infty}$ o $\frac00$) e si riconducono, con
le regole delle potenze e qualche cambio di variabile, al limite
$(1+x)^{1/x}\to e$. Per «portare il limite dentro» le funzioni $\log$, $e^{t}$, $t^{b}$ si
usa la loro continuità (\cref{cor:limite-composta-continua}); per i cambi di variabile
si usa il \cref{teo:limite-composta}, verificando ogni volta l'ipotesi $g(x)\neq y_{0}$.
\end{strategia}

\begin{dimostrazione}
\eqref{eq:ln-e-b}: Se $b=0$ la funzione vale costantemente $1=e^{0}$. Se $b\neq0$, per
$0<\abs x<\frac1{\abs b}$ si ha $1+bx>0$ e, per il \cref{teo:potenze-reali},
\[
  (1+bx)^{1/x}=\Bigl[(1+bx)^{\frac1{bx}}\Bigr]^{b}.
\]
Con $g(x)=bx$, che tende a $0$ ed è diversa da $0$ per $x\neq0$, il
\cref{teo:limite-composta} e la \eqref{eq:lim-1-piu-x} danno $(1+bx)^{1/(bx)}\to e$;
per la continuità della funzione $t\mapsto t^{b}$ nel punto $t=e$
(\cref{teo:continuita-esponenziale}(iii) e \cref{cor:limite-composta-continua}) il
secondo membro tende a $e^{b}$.

\eqref{eq:ln-log}: Per $x>-1$, $x\neq0$, si ha
$\frac{\log(1+x)}x=\frac1x\log(1+x)=\log\bigl[(1+x)^{1/x}\bigr]$ (regola $\log y^{c}=c\log y$).
Poiché $(1+x)^{1/x}\to e$ e il logaritmo è continuo nel punto $e$, il limite è
$\log e=1$.

\eqref{eq:ln-exp}: Poniamo $y=g(x)=e^{x}-1$. Per la continuità dell'esponenziale,
$g(x)\to e^{0}-1=0$ per $x\to0$, e $g(x)\neq0$ per $x\neq0$ perché l'esponenziale è
iniettiva; inoltre $x=\log(1+y)$. Dunque
\[
  \frac{e^{x}-1}{x}=\frac{y}{\log(1+y)}=f(g(x)), \qquad\text{dove}\quad
  f(y)=\frac{y}{\log(1+y)},
\]
e $f(y)\to\frac11=1$ per $y\to0$
per la \eqref{eq:ln-log}; il \cref{teo:limite-composta} conclude. Per la seconda, se
$a=1$ entrambi i membri sono nulli; se $a\neq1$, per la \eqref{eq:potenza-esponenziale}
$a^{x}=e^{x\log a}$ e
\[
  \frac{a^{x}-1}{x}=\log a\cdot\frac{e^{x\log a}-1}{x\log a}\ \to\ \log a\cdot1,
\]
per il \cref{teo:limite-composta} con $g(x)=x\log a$, che si annulla solo per $x=0$.

\eqref{eq:ln-pot}: Se $b=0$ il limite è $0=b$. Se $b\neq0$, per $x>-1$, $x\neq0$, si
ha $(1+x)^{b}=e^{b\log(1+x)}$, e la quantità $y=b\log(1+x)$ è diversa da zero e tende a
zero; quindi
\[
  \frac{(1+x)^{b}-1}{x}=\frac{e^{y}-1}{y}\cdot\frac{b\log(1+x)}{x}\ \to\ 1\cdot b=b,
\]
per la \eqref{eq:ln-exp}, il \cref{teo:limite-composta} e la \eqref{eq:ln-log}.
\end{dimostrazione}

\begin{osservazione}
Con la continuità del logaritmo si ritrova anche il limite $\frac{\log n}{n}\to0$,
dimostrato nel \cref{teo:gerarchia}: infatti
$\frac{\log n}{n}=\log n^{1/n}=\log\sqrt[n]n\to\log1=0$, per la \cref{prop:radice-n-n} e la
continuità del logaritmo nel punto $1$.
\end{osservazione}

I limiti \eqref{eq:ln-log} e \eqref{eq:ln-exp} dicono che, per $x$ piccolo,
$\log(1+x)\approx x$ ed $e^{x}-1\approx x$, con un errore relativo che tende a zero. Ma,
come per $1-\cos x$, il calcolo numerico diretto di queste espressioni soffre di
cancellazione: per questo le librerie di calcolo forniscono le funzioni
\texttt{log1p}$(x)=\log(1+x)$ ed \texttt{expm1}$(x)=e^{x}-1$, calcolate senza formare la
differenza.

\begin{python}[I limiti notevoli e le funzioni \texttt{log1p}, \texttt{expm1}]
import math

# log(1+x)/x -> 1 e (e^x - 1)/x -> 1 per x -> 0. Le funzioni
# log1p(x) = log(1+x), expm1(x) = e^x - 1 evitano la cancellazione.
print("    x     log(1+x)/x   log1p(x)/x   (e^x-1)/x    expm1(x)/x")
for k in (1, 3, 5, 8, 11, 14, 16):
    x = 10.0**(-k)
    a, b = math.log(1 + x) / x, math.log1p(x) / x
    c, d = (math.exp(x) - 1) / x, math.expm1(x) / x
    print(f"1e-{k:<3d} {a:12.10f} {b:12.10f} {c:12.10f} {d:12.10f}")
\end{python}
\begin{risultato}
    x     log(1+x)/x   log1p(x)/x   (e^x-1)/x    expm1(x)/x
1e-1   0.9531017980 0.9531017980 1.0517091808 1.0517091808
1e-3   0.9995003331 0.9995003331 1.0005001667 1.0005001667
1e-5   0.9999950000 0.9999950000 1.0000050000 1.0000050000
1e-8   0.9999999889 0.9999999950 0.9999999939 1.0000000050
1e-11  1.0000000827 1.0000000000 1.0000000827 1.0000000000
1e-14  0.9992007222 1.0000000000 0.9992007222 1.0000000000
1e-16  0.0000000000 1.0000000000 0.0000000000 1.0000000000
\end{risultato}

\begin{idea}[L'equazione del razzo per piccole frazioni di propellente]
L'equazione di Tsiolkovsky dà l'incremento di velocità di un razzo che brucia
propellente: $\Delta v=v_{e}\log\frac{m_{0}}{m_{1}}$, dove $v_{e}$ è la velocità efficace di
scarico, $m_{0}$ la massa iniziale e $m_{1}$ quella finale. Posto
$\mu=\frac{m_{0}-m_{1}}{m_{1}}$, si ha $\Delta v=v_{e}\log(1+\mu)$, e per la
\eqref{eq:ln-log}, quando la massa di propellente consumata è piccola rispetto a quella
del veicolo, $\Delta v\approx v_{e}\,\mu$: è la conservazione della quantità di moto «a
un passo», come se il propellente fosse espulso tutto insieme. Con $v_{e}=3000$~m/s e
$\mu=0{,}05$ la formula esatta dà $\Delta v=146{,}4$~m/s, quella approssimata $150$~m/s:
un errore del $2{,}5\%$, che diventa trascurabile per le piccole correzioni d'orbita e
cresce rapidamente per le manovre importanti. La differenza tra le due formule, il
logaritmo al posto del termine lineare $\mu$, è il prezzo da pagare per dover accelerare
anche il propellente non ancora bruciato.
\end{idea}

\subsection{Approfondimento: funzioni continue e iniettive}
\label{subsec:continue-iniettive}

Il criterio di invertibilità (\cref{teo:invertibilita}) dice che una funzione continua e
strettamente monotòna su un intervallo è iniettiva. Vale anche il viceversa: su un
intervallo, una funzione continua può essere iniettiva soltanto se è strettamente
monotòna. L'idea intuitiva è chiara: un grafico che si traccia senza staccare la matita,
e che prima sale e poi scende, ripassa per forza a una stessa altezza.

\begin{proposizione}[Le funzioni continue e iniettive sono monotòne]\label{prop:continua-iniettiva-monotona}
Sia $I$ un intervallo e sia $f:I\to\R$ continua e iniettiva. Allora $f$ è strettamente
monotòna in $I$.
\end{proposizione}

\begin{schema}
\Ipotesi{$I$ intervallo; \ $f$ continua in $I$; \ $f(x_{1})\neq f(x_{2})$ se $x_{1}\neq x_{2}$.}
\Tesi{$f$ è strettamente crescente oppure strettamente decrescente in $I$.}
\end{schema}

\begin{strategia}
Ragioniamo per assurdo. Se $f$ non fosse strettamente monotòna, ci sarebbero una coppia
di punti $a_{1}<a_{2}$ su cui $f$ «sale» e una coppia $x_{1}<x_{2}$ su cui $f$
«scende». Facciamo scorrere con continuità la prima coppia fino alla seconda, con due
punti mobili $u(t)<v(t)$: la differenza $f(v(t))-f(u(t))$ passa da un valore positivo a
uno negativo, e per il teorema degli zeri si annulla. Ma allora $f$ assume lo stesso
valore in due punti distinti.
\end{strategia}

\begin{dimostrazione}
Poiché $f$ è iniettiva, in due punti distinti i valori sono sempre diversi. Dire che $f$
non è strettamente crescente significa allora che esistono $x_{1}<x_{2}$ in $I$ con
$f(x_{1})>f(x_{2})$; dire che non è strettamente decrescente, che esistono $a_{1}<a_{2}$
in $I$ con $f(a_{1})<f(a_{2})$. Supponiamo per assurdo che accadano entrambe le cose, e
poniamo per $t\in[0,1]$
\begin{equation}\label{eq:iniettiva-uv}
  u(t)=(1-t)\,a_{1}+t\,x_{1}, \qquad v(t)=(1-t)\,a_{2}+t\,x_{2} .
\end{equation}
Il punto $u(t)$ è compreso tra $a_{1}$ e $x_{1}$, e quindi sta in $I$, perché un
intervallo contiene tutti i punti compresi tra due suoi punti (\cref{def:intervalli});
lo stesso vale per $v(t)$. Inoltre
\[
  v(t)-u(t)=(1-t)\,(a_{2}-a_{1})+t\,(x_{2}-x_{1})>0 \qquad\perogni t\in[0,1],
\]
perché è la somma di due termini non negativi che non si annullano insieme (il primo si
annulla solo per $t=1$, il secondo solo per $t=0$). La funzione
\[
  \phi(t)=f\bigl(v(t)\bigr)-f\bigl(u(t)\bigr), \qquad t\in[0,1],
\]
è continua, perché $u$ e $v$ sono polinomi di primo grado e $f$ è continua
(\cref{teo:composta-continua,teo:operazioni-continue}); inoltre
\[
  \phi(0)=f(a_{2})-f(a_{1})>0, \qquad \phi(1)=f(x_{2})-f(x_{1})<0 .
\]
Per il teorema degli zeri, applicato a $-\phi$ (\cref{teo:zeri}), esiste $t_{0}\in(0,1)$
con $\phi(t_{0})=0$, cioè $f\bigl(u(t_{0})\bigr)=f\bigl(v(t_{0})\bigr)$, mentre
$u(t_{0})<v(t_{0})$: ciò contraddice l'iniettività di $f$. Dunque $f$ è strettamente
crescente oppure strettamente decrescente.
\end{dimostrazione}

\begin{corollario}\label{cor:continua-invertibile}
Sia $f$ continua nell'intervallo $[a,b]$. Allora $f$ è iniettiva se e solo se è
strettamente monotòna; in tal caso $f([a,b])$ è l'intervallo di estremi $f(a)$, $f(b)$,
e la funzione inversa è continua.
\end{corollario}

\begin{schema}
\Ipotesi{$f$ continua in $[a,b]$.}
\Tesi{$f$ iniettiva $\sse$ $f$ strettamente monotòna; \ in tal caso $f^{-1}$ è continua
sull'intervallo di estremi $f(a)$, $f(b)$.}
\end{schema}

\begin{strategia}
Basta mettere insieme i risultati: la \cref{prop:continua-iniettiva-monotona} dà la
monotonia, il criterio di invertibilità l'immagine, il \cref{teo:continuita-inversa} la
continuità dell'inversa.
\end{strategia}

\begin{dimostrazione}
Una funzione strettamente monotòna è iniettiva, perché punti distinti hanno valori
ordinati in senso stretto; viceversa, se $f$ è iniettiva, è strettamente monotòna per la
\cref{prop:continua-iniettiva-monotona}. In tal caso il \cref{teo:invertibilita} dà
$f([a,b])$, e il \cref{teo:continuita-inversa} la continuità di $f^{-1}$.
\end{dimostrazione}

Per le funzioni continue su un intervallo, dunque, «invertibile» e «strettamente
monotòna» sono la stessa cosa. In pratica: se la curva di taratura di uno strumento, come
quella dell'anemometro vista sopra, è continua e non assume mai due volte lo stesso
valore, è automaticamente monotòna, e la sua inversa è continua. Gli esempi che seguono
mostrano che l'ipotesi che il dominio sia un intervallo non si può omettere.

\begin{figure}[htb]
\centering
\begin{tikzpicture}
\begin{axis}[width=0.49\textwidth, height=4.6cm, xmin=-0.2, xmax=3.5, ymin=-0.2,
  ymax=3.4, xtick={1,2,3}, ytick={1,2,3}, xlabel={$x$}, ylabel={$y$},
  title={\footnotesize $y=f(x)$}]
  \addplot[curva, domain=0:1] {x};
  \addplot[curva, domain=2:3] {x-1};
  \addplot[tratteggio] coordinates {(0,1) (2,1)};
  \draw[blu, thick, fill=white] (axis cs:0,0) circle (1.8pt);
  \fill[blu] (axis cs:1,1) circle (1.8pt);
  \draw[blu, thick, fill=white] (axis cs:2,1) circle (1.8pt);
  \fill[blu] (axis cs:3,2) circle (1.8pt);
\end{axis}
\end{tikzpicture}\hfill
\begin{tikzpicture}
\begin{axis}[width=0.49\textwidth, height=4.6cm, xmin=-0.2, xmax=3.5, ymin=-0.2,
  ymax=3.4, xtick={1,2,3}, ytick={1,2,3}, xlabel={$y$}, ylabel={$x$},
  title={\footnotesize $x=f^{-1}(y)$}]
  \addplot[curva2, domain=0:1] {x};
  \addplot[curva2, domain=1:2] {x+1};
  \addplot[tratteggio] coordinates {(1,0) (1,2)};
  \draw[rame, thick, fill=white] (axis cs:0,0) circle (1.8pt);
  \fill[rame] (axis cs:1,1) circle (1.8pt);
  \draw[rame, thick, fill=white] (axis cs:1,2) circle (1.8pt);
  \fill[rame] (axis cs:2,3) circle (1.8pt);
\end{axis}
\end{tikzpicture}
\caption{La funzione dell'\cref{es:inversa-discontinua}(ii) e la sua inversa. La $f$ è
continua e strettamente crescente sull'insieme $(0,1]\cup(2,3]$, che non è un
intervallo; la sua inversa, definita sull'intervallo $(0,2]$, ha un salto nel punto
$1$.}
\label{fig:inversa-discontinua}
\end{figure}

\begin{esempio}[Se il dominio non è un intervallo]\label{es:inversa-discontinua}
\begin{enumerate}[label=\textup{(\roman*)}]
  \item La funzione $\frac1x$ è continua e iniettiva in $\R\setminus\{0\}$, ma non è
  monotòna: decresce su ciascuna delle due semirette, eppure $f(-1)=-1<1=f(1)$. La
  \cref{prop:continua-iniettiva-monotona} non si applica, perché il dominio non è un
  intervallo.
  \item Siano $A=(0,1]\cup(2,3]$ e
  \[
    f(x)=
    \begin{cases}
      x & \text{se } 0<x\le1,\\
      x-1 & \text{se } 2<x\le3 .
    \end{cases}
  \]
  La funzione è continua in $A$, perché in un intorno di ogni punto di $A$ coincide con
  un polinomio (\cref{lem:localita}), ed è strettamente crescente, perché i valori del
  primo tratto stanno in $(0,1]$ e quelli del secondo in $(1,2]$. La sua immagine è
  $(0,2]$, e la funzione inversa
  \[
    f^{-1}(y)=
    \begin{cases}
      y & \text{se } 0<y\le1,\\
      y+1 & \text{se } 1<y\le2
    \end{cases}
  \]
  ha nel punto $y=1$ un salto: $f^{-1}(1)=1$, mentre $f^{-1}(y)\to2$ per $y\to1^{+}$
  (\cref{fig:inversa-discontinua}). Il «buco» $(1,2]$ nel dominio di $f$ diventa un
  salto dell'inversa. Il dominio $A$ non è un intervallo, e questo impedisce la prima
  dimostrazione del \cref{teo:continuita-inversa}; non è chiuso, e questo impedisce la
  seconda: la successione $y_{n}=1+\frac1n$ tende a $1$, ma le controimmagini
  $x_{n}=2+\frac1n$ tendono a $2$, che non sta in $A$.
\end{enumerate}
\end{esempio}


%% ----------------------------------------------------------------------------
%%  §4.11  Esercizi di riepilogo
%% ----------------------------------------------------------------------------
\section{Esercizi di riepilogo}
\label{sec:esercizi-cap4}

Come nei capitoli precedenti, ogni soluzione è accompagnata da una nota sulla
\emph{scelta logica} che la guida: quale definizione o quale teorema conviene usare, e
perché.

\needspace{8\baselineskip}
\subsection*{La definizione di limite}

\begin{esercizio}\label{es:eps-delta-radice}
Verificare con la definizione \eqref{eq:eps-delta} che $\lim_{x\to2}(3x-1)=5$ e che
$\lim_{x\to4}\sqrt x=2$.
\begin{soluzione}
Nel primo caso $\abs{(3x-1)-5}=3\abs{x-2}$, che è minore di $\epsilon$ se
$\abs{x-2}<\delta=\frac\epsilon3$. Nel secondo caso razionalizziamo: per $x\ge0$,
\[
  \abs{\sqrt x-2}=\frac{\abs{x-4}}{\sqrt x+2}\le\frac{\abs{x-4}}2,
\]
che è minore di $\epsilon$ se $\abs{x-4}<\delta=2\epsilon$ (e $x\ge0$).
\emph{Scelta logica:} si cerca di maggiorare $\abs{f(x)-\ell}$ con una costante per $\abs{x-x_{0}}$;
la razionalizzazione trasforma una differenza di radici, difficile da stimare, in un
quoziente il cui denominatore è facile da minorare.
\end{soluzione}
\end{esercizio}

\begin{esercizio}\label{es:ponte-infinito}
Scrivere per esteso la dimostrazione dell'equivalenza tra le due formulazioni di
$\lim_{x\to+\infty}f(x)=\ell$, con $\ell\in\R$ e $A$ non limitato superiormente
(\cref{oss:ponte-infinito}).
\begin{soluzione}
Le due formulazioni sono: (S) per ogni $x_{n}\in A$ con $x_{n}\to+\infty$ si ha
$f(x_{n})\to\ell$; (D) per ogni $\epsilon>0$ esiste $k$ tale che $\abs{f(x)-\ell}<\epsilon$
per ogni $x\in A$ con $x>k$. Se vale (D) e $x_{n}\to+\infty$, fissato $\epsilon$ si
prende $k$ come in (D); esiste $\nu$ tale che $x_{n}>k$ per $n>\nu$, e quindi
$\abs{f(x_{n})-\ell}<\epsilon$ per $n>\nu$: vale (S). Se (D) fosse falsa, esisterebbe
$\epsilon_{0}>0$ tale che per ogni $k$ c'è un $x\in A$ con $x>k$ e
$\abs{f(x)-\ell}\ge\epsilon_{0}$; scegliendo $k=n$ si ottiene $x_{n}\in A$ con $x_{n}>n$,
quindi $x_{n}\to+\infty$, e $\abs{f(x_{n})-\ell}\ge\epsilon_{0}$ per ogni $n$: allora
$f(x_{n})\not\to\ell$ e (S) è falsa. \emph{Scelta logica:} è la dimostrazione del
\cref{teo:eps-delta}, in cui gli intorni di raggio $\frac1n$ di $x_{0}$ sono sostituiti
dalle semirette $(n,+\infty)$, gli intorni «di raggio $\frac1n$» di $+\infty$.
\end{soluzione}
\end{esercizio}

\needspace{8\baselineskip}
\subsection*{Calcolo di limiti}

\begin{esercizio}\label{es:calcolo-limiti-4}
Calcolare i limiti
\begin{gather*}
  \text{(a)}\ \lim_{x\to0}\frac{\sen5x}{\sen3x};\qquad
  \text{(b)}\ \lim_{x\to0}\frac{\tg x-\sen x}{x^{3}};\\
  \text{(c)}\ \lim_{x\to0}\frac{e^{2x}-1}{\sen x};\qquad
  \text{(d)}\ \lim_{x\to+\infty}\Bigl(1+\frac3x\Bigr)^{x}.
\end{gather*}
\begin{soluzione}
(a) $\frac{\sen5x}{\sen3x}=\frac{\sen5x}{x}\cdot\frac{x}{\sen3x}\to5\cdot\frac13=\frac53$,
per l'\cref{es:composte}(i). (b) Per $0<\abs x<\frac\pi2$,
$\tg x-\sen x=\sen x\Bigl(\frac1{\cos x}-1\Bigr)=\frac{\sen x\,(1-\cos x)}{\cos x}$, e quindi
\[
  \frac{\tg x-\sen x}{x^{3}}=\frac{\sen x}{x}\cdot\frac{1-\cos x}{x^{2}}\cdot\frac1{\cos x}
  \to1\cdot\frac12\cdot1=\frac12,
\]
per le \eqref{eq:lim-senx-x}, \eqref{eq:lim-1-cos} e la continuità del coseno.
(c) $\frac{e^{2x}-1}{\sen x}=2\cdot\frac{e^{2x}-1}{2x}\cdot\frac{x}{\sen x}\to2\cdot1\cdot1=2$,
per la \eqref{eq:ln-exp} con il cambio di variabile $y=2x$. (d) Con $y=\frac1x$, che
tende a $0^{+}$ ed è diverso da zero, $\bigl(1+\frac3x\bigr)^{x}=(1+3y)^{1/y}\to e^{3}$ per
la \eqref{eq:ln-e-b}. \emph{Scelta logica:} in tutti i casi la forma indeterminata si
scioglie riconducendo l'espressione, con semplici manipolazioni algebriche, a un prodotto
di limiti notevoli; il cambio di variabile è giustificato dal \cref{teo:limite-composta}.
\end{soluzione}
\end{esercizio}

\begin{esercizio}
Calcolare $\displaystyle\lim_{x\to+\infty}\bigl(\sqrt{x^{2}+x}-x\bigr)$.
\begin{soluzione}
È una forma $\infty-\infty$. Razionalizzando, per $x>0$,
\begin{align*}
  \sqrt{x^{2}+x}-x&=\frac{x^{2}+x-x^{2}}{\sqrt{x^{2}+x}+x}=\frac{x}{x\Bigl(\sqrt{1+\frac1x}+1\Bigr)}\\
  &=\frac1{\sqrt{1+\frac1x}+1}\ \to\ \frac1{1+1}=\frac12,
\end{align*}
perché $1+\frac1x\to1$ e la radice quadrata è continua nel punto $1$.
\emph{Scelta logica:} è lo stesso calcolo svolto per la successione $\sqrt{n^{2}+n}-n$ nel
\cref{cap:successioni}; la definizione «per successioni» garantisce che il risultato vale
lungo ogni $x_{n}\to+\infty$, non solo lungo i naturali.
\end{soluzione}
\end{esercizio}

\needspace{8\baselineskip}
\subsection*{Continuità e discontinuità}

\begin{esercizio}
Determinare il numero $k$ per cui la funzione $f(x)=\frac{\sen kx}{x}$ per $x\neq0$,
$f(0)=3$, è continua in $0$; determinare il numero $a$ per cui la funzione
$g(x)=ax+1$ per $x\le1$, $g(x)=x^{2}+2$ per $x>1$, è continua su $\R$.
\begin{soluzione}
Per l'\cref{es:composte}(i), $f(x)\to k$ per $x\to0$ (anche per $k=0$, quando
$f(x)=0$); $f$ è continua in $0$ se e solo se $k=3$. Per $g$, ciascun tratto è un
polinomio, e la sola continuità da verificare è in $x_{0}=1$: il limite sinistro è
$g(1)=a+1$, il limite destro è $1+2=3$; per la \cref{prop:dx-sx} serve $a+1=3$, cioè
$a=2$. \emph{Scelta logica:} la continuità in un punto di raccordo si verifica
confrontando i due limiti unilaterali con il valore.
\end{soluzione}
\end{esercizio}

\begin{esercizio}
Classificare le discontinuità delle funzioni $f(x)=\frac{x^{2}-1}{\abs{x-1}}$ e
$g(x)=\lfloor x\rfloor+\lfloor-x\rfloor$.
\begin{soluzione}
Per $x\neq1$ si ha $f(x)=\frac{(x-1)(x+1)}{\abs{x-1}}$, che vale $x+1$ per $x>1$ e
$-(x+1)$ per $x<1$: i limiti unilaterali in $1$ sono $2$ e $-2$, e $f$ ha in $1$ una
singolarità di prima specie (salto $4$). Per $g$: se $x\in\Z$, $g(x)=x-x=0$; se
$x\notin\Z$, posto $k=\lfloor x\rfloor$ si ha $k<x<k+1$, quindi $-k-1<-x<-k$ e
$\lfloor-x\rfloor=-k-1$, cioè $g(x)=-1$. Dunque $g$ vale $-1$ fuori dagli interi e $0$
sugli interi: in ogni intero il limite esiste ed è $-1\neq g(k)$, e la discontinuità è
eliminabile. \emph{Scelta logica:} conviene sempre semplificare l'espressione nei vari
intervalli prima di calcolare i limiti. Si noti che $g$ è somma di due funzioni con
salti, ma i salti si compensano: la somma di funzioni discontinue può essere «quasi»
continua.
\end{soluzione}
\end{esercizio}

\begin{esercizio}
Dimostrare che la funzione $f(x)=x$ per $x\in\Q$, $f(x)=0$ per $x\notin\Q$ è continua in
$x_{0}=0$ e discontinua in ogni $x_{0}\neq0$.
\begin{soluzione}
Si ha $\abs{f(x)}\le\abs x$ per ogni $x$, quindi $f(x)\to0=f(0)$ per $x\to0$ (carabinieri).
Se $x_{0}\neq0$, per la densità (\cref{teo:densita-Q}, \cref{cor:densita-irrazionali})
esistono successioni di razionali $q_{n}\to x_{0}$ e di irrazionali $r_{n}\to x_{0}$, con
$q_{n},r_{n}\neq x_{0}$; allora $f(q_{n})=q_{n}\to x_{0}$ e $f(r_{n})=0\to0\neq x_{0}$, e il
limite non esiste (\cref{cor:non-esistenza}). \emph{Scelta logica:} la continuità in $0$
viene da una stima; la discontinuità altrove da due successioni con limiti diversi.
\end{soluzione}
\end{esercizio}

\needspace{8\baselineskip}
\subsection*{Teoremi di esistenza}

\begin{esercizio}\label{es:grado-dispari}
Dimostrare che ogni polinomio di grado dispari ha almeno uno zero reale.
\begin{soluzione}
Sia $P(x)=a_{m}x^{m}+\dots+a_{0}$ con $m$ dispari e, ad esempio, $a_{m}>0$ (altrimenti si
considera $-P$). Per $x\neq0$
\[
  P(x)=x^{m}\Bigl(a_{m}+\frac{a_{m-1}}x+\dots+\frac{a_{0}}{x^{m}}\Bigr),
\]
e la parentesi tende ad $a_{m}>0$ per $x\to\pm\infty$. Poiché $m$ è dispari,
$x^{m}\to+\infty$ per $x\to+\infty$ e $x^{m}\to-\infty$ per $x\to-\infty$; quindi
$P(x)\to+\infty$ per $x\to+\infty$ e $P(x)\to-\infty$ per $x\to-\infty$. Per la
permanenza del segno (\cref{teo:confronto-funzioni}(i)) esistono $a<b$ con $P(a)<0<P(b)$,
e il teorema degli zeri, applicato al polinomio continuo $P$ su $[a,b]$, conclude.
\emph{Scelta logica:} i limiti all'infinito servono soltanto a trovare due punti in cui
il polinomio ha segno opposto. Per i polinomi di grado pari la conclusione è falsa:
$x^{2}+1$ non ha zeri reali.
\end{soluzione}
\end{esercizio}

\begin{esercizio}[Punto fisso]\label{es:punto-fisso}
Sia $f:[0,1]\to[0,1]$ continua. Dimostrare che esiste $x\in[0,1]$ con $f(x)=x$.
Applicare il risultato all'equazione $\cos x=x$, e calcolarne la soluzione con il metodo
di bisezione.
\begin{soluzione}
Sia $g(x)=f(x)-x$, continua su $[0,1]$. Poiché $0\le f(x)\le1$, si ha $g(0)=f(0)\ge0$ e
$g(1)=f(1)-1\le0$. Se $g(0)=0$ o $g(1)=0$ abbiamo finito; altrimenti $g(0)>0>g(1)$ e il
teorema degli zeri fornisce $x\in(0,1)$ con $g(x)=0$. La funzione $\cos x$ è continua e,
essendo decrescente su $[0,1]\subseteq[0,\pi]$, manda $[0,1]$ in
$[\cos1,1]\subseteq[0,1]$; dunque esiste $x$ con $\cos x=x$. Tale soluzione è unica,
perché $\cos x-x$ è strettamente decrescente su $[0,1]$. \emph{Scelta logica:} un
problema di punto fisso diventa un problema di zeri considerando la differenza
$f(x)-x$.
\end{soluzione}
\end{esercizio}

\begin{python}[La soluzione di $\cos x=x$]
import math

g = lambda x: math.cos(x) - x      # g(0) = 1 > 0, g(1) < 0
a, b = 0.0, 1.0
n = 0
while (b - a) / 2 >= 1e-12:        # errore del punto medio < 1e-12
    c = (a + b) / 2
    if g(c) > 0:                   # g decrescente: zero a destra
        a = c
    else:
        b = c
    n += 1
x = (a + b) / 2
print(f"x = {x:.12f}   cos x = {math.cos(x):.12f}   ({n} passi)")
\end{python}
\begin{risultato}
x = 0.739085133216   cos x = 0.739085133215   (39 passi)
\end{risultato}

Il numero $0{,}739085\ldots$ si ottiene anche premendo ripetutamente il tasto «cos» di
una calcolatrice (in radianti), partendo da un numero qualsiasi: le iterate
$x_{n+1}=\cos x_{n}$ convergono al punto fisso. Perché ciò accada, e con quale velocità, lo
capiremo con il calcolo differenziale.

\begin{esercizio}[L'equazione di Keplero]\label{es:kepler}
In un'orbita ellittica di eccentricità $e\in[0,1)$, la posizione del satellite al tempo
$t$ si ricava dall'\emph{anomalia eccentrica} $E$, soluzione dell'equazione di Keplero
\[
  E-e\sen E=M,
\]
dove l'\emph{anomalia media} $M=\frac{2\pi}{T}(t-t_{0})$ è proporzionale al tempo trascorso
dal passaggio al perigeo. Dimostrare che per ogni $M\in[0,\pi]$ l'equazione ha una e una
sola soluzione $E\in[0,\pi]$, e calcolarla per $e=0{,}5$, $M=1$ con errore minore di
$10^{-10}$.
\begin{soluzione}
La funzione $f(E)=E-e\sen E-M$ è continua, con $f(0)=-M\le0$ e
$f(\pi)=\pi-M\ge0$: per il teorema degli zeri (o direttamente, se $M=0$ o $M=\pi$) esiste
una soluzione in $[0,\pi]$. Per l'unicità basta osservare che $f$ è strettamente
crescente: se $E_{1}<E_{2}$, allora
$f(E_{2})-f(E_{1})=(E_{2}-E_{1})-e(\sen E_{2}-\sen E_{1})\ge(E_{2}-E_{1})(1-e)>0$, perché
$\abs{\sen E_{2}-\sen E_{1}}\le\abs{E_{2}-E_{1}}$ (si veda l'\cref{es:seno-lipschitz}).
Il numero di dimezzamenti necessari si conosce in anticipo: si deve avere
$\frac{\pi}{2^{n+1}}<10^{-10}$, cioè $n\ge34$. \emph{Scelta logica:} l'equazione è
trascendente e non ha soluzioni in forma chiusa; il teorema degli zeri ne garantisce
l'esistenza, la monotonia l'unicità, la bisezione il calcolo con errore controllato.
\end{soluzione}
\end{esercizio}

\begin{python}[L'equazione di Keplero con la bisezione]
import math

def bisezione(f, a, b, tol):
    """Approssima uno zero di f in [a, b] con errore < tol."""
    assert f(a) <= 0 <= f(b)
    passi = 0
    while (b - a) / 2 >= tol:      # errore del punto medio
        c = (a + b) / 2
        if f(c) > 0:
            b = c
        else:
            a = c
        passi += 1
    return (a + b) / 2, passi

e, M = 0.5, 1.0                    # eccentricità, anomalia media
kepler = lambda E: E - e * math.sin(E) - M
E, passi = bisezione(kepler, 0, math.pi, 1e-10)
print(f"E = {E:.10f} rad  dopo {passi} dimezzamenti")
print(f"verifica: E - e sen E = {E - e * math.sin(E):.10f}")
\end{python}
\begin{risultato}
E = 1.4987011335 rad  dopo 34 dimezzamenti
verifica: E - e sen E = 1.0000000000
\end{risultato}

\begin{esercizio}\label{es:seno-lipschitz}
Dimostrare che $\abs{\sen x-\sen y}\le\abs{x-y}$ per ogni $x,y\in\R$.
\begin{soluzione}
Per le formule di prostaferesi, che si ricavano dalle formule di addizione
\eqref{eq:addizione-sen} scrivendo $x=\frac{x+y}2+\frac{x-y}2$ e
$y=\frac{x+y}2-\frac{x-y}2$, si ha
$\sen x-\sen y=2\cos\frac{x+y}2\,\sen\frac{x-y}2$. Poiché $\abs{\cos t}\le1$ e
$\abs{\sen t}\le\abs t$ (\cref{lem:sen-minore-x}), risulta
$\abs{\sen x-\sen y}\le2\cdot1\cdot\frac{\abs{x-y}}2=\abs{x-y}$. \emph{Scelta logica:}
la disuguaglianza dà un'altra dimostrazione, con $\delta=\epsilon$, della continuità del
seno; funzioni con una stima di questo tipo si dicono \emph{lipschitziane}.
\end{soluzione}
\end{esercizio}

\begin{esercizio}
Quanti passi del metodo di bisezione servono per approssimare una radice in $[0,1]$
con errore minore di $10^{-6}$?
\begin{soluzione}
Per la \eqref{eq:bisezione-errore} basta $\frac1{2^{n+1}}<10^{-6}$, cioè
$2^{n+1}>10^{6}$. Poiché $2^{19}=524\,288<10^{6}<2^{20}=1\,048\,576$, serve $n+1\ge20$:
il punto medio $c_{19}$ ha la precisione richiesta; per individuarlo occorre conoscere
il segno di $f$ nei diciannove punti $c_{0},\dots,c_{18}$. \emph{Scelta logica:} la
stima dell'errore della bisezione non dipende dalla funzione, ma solo dalla lunghezza
dell'intervallo iniziale.
\end{soluzione}
\end{esercizio}

\begin{esercizio}
Dimostrare che una funzione continua su $[a,b]$ che assume soltanto valori razionali è
costante.
\begin{soluzione}
Se esistessero $x_{1},x_{2}$ con $f(x_{1})<f(x_{2})$, tra questi due valori ci sarebbe un
irrazionale $y$ (\cref{cor:densita-irrazionali}); per il teorema dei valori intermedi,
applicato su $[x_{1},x_{2}]$ o $[x_{2},x_{1}]$, $f$ assumerebbe il valore $y$, contro
l'ipotesi. \emph{Scelta logica:} l'immagine di un intervallo mediante una funzione
continua è un intervallo, e un intervallo con più di un punto contiene irrazionali.
\end{soluzione}
\end{esercizio}

\begin{esercizio}[Il monaco e la montagna]
Un monaco parte alle $6$ del mattino dal monastero a valle e sale lungo un sentiero fino
all'eremo in vetta, dove arriva alle $18$; il giorno dopo parte alle $6$ dall'eremo e
ridiscende lungo lo stesso sentiero, arrivando al monastero alle $18$. Dimostrare che
esiste un punto del sentiero in cui il monaco si trova alla stessa ora in entrambi i
giorni.
\begin{soluzione}
Sia $L$ la lunghezza del sentiero e siano $s_{1}(t)$, $s_{2}(t)$ le distanze dal
monastero, misurate lungo il sentiero, all'ora $t\in[6,18]$ del primo e del secondo
giorno; sono funzioni continue (il monaco non si teletrasporta). La funzione
$h(t)=s_{1}(t)-s_{2}(t)$ è continua, con $h(6)=0-L<0$ e $h(18)=L-0>0$: per il teorema degli
zeri esiste $\bar t$ con $s_{1}(\bar t)=s_{2}(\bar t)$. Si noti che non serve alcuna
ipotesi sulla velocità, che può anche annullarsi o cambiare segno. \emph{Scelta logica:}
si immagina che due monaci percorrano il sentiero lo stesso giorno, uno in salita e uno
in discesa; è evidente che si incontrano, e il teorema degli zeri rende rigoroso questo
«evidente».
\end{soluzione}
\end{esercizio}

\begin{esercizio}
Dimostrare che la funzione $f(x)=\frac{x}{1+x^{2}}$ ha massimo e minimo sull'intervallo
illimitato $[0,+\infty)$.
\begin{soluzione}
Per $x\ge0$ si ha $f(x)\ge0=f(0)$, quindi $0$ è il minimo. Inoltre $1+x^{2}\ge2x$, perché
equivale a $(x-1)^{2}\ge0$, e quindi $f(x)\le\frac12=f(1)$: $\frac12$ è il massimo.
\emph{Scelta logica:} il teorema di Weierstrass non si applica, perché l'intervallo non è
limitato; ma le sue ipotesi sono sufficienti, non necessarie, e qui massimo e minimo si
trovano con una disuguaglianza elementare.
\end{soluzione}
\end{esercizio}

\begin{esercizio}
Dimostrare che la funzione $f(x)=x^{3}+x$ è una biiezione continua e strettamente
crescente di $\R$ su $\R$, che la sua inversa è continua, e calcolare $f^{-1}(1)$.
\begin{soluzione}
$f$ è continua (polinomio) e strettamente crescente (somma di due funzioni strettamente
crescenti). Inoltre $f(x)\to\pm\infty$ per $x\to\pm\infty$, quindi per ogni $y\in\R$
esistono $a<b$ con $f(a)<y<f(b)$, e per il teorema dei valori intermedi $y$ è assunto: $f$
è suriettiva. Per ogni $[a,b]$ la restrizione di $f$ ha inversa continua su $[f(a),f(b)]$
(\cref{teo:continuita-inversa}), e ogni $y_{0}\in\R$ è interno a un tale intervallo:
per la località del limite $f^{-1}$ è continua su $\R$. Infine $f^{-1}(1)$ è la soluzione
di $x^{3}+x=1$, cioè della \eqref{eq:cubica}: $f^{-1}(1)=0{,}682327\ldots$, calcolata nel
\cref{sec:bisezione}. \emph{Scelta logica:} calcolare un valore della funzione inversa
equivale a risolvere un'equazione; la bisezione è dunque anche un algoritmo per
calcolare funzioni inverse.
\end{soluzione}
\end{esercizio}

%% ----------------------------------------------------------------------------
%%  §4.12  Punti fondamentali del capitolo
%% ----------------------------------------------------------------------------
\section{Punti fondamentali del capitolo}
\label{sec:sintesi-cap04}

L'idea centrale del capitolo è che il limite di una funzione si riduce ai limiti delle
successioni: $f(x)\to\ell$ se $f(x_{n})\to\ell$ lungo \emph{ogni} successione che si
avvicina a $x_{0}$. Così tutto il \cref{cap:successioni} si trasferisce alle funzioni.
La continuità chiede che il limite coincida con il valore; su un intervallo chiuso e
limitato, insieme alla completezza di $\R$, essa produce i grandi teoremi di esistenza.

\begin{puntichiave}
\item \textbf{Il limite per successioni.} Se $x_{0}$ è di accumulazione per $A$ (cioè
esiste una successione di punti di $A\setminus\{x_{0}\}$ che tende a $x_{0}$,
\cref{prop:caratt-accumulazione}), si ha $\lim_{x\to x_{0}}f(x)=\ell$ quando
$f(x_{n})\to\ell$ per ogni $x_{n}\in A\setminus\{x_{0}\}$ con $x_{n}\to x_{0}$
(\cref{def:limite-funzione}). Il valore $f(x_{0})$ non conta, e il limite, se esiste, è
unico (\cref{prop:unicita-limite-funzioni}).
\item \textbf{La formulazione con $\epsilon$ e $\delta$.} La definizione equivale a:
per ogni $\epsilon>0$ esiste $\delta>0$ tale che $\abs{f(x)-\ell}<\epsilon$ per ogni
$x\in A$ con $0<\abs{x-x_{0}}<\delta$ (\cref{teo:eps-delta}); $\delta$ dipende da
$\epsilon$. Per i punti e i limiti infiniti si usano gli intorni della
\cref{tab:intorni}.
\item \textbf{Località e limiti unilaterali.} Il limite dipende solo dai valori di $f$
vicino a $x_{0}$ (\cref{lem:localita}); se $x_{0}$ è di accumulazione da entrambi i
lati, il limite esiste se e solo se esistono e coincidono il limite destro e il limite
sinistro (\cref{prop:dx-sx}).
\item \textbf{Come si prova che un limite non esiste.} Bastano due successioni
ammissibili lungo le quali $f$ ha limiti diversi (\cref{cor:non-esistenza}): così
$\sen\frac1x$ non ha limite per $x\to0$ (\cref{es:sen-1-x}) e $\sen x$ non ha limite
per $x\to+\infty$ (\cref{es:sen-cos-infinito}).
\item \textbf{Regole di calcolo.} Somme, prodotti e quozienti (con limite del
denominatore non nullo), permanenza del segno, confronto, carabinieri e «limitata per
infinitesima» valgono come per le successioni, esclusi i casi di forma indeterminata
(\cref{teo:operazioni-funzioni,teo:confronto-funzioni}).
\item \textbf{Limite di una funzione composta.} Se $g(x)\to y_{0}$ e $f(y)\to\ell$,
allora $f(g(x))\to\ell$ purché $g(x)\neq y_{0}$ per $x\neq x_{0}$ vicino a $x_{0}$
(\cref{teo:limite-composta}); senza questa ipotesi la tesi può essere falsa
(\cref{oss:composta-controesempio}), ma è vera se $f$ è continua in $y_{0}$
(\cref{cor:limite-composta-continua}).
\item \textbf{Limiti notevoli.} Per $x\to0$: $\frac{\sen x}{x}\to1$,
$\frac{1-\cos x}{x^{2}}\to\frac12$, $(1+x)^{1/x}\to e$ (\cref{es:composte}),
$\frac{\log(1+x)}{x}\to1$, $\frac{e^{x}-1}{x}\to1$, $\frac{a^{x}-1}{x}\to\log a$,
$\frac{(1+x)^{b}-1}{x}\to b$ (\cref{prop:limiti-notevoli-exp}).
\item \textbf{Funzioni continue.} $f$ è continua in $x_{0}\in A$ se
$\lim_{x\to x_{0}}f(x)=f(x_{0})$, cioè se $f(x_{n})\to f(x_{0})$ per ogni successione di
punti di $A$ convergente a $x_{0}$, anche con $x_{n}=x_{0}$
(\cref{prop:caratt-continuita}). Somme, prodotti, quozienti e composte di funzioni
continue sono continue, e le funzioni elementari sono continue nel loro dominio
(\cref{tab:elementari-continue}).
\item \textbf{Discontinuità e prolungamento.} Le discontinuità sono eliminabili, di
prima specie (salto) o di seconda specie (\cref{def:discontinuita}); una funzione si
prolunga per continuità in $x_{0}$ se e solo se esiste finito il limite in $x_{0}$
(\cref{def:prolungamento}), come $\frac{\sen x}{x}$ in $0$.
\item \textbf{Teorema degli zeri e bisezione.} Se $f$ è continua in $[a,b]$ e $f(a)$,
$f(b)$ hanno segno discorde, esiste $x_{0}\in(a,b)$ con $f(x_{0})=0$ (\cref{teo:zeri}).
La dimostrazione per bisezione è un algoritmo, con errore
$\abs{c_{n}-x_{0}}\le\frac{b-a}{2^{n+1}}$, e poggia sulla completezza: in $\Q$ il
teorema è falso.
\item \textbf{Valori intermedi e Weierstrass.} Una funzione continua nell'intervallo
\emph{chiuso e limitato} $[a,b]$ ha minimo $m$ e massimo $M$ (\cref{teo:weierstrass}) e
assume tutti i valori compresi tra essi: $f([a,b])=[m,M]$
(\cref{teo:valori-intermedi-2,cor:immagine-intervallo}). Le ipotesi sono sufficienti, non
necessarie, e nessuna può essere omessa (\cref{fig:controesempi-weierstrass}).
\item \textbf{Monotone e inverse.} Nei punti interni di $[a,b]$ una funzione monotòna
può avere solo discontinuità di salto (\cref{cor:monotone-salti}); su $[a,b]$ è continua se e solo se la sua immagine è
l'intervallo di estremi $f(a)$, $f(b)$ (\cref{teo:criterio-monotone}). L'inversa di una
funzione continua e strettamente monotòna su $[a,b]$ è continua (\cref{teo:continuita-inversa}):
da qui la continuità di $a^{x}$, $\log_{a}x$, $x^{b}$ (\cref{teo:continuita-esponenziale}).
\item \textbf{Errori tipici.} Confondere il limite con il valore; dedurre un limite da
una tabella o da un calcolo numerico, che può soffrire di cancellazione; far passare al
limite una disuguaglianza stretta ($x^{2}>0$, ma $x^{2}\to0$); applicare il teorema
degli zeri fuori da un intervallo ($\frac1x$ cambia segno su $\R\setminus\{0\}$ senza
annullarsi).
\item \textbf{Applicazioni aeronautiche.} I valori intermedi garantiscono un'incidenza
di equilibrio portanza--peso; Weierstrass garantisce l'efficienza massima della polare
parabolica, $E_{\max}=\frac1{2\sqrt{kC_{D0}}}$ per $C_{L}^{*}=\sqrt{C_{D0}/k}$
(\cref{fig:efficienza}); la bisezione risolve l'equazione di Keplero
(\cref{es:kepler}); l'onda d'urto normale è un esempio fisico di salto.
\end{puntichiave}

\paragraph{Il filo del discorso.} La continuità dice che a piccole variazioni della
variabile corrispondono piccole variazioni della funzione, ma non \emph{quanto}
piccole: non distingue il raccordo liscio dallo spigolo della tropopausa, né spiega la
diversa precisione dell'anemometro alle varie velocità. I teoremi di esistenza, poi,
non dicono dove si trovi un massimo. Nel \cref{cap:derivate} il limite del rapporto
incrementale, già intravisto con la velocità istantanea, misurerà la rapidità di
variazione e darà gli strumenti per localizzare massimi e minimi.

